% Lesson 11 — Vocabulary: Words and expressions related to taking part in recreational activities — Anglais 1ère A — Cours signé OnBuch+
% Programme officiel MINESEC. Compiler avec : tectonic ENA11-vocabulary-words-and-expressions-related-to-taking-part.tex
\newcommand{\DOCMATIERE}{Anglais}
\newcommand{\DOCNIVEAU}{1ère A}
\newcommand{\DOCMODULE}{Chapter 1 : Family and Social Life}
\newcommand{\DOCLECON}{ENA11}
\newcommand{\DOCTITRE}{Lesson 11 — Vocabulary: Words and expressions related to taking part in recreational activities}
% OnBuch+ — préambule « cours signés » ENRICHI (compilé avec tectonic / XeLaTeX)
% Système visuel « Pop » d'OnBuch+ : fond crème (jamais blanc), bordures encre
% épaisses, ombres « dures » décalées, titres Archivo Black en capitales.
% Version MATHÉMATIQUES : graphes (pgfplots/tikz), boîtes pédagogiques
% (définition, propriété, méthode, attention, le savais-tu, exercice résolu,
% exercice, TP, bilan), page de garde et signature OnBuch+ ancrée en bas.
%
% Macros attendues dans main.tex AVANT \input{preamble} :
%   \DOCMATIERE  \DOCNIVEAU  \DOCMODULE  \DOCLECON (numéro)  \DOCTITRE
\documentclass[11pt]{article}

\usepackage{fontspec}
\usepackage[english,french]{babel} % français = langue principale ; l'anglais via \eng{..} / textebox
\usepackage[a4paper,margin=2cm,top=3.4cm,bottom=2.8cm,headheight=1.4cm,footskip=1.3cm]{geometry}
\usepackage{amsmath,amssymb,amsfonts}
\usepackage{mathtools} % \xmapsto, \coloneqq... souvent utilisés par les modèles
\usepackage{cancel} % simplifications barrées (bilans de forces)
% \abs{z} (module, valeur absolue) et \norm{u} (norme), utilisés par les modèles
\providecommand{\abs}{}\let\abs\relax\DeclarePairedDelimiter{\abs}{\lvert}{\rvert}
\providecommand{\norm}{}\let\norm\relax\DeclarePairedDelimiter{\norm}{\lVert}{\rVert}
\usepackage[table]{xcolor} % [table] : fournit aussi \rowcolors (alternance de lignes)
\usepackage{pagecolor}
\usepackage{tikz}
\usetikzlibrary{arrows.meta,positioning,calc,shapes.geometric,shapes.misc,shapes.arrows,decorations.pathmorphing,decorations.pathreplacing,decorations.markings,patterns,shadows,mindmap,trees,fit,backgrounds,matrix,intersections,angles,quotes}
\usepackage{pgfplots}
\usepackage[version=4]{mhchem} % équations nucléaires \ce{^{235}_{92}U}
\usepackage[european,siunitx]{circuitikz} % schémas électriques
\pgfplotsset{compat=1.18}
\usepgfplotslibrary{fillbetween,groupplots} % aires entre courbes (calcul intégral)
\usepackage{enumitem}
\usepackage{fancyhdr}
% [most] charge toutes les bibliothèques tcolorbox (listings, minted, raster,
% documentation...) et gonfle inutilement la mémoire TeX sur un document long
% (~300 boîtes) : on ne charge que ce qui est réellement utilisé.
\usepackage[skins,breakable]{tcolorbox}
\usepackage{graphicx}
\usepackage{colortbl}
\usepackage{booktabs}
\usepackage{tabularx}
\usepackage{diagbox} % case d'en-tête diagonale des tableaux à double entrée
% Colonnes extensibles centrée / gauche / droite (C, L, R), souvent utilisées par les modèles
\newcolumntype{C}{>{\centering\arraybackslash}X}
\newcolumntype{L}{>{\raggedright\arraybackslash}X}
\newcolumntype{R}{>{\raggedleft\arraybackslash}X}
\usepackage{array}
\usepackage{multicol}
\usepackage{multirow}
\usepackage{siunitx}
% parse-numbers=false : accepte aussi \SI{2,5 \times 10^{-3}}{...}
\sisetup{locale=FR,output-decimal-marker={,},group-minimum-digits=4,parse-numbers=false}
\DeclareSIUnit\ohm{\ensuremath{\symup{\Omega}}} % U+2126 absent de la police
\DeclareSIUnit\division{div} % réglages d'oscilloscope (ms/div, V/div)
\DeclareSIUnit\tr{tr} % tours (vitesse de rotation, ex. \SI{1500}{\tr\per\minute})
\DeclareSIUnit\molar{M} % concentration molaire (ex. \SI{3}{\molar}, \SI{1}{\micro\molar})
\DeclareSIUnit\an{an} % année (le modèle confond parfois avec \year, un compteur TeX) — ex. \SI{5}{\kilo\meter\per\an}
\DeclareSIUnit\FCFA{FCFA} % franc CFA (ex. \SI{500}{\FCFA\per\kilo\gram})
\DeclareSIUnit\degreeBrix{\textdegree Brix} % teneur en sucre d'un sirop/jus (réfractométrie)
\DeclareSIUnit\month{mois} % le modèle utilise parfois \month, un compteur TeX, comme unité de durée
\usepackage{qrcode}
% \degree : commande fréquemment utilisée par les modèles pour un angle ou une
% température en dehors de \SI{}{} (non fournie par les paquets chargés).
\providecommand{\degree}{^{\circ}}
% \qed (fin de démonstration), utilisé par les modèles sans amsthm
\providecommand{\qed}{\hfill$\square$}
% \barre : double barre des valeurs interdites dans un tableau de variations (tkz-tab)
\providecommand{\barre}{\ensuremath{\|}}
% environnement proof (amsthm non chargé) utilisé par les modèles
\ifdefined\proof\else\newenvironment{proof}[1][Démonstration]{\par\noindent\textit{#1.}\ }{\qed\par}\fi
\usepackage{lastpage}
\usepackage{hyperref}
\hypersetup{hidelinks}

% ============================================================
% Palette « Pop » OnBuch+ — mêmes hex que la classe P (plus_ui.dart)
% ============================================================
\definecolor{poporange}{HTML}{F59321}
\definecolor{popdark}{HTML}{E07A0C}
\definecolor{popcream}{HTML}{FFF6E8}
\definecolor{popink}{HTML}{1C1714}
\definecolor{poppurple}{HTML}{7A5AE0}
\definecolor{popgreen}{HTML}{2E9E6A}
\definecolor{poppink}{HTML}{E0457B}
\definecolor{popblue}{HTML}{2F7DE1}
\definecolor{popgold}{HTML}{C98A12}
\definecolor{popmuted}{HTML}{8A7F76}
\definecolor{popline}{HTML}{EADFD2}
\definecolor{popgray}{HTML}{8A7F76}
\colorlet{popgrey}{popgray}
% Alias tolérants (le modèle nomme parfois les couleurs différemment)
\colorlet{popwhite}{white}
\colorlet{popblack}{popink}
\colorlet{popred}{poppink}
\colorlet{popyellow}{popgold}
% Teintes claires (fonds de tableaux/encadrés) — le modèle nomme parfois
% les couleurs avec un suffixe L (ex. popblueL) sans les avoir définies.
\colorlet{poporangeL}{poporange!20}
\colorlet{popdarkL}{popdark!20}
\colorlet{poppurpleL}{poppurple!20}
\colorlet{popgreenL}{popgreen!20}
\colorlet{poppinkL}{poppink!20}
\colorlet{popblueL}{popblue!20}
\colorlet{popgoldL}{popgold!20}
\colorlet{popredL}{popred!20}
\colorlet{popgrayL}{popgray!20}
% Teintes claires pour fonds de tableaux / zones
\colorlet{popblueL}{popblue!10!white}
\colorlet{poporangeL}{poporange!14!white}
\colorlet{popgreenL}{popgreen!12!white}
\colorlet{poppurpleL}{poppurple!12!white}
\colorlet{poppinkL}{poppink!12!white}
\colorlet{popgoldL}{popgold!14!white}

\pagecolor{popcream}

% ============================================================
% Typographie
% ============================================================
\defaultfontfeatures{Ligatures=TeX}
\setmainfont{PlusJakartaSans-Medium.ttf}[Path=../../../fonts/,BoldFont=PlusJakartaSans-Bold.ttf]
% Maths en sans-serif (Fira Math) avec chiffres et lettres droites de Plus
% Jakarta, pour que les formules se fondent dans le texte.
\usepackage[math-style=ISO,bold-style=ISO]{unicode-math}
\setmathfont{FiraMath-Regular.otf}[Path=../../../fonts/]
\setmathfont{PlusJakartaSans-Medium.ttf}[Path=../../../fonts/,range={up/num,up/{Latin,latin},\mathup}]
\usepackage{icomma} % virgule décimale sans espace en mode math
% Caractères mathématiques Unicode absents des polices de texte (Plus Jakarta,
% Archivo Black) : rendus par leur équivalent mathématique, en texte comme en
% formule — sinon ils s'affichent en carré vide (ex. « Arithmétique dans ℤ »).
\usepackage{newunicodechar}
\newunicodechar{ℝ}{\ensuremath{\mathbb{R}}}
\newunicodechar{ℤ}{\ensuremath{\mathbb{Z}}}
\newunicodechar{ℂ}{\ensuremath{\mathbb{C}}}
\newunicodechar{ℕ}{\ensuremath{\mathbb{N}}}
\newunicodechar{ℚ}{\ensuremath{\mathbb{Q}}}
\newunicodechar{𝔻}{\ensuremath{\mathbb{D}}}
\newunicodechar{α}{\ensuremath{\alpha}}
\newunicodechar{β}{\ensuremath{\beta}}
\newunicodechar{γ}{\ensuremath{\gamma}}
\newunicodechar{δ}{\ensuremath{\delta}}
\newunicodechar{ε}{\ensuremath{\varepsilon}}
\newunicodechar{θ}{\ensuremath{\theta}}
\newunicodechar{λ}{\ensuremath{\lambda}}
\newunicodechar{μ}{\ensuremath{\mu}}
\newunicodechar{σ}{\ensuremath{\sigma}}
\newunicodechar{τ}{\ensuremath{\tau}}
\newunicodechar{φ}{\ensuremath{\varphi}}
\newunicodechar{ψ}{\ensuremath{\psi}}
\newunicodechar{ω}{\ensuremath{\omega}}
\newunicodechar{Σ}{\ensuremath{\Sigma}}
% majuscules produites par \MakeUppercase dans les titres (α-aminés -> Α)
\newunicodechar{Α}{\ensuremath{\alpha}}
\newunicodechar{Β}{\ensuremath{\beta}}
\newunicodechar{Γ}{\ensuremath{\Gamma}}
\newunicodechar{Δ}{\ensuremath{\Delta}}
\newunicodechar{Ε}{\ensuremath{\varepsilon}}
\newunicodechar{Θ}{\ensuremath{\Theta}}
\newunicodechar{Λ}{\ensuremath{\Lambda}}
\newunicodechar{Μ}{\ensuremath{\mu}}
\newunicodechar{Τ}{\ensuremath{\tau}}
\newunicodechar{Φ}{\ensuremath{\Phi}}
\newunicodechar{Ψ}{\ensuremath{\Psi}}
\newunicodechar{Ω}{\ensuremath{\Omega}}
\newunicodechar{↦}{\ensuremath{\mapsto}}
\newunicodechar{∈}{\ensuremath{\in}}
\newunicodechar{∉}{\ensuremath{\notin}}
\newunicodechar{∧}{\ensuremath{\wedge}}
\newunicodechar{⇔}{\ensuremath{\Leftrightarrow}}
\newunicodechar{⟺}{\ensuremath{\Longleftrightarrow}}
\newunicodechar{⇒}{\ensuremath{\Rightarrow}}
\newunicodechar{⟹}{\ensuremath{\Longrightarrow}}
\newunicodechar{∘}{\ensuremath{\circ}}
\newunicodechar{∪}{\ensuremath{\cup}}
\newunicodechar{∩}{\ensuremath{\cap}}
\newunicodechar{≡}{\ensuremath{\equiv}}
\newunicodechar{⊂}{\ensuremath{\subset}}
\newunicodechar{⊥}{\ensuremath{\perp}}
\newunicodechar{∃}{\ensuremath{\exists}}
\newunicodechar{∀}{\ensuremath{\forall}}
\newunicodechar{∛}{\ensuremath{\sqrt[3]{\,}}}
\newunicodechar{∜}{\ensuremath{\sqrt[4]{\,}}}
\newunicodechar{⃗}{\ensuremath{{}^{\to}}}
\newunicodechar{ⁿ}{\ifmmode^{n}\else\textsuperscript{\ensuremath{n}}\fi}
\newunicodechar{ˣ}{\ifmmode^{x}\else\textsuperscript{\ensuremath{x}}\fi}
\newunicodechar{ᵇ}{\ifmmode^{b}\else\textsuperscript{\ensuremath{b}}\fi}
\newunicodechar{ʳ}{\ifmmode^{r}\else\textsuperscript{\ensuremath{r}}\fi}
\newunicodechar{ᵘ}{\ifmmode^{u}\else\textsuperscript{\ensuremath{u}}\fi}
\newunicodechar{ʸ}{\ifmmode^{y}\else\textsuperscript{\ensuremath{y}}\fi}
\newunicodechar{ᵃ}{\ifmmode^{a}\else\textsuperscript{\ensuremath{a}}\fi}
\newunicodechar{ᵉ}{\ifmmode^{e}\else\textsuperscript{\ensuremath{e}}\fi}
\newunicodechar{ᵏ}{\ifmmode^{k}\else\textsuperscript{\ensuremath{k}}\fi}
\newunicodechar{ᵖ}{\ifmmode^{p}\else\textsuperscript{\ensuremath{p}}\fi}
\newunicodechar{ᴺ}{\ifmmode^{N}\else\textsuperscript{\ensuremath{N}}\fi}
\newunicodechar{ᶜ}{\ifmmode^{c}\else\textsuperscript{\ensuremath{c}}\fi}
\newunicodechar{ʰ}{\ifmmode^{h}\else\textsuperscript{\ensuremath{h}}\fi}
\newunicodechar{ᵠ}{\ifmmode^{q}\else\textsuperscript{\ensuremath{q}}\fi}
\newunicodechar{ₙ}{\ifmmode_{n}\else\textsubscript{\ensuremath{n}}\fi}
\newunicodechar{ₐ}{\ifmmode_{a}\else\textsubscript{\ensuremath{a}}\fi}
\newunicodechar{ᵢ}{\ifmmode_{i}\else\textsubscript{\ensuremath{i}}\fi}
\newunicodechar{ᵣ}{\ifmmode_{r}\else\textsubscript{\ensuremath{r}}\fi}
\newunicodechar{ₖ}{\ifmmode_{k}\else\textsubscript{\ensuremath{k}}\fi}
\newunicodechar{ᵦ}{\ifmmode_{\beta}\else\textsubscript{\ensuremath{\beta}}\fi}
\newunicodechar{ₓ}{\ifmmode_{x}\else\textsubscript{\ensuremath{x}}\fi}
\newfontfamily\popdisplay{ArchivoBlack-Regular.ttf}[Path=../../../fonts/]
\newfontfamily\poplogo{SpaceGrotesk-Bold.ttf}[Path=../../../fonts/]
\newfontfamily\popsemi{PlusJakartaSans-SemiBold.ttf}[Path=../../../fonts/]
\newfontfamily\popxbold{PlusJakartaSans-ExtraBold.ttf}[Path=../../../fonts/]
\newfontfamily\popxboldls{PlusJakartaSans-ExtraBold.ttf}[Path=../../../fonts/,LetterSpace=3.0]

\setlist[enumerate]{leftmargin=1.7em,itemsep=5pt,topsep=4pt}
\setlist[itemize]{leftmargin=1.7em,itemsep=4pt,topsep=4pt,label={\color{poporange}\textbullet}}
\setlength{\parindent}{0pt}
\setlength{\parskip}{5pt}
\renewcommand{\arraystretch}{1.25}

% pgfplots : style Pop par défaut
\pgfplotsset{
  every axis/.append style={axis line style={popink,line width=1pt},tick style={popink},
    grid style={popline,line width=0.5pt},label style={font=\small},tick label style={font=\footnotesize}},
  popcurve/.style={poporange,line width=1.8pt,no markers,smooth},
}
% Même style disponible dans un \draw TikZ simple (hors pgfplots)
\tikzset{popcurve/.style={poporange,line width=1.8pt}}
\tikzset{popgrid/.style={popline,very thin}}
\colorlet{popcurve}{poporange} % le nom du style est parfois utilisé comme couleur

% ============================================================
% Pill / squiggle
% ============================================================
\newcommand{\poppill}[3]{%
  \tikz[baseline=-0.55ex]\node[draw=popink,line width=1.2pt,rounded corners=2.6pt,%
    fill=#2,inner xsep=6.5pt,inner ysep=3pt]{%
    \popxboldls\fontsize{7}{7}\selectfont\color{#3}\MakeUppercase{#1}};%
}
\newcommand{\popsquiggle}[1]{%
  \tikz[baseline=-0.4ex]{
    \draw[#1,line width=2.6pt,line cap=round]
      plot [smooth] coordinates {(0,0.09) (0.34,0.01) (0.68,0.21) (1.02,0.01) (1.36,0.10)};
  }%
}
% Mot-clé surligné (marqueur orange sous le texte)
\newcommand{\cle}[1]{{\popxbold\color{popdark}#1}}

% ============================================================
% En-tête / pied
% ============================================================
\pagestyle{fancy}
\fancyhf{}
\renewcommand{\headrulewidth}{0pt}
\renewcommand{\footrulewidth}{0pt}
\lhead{\raisebox{-0.15ex}{\poplogo\fontsize{13}{13}\selectfont\color{popink}OnBuch\color{poporange}+}}
\rhead{\poppill{\DOCMATIERE\ \textbullet\ \DOCNIVEAU\ \textbullet\ Leçon \DOCLECON}{white}{popink}}
\lfoot{\popxbold\fontsize{7.6}{7.6}\selectfont\color{popmuted}\MakeUppercase{Cours signé OnBuch+}\ \textbullet\ {\popsemi\DOCTITRE}}
\rfoot{\tikz[baseline=-0.6ex]\node[rounded corners=8.5pt,fill=popink,minimum height=17pt,minimum width=17pt,inner xsep=5pt,inner sep=0]{\popxbold\fontsize{8}{8}\selectfont\color{popcream}\thepage\,/\,\pageref*{LastPage}};}

% ============================================================
% Titres
% ============================================================
\newcommand{\chaptitre}[2]{%
  \begin{tcolorbox}[enhanced,colback=poporange,colframe=popink,boxrule=2.6pt,arc=11pt,
      drop shadow=popink,left=15pt,right=15pt,top=12pt,bottom=11pt,before skip=3pt,after skip=18pt]
    \poppill{#2}{popink}{popcream}\par\vspace{9pt}
    {\raggedright\hyphenpenalty=10000\exhyphenpenalty=10000\popdisplay\fontsize{22}{24}\selectfont\color{popcream}\MakeUppercase{#1}\par}
  \end{tcolorbox}
}
% Section numérotée : pastille orange avec le numéro + titre Archivo
\newcounter{popsec}
\newcommand{\coursec}[1]{%
  \stepcounter{popsec}\setcounter{popsub}{0}%
  \par\vspace{16pt}\noindent
  \begin{minipage}{\linewidth}
  \tikz[baseline=-0.7ex]\node[draw=popink,line width=1.6pt,fill=poporange,rounded corners=4pt,minimum size=22pt,inner sep=0]{\popdisplay\fontsize{12}{12}\selectfont\color{popcream}\thepopsec};%
  \hspace{8pt}{\popdisplay\fontsize{15}{17}\selectfont\color{popink}\MakeUppercase{#1}}\par
  \vspace{4pt}\noindent{\color{poporange}\rule{36pt}{3.6pt}}
  \end{minipage}\par\nopagebreak\vspace{8pt}
}
\newcounter{popsub}
\newcommand{\courssub}[1]{%
  \stepcounter{popsub}%
  \par\vspace{9pt}\noindent{\popxbold\fontsize{11.5}{13}\selectfont\color{popink}{\color{poporange}\thepopsec.\thepopsub}\hspace{6pt}#1}\par\nopagebreak\vspace{4pt}
}

% ============================================================
% Boîtes pédagogiques — même recette : fond blanc, bordure épaisse colorée,
% ombre dure encre, badge plein. Titre optionnel [#1].
% ============================================================
\tcbset{popbase/.style={breakable,enhanced,colback=white,boxrule=2.3pt,arc=9pt,
  shadow={3pt}{-3pt}{0pt}{popink},left=14pt,right=14pt,top=11pt,bottom=11pt,before skip=11pt,after skip=15pt,
  fontupper=\popsemi\fontsize{10.6}{14.5}\selectfont\color{popink},
  fontlower=\popsemi\fontsize{10.6}{14.5}\selectfont\color{popink}}}
% Coche dessinée (le glyphe ✓ n'existe pas dans les polices du document)
\newcommand{\popcheck}[1]{\tikz[baseline=-0.5ex]\draw[#1,line width=1.6pt,line cap=round,line join=round](-0.13,0.0)--(-0.04,-0.09)--(0.13,0.1);}
\providecommand{\coche}{\popcheck{popgreen}}
\providecommand{\resultat}[1]{\par\smallskip\noindent\fcolorbox{popgreen}{popgreenL}{\parbox{\dimexpr\linewidth-2\fboxsep-2\fboxrule\relax}{\textbf{Résultat.} #1}}\par\smallskip}
\newcommand{\popboxhead}[3]{\poppill{#1}{#2}{white}\ifx\relax#3\relax\else\hspace{6pt}{\popxbold\fontsize{10.6}{12}\selectfont\color{popink}#3}\fi\par\vspace{7pt}}

\newtcolorbox{aretenir}[1][]{popbase,colframe=popgreen,before upper={\popboxhead{À retenir}{popgreen}{#1}}}
% Texte en anglais (document de lecture, dialogue, extrait) : cadre
% d'étude. Un titre optionnel (source, auteur) s'affiche dans l'en-tête.
\newtcolorbox{textebox}[1][]{popbase,colframe=popblue,before upper={\popboxhead{Text}{popblue}{#1}\begin{otherlanguage}{english}},after upper={\end{otherlanguage}}}
% Marques ✓ / ✗ (forme correcte / fautive) souvent utilisées par les modèles
\providecommand{\cross}{\textcolor{poppink}{\ensuremath{\times}}}
\providecommand{\xmark}{\cross}\providecommand{\crossmark}{\cross}\providecommand{\cmark}{\ensuremath{\checkmark}}
% Ligne de réponse / justification à compléter (inventée par les modèles)
\providecommand{\justification}[1]{\par\noindent\textit{Justification~:} \dotfill\par}
% \textipa (package tipa non chargé) : la police ne couvre pas l'API -> simple passage
\providecommand{\textipa}[1]{#1}
% Soulignement (ulem non chargé) : \uline, \ul
\providecommand{\uline}[1]{\underline{#1}}\providecommand{\ul}[1]{\underline{#1}}
% \glossaire{\item ...} inventée par les modèles : liste de glossaire
\providecommand{\glossaire}[1]{\begin{itemize}#1\end{itemize}}
% \alert (beamer) inventée par les modèles : texte mis en évidence
\providecommand{\alert}[1]{\textbf{\textcolor{poppink}{#1}}}
% variantes ASCII d'ordinaux écrites en macro
\providecommand{\iere}{\textsuperscript{re}}\providecommand{\ieme}{\textsuperscript{e}}\providecommand{\ere}{\textsuperscript{re}}\providecommand{\eme}{\textsuperscript{e}}\providecommand{\er}{\textsuperscript{er}}\providecommand{\ie}{\textsuperscript{e}}
% Ordinaux français écrits en macro par les modèles : 1\ière, 3\ième
\newcommand{\ière}{\textsuperscript{re}}\newcommand{\ième}{\textsuperscript{e}}
% \exemple (inventée par les modèles) : étiquette « Exemple : »
\providecommand{\exemple}{\textbf{Exemple~:}\ }
% \square utilisable aussi hors mode math (cases à cocher)
\AtBeginDocument{\let\squareMath\square\renewcommand{\square}{\ensuremath{\squareMath}}}
% Mot ou phrase en anglais à l'intérieur d'un paragraphe français
\newcommand{\eng}[1]{\ifmmode\text{\textit{#1}}\else\begingroup\selectlanguage{english}\itshape #1\endgroup\fi}
\let\esp\eng % alias toléré
% flèches d'intonation inventées par les modèles
\providecommand{\textsearrow}{}\renewcommand{\textsearrow}{\ensuremath{\searrow}}\providecommand{\textnearrow}{}\renewcommand{\textnearrow}{\ensuremath{\nearrow}}
% \fr{..} : mot français (inventé par les modèles)
\providecommand{\fr}[1]{\textit{#1}}
\newtcolorbox{exemplebox}[1][]{popbase,colframe=poporange,before upper={\popboxhead{Exemple}{poporange}{#1}}}
\newtcolorbox{definition}[1][]{popbase,colframe=popblue,before upper={\popboxhead{Définition}{popblue}{#1}}}
\newtcolorbox{propriete}[1][]{popbase,colframe=poppurple,before upper={\popboxhead{Propriété}{poppurple}{#1}}}
\newtcolorbox{methode}[1][]{popbase,colframe=popgold,before upper={\popboxhead{Méthode}{popgold}{#1}}}
\newtcolorbox{attention}[1][]{popbase,colframe=poppink,before upper={\popboxhead{Attention}{poppink}{#1}}}
\newtcolorbox{experience}[1][]{popbase,colframe=popblue,colback=popblueL,before upper={\popboxhead{Expérience}{popblue}{#1}}}
% « Le savais-tu ? » : carte violette pleine (contexte, culture, Cameroun)
\newtcolorbox{savaistu}[1][]{popbase,colback=poppurple,colframe=popink,
  fontupper=\popsemi\fontsize{10.4}{14.2}\selectfont\color{white},
  before upper={\poppill{Le savais-tu\,?}{white}{poppurple}\ifx\relax#1\relax\else\hspace{6pt}{\popxbold\color{white}#1}\fi\par\vspace{7pt}}}
% Exercice résolu : énoncé en haut, \tcblower puis la solution sur fond crème
\newcounter{popexo}\newcounter{popexr}
\newtcolorbox{exoresolu}[1][]{popbase,colframe=poporange,colbacklower=popcream,
  segmentation style={popink,line width=1.2pt,dashed},
  before upper={\stepcounter{popexr}\popboxhead{Exercice résolu \thepopexr}{poporange}{#1}},
  before lower={\poppill{Solution}{popink}{popcream}\par\vspace{6pt}}}
% Exercice (non résolu en place) — la correction est dans la partie Corrigés
\newtcolorbox{exercice}[1][]{popbase,colframe=popink,
  before upper={\stepcounter{popexo}\popboxhead{Exercice \thepopexo}{popink}{#1}}}
% Corrigé
\newtcolorbox{corrige}[1][]{popbase,colframe=popgreen,colback=popgreenL,
  before upper={\popboxhead{Corrigé}{popgreen}{#1}}}
% Objectifs / prérequis / bilan
\newtcolorbox{objectifs}{popbase,colframe=popink,colback=poporangeL,
  before upper={\popboxhead{Objectifs de la leçon}{popink}{}}}
\newtcolorbox{prerequis}{popbase,colframe=popink,colback=popblueL,
  before upper={\popboxhead{Prérequis}{popblue}{}}}
\newtcolorbox{bilan}{popbase,colframe=popink,colback=white,boxrule=2.8pt,
  before upper={\popboxhead{Fiche bilan}{poporange}{L'essentiel à connaître pour l'examen}}}
% Figure encadrée avec légende
\newtcolorbox{popfigure}[1][]{enhanced,breakable=false,colback=white,colframe=popline,boxrule=1.4pt,arc=8pt,
  left=10pt,right=10pt,top=10pt,bottom=8pt,before skip=10pt,after skip=12pt,halign=center,
  lower separated=false,halign lower=center,
  fontlower=\popsemi\fontsize{9}{11.5}\selectfont\color{popmuted},
  #1}
\newcommand{\legende}[1]{\par\vspace{6pt}{\popsemi\fontsize{9}{11.5}\selectfont\color{popmuted}#1}}

% ============================================================
% Signature OnBuch+
% ============================================================
\newcommand{\poplogoinline}[1]{{\poplogo\fontsize{#1}{#1}\selectfont\color{popink}OnBuch\color{poporange}+}}
\newcommand{\signatureonbuch}{%
  \begin{tcolorbox}[enhanced,colback=popink,colframe=popink,boxrule=2.2pt,arc=10pt,
      drop shadow=poporange,left=14pt,right=14pt,top=10pt,bottom=10pt,before skip=0pt,after skip=0pt]
    \tikz[baseline=-0.7ex]\node[circle,fill=popgreen,draw=popcream,line width=1.3pt,minimum size=22pt,inner sep=0]{\popcheck{white}};%
    \hspace{9pt}\begin{minipage}[c]{0.62\linewidth}
      {\popxbold\fontsize{12}{13}\selectfont\color{popcream}Signé\ \textbullet\ {\poplogo OnBuch\color{poporange}+}}\par\vspace{2pt}
      {\raggedright\popsemi\fontsize{8}{10.5}\selectfont\color{popline}\MakeUppercase{Équipe pédagogique OnBuch+}\\\MakeUppercase{Conforme au programme officiel MINESEC}\par}
    \end{minipage}\hfill
    {\poplogo\fontsize{22}{22}\selectfont\color{popcream}OnBuch\color{poporange}+}
  \end{tcolorbox}%
}

% ============================================================
% Page de garde
% #1 titre · #2 matière · #3 niveau · #4 module · #5 n° leçon · #6 durée ·
% #7 description / objectifs (texte ou itemize)
% ============================================================
\newcommand{\pagegarde}[7]{%
  \thispagestyle{empty}
  \newgeometry{a4paper,margin=1.7cm,top=1.6cm,bottom=1.4cm}%
  \begin{tikzpicture}[remember picture,overlay]
    \fill[poporange!18] ([xshift=-3.2cm,yshift=-3.6cm]current page.north east) circle (4.6cm);
    \fill[poppurple!14] ([xshift=2.4cm,yshift=7.6cm]current page.south west) circle (3.8cm);
  \end{tikzpicture}%
  {\poplogo\fontsize{30}{30}\selectfont\color{popink}OnBuch\color{poporange}+}\hfill
  \raisebox{0.6ex}{\poppill{Cours signé \textbullet\ Premium}{popink}{popcream}}\par
  \vspace{1pt}\popsquiggle{poppurple}\par
  \vspace{8pt}
  \begin{tcolorbox}[enhanced,colback=poporange,colframe=popink,boxrule=2.8pt,arc=13pt,
      drop shadow=popink,left=18pt,right=18pt,top=12pt,bottom=13pt,after skip=10pt]
    \poppill{#2 \textbullet\ #3}{popink}{popcream}\hspace{5pt}\poppill{#4}{white}{popink}\par\vspace{8pt}
    {\popdisplay\fontsize{14}{14}\selectfont\color{popink}LEÇON #5}\par\vspace{4pt}
    {\raggedright\hyphenpenalty=10000\exhyphenpenalty=10000\popdisplay\fontsize{25}{27}\selectfont\color{popcream}\MakeUppercase{#1}\par}
    \vspace{8pt}
    {\popxbold\fontsize{9.5}{9.5}\selectfont\color{popink}\MakeUppercase{Durée conseillée : #6}}
  \end{tcolorbox}
  \begin{tcolorbox}[enhanced,colback=white,colframe=popink,boxrule=2.2pt,arc=10pt,
      drop shadow=popink,left=16pt,right=16pt,top=10pt,bottom=10pt,after skip=8pt]
    \poppill{Dans cette leçon}{poppurple}{white}\par\vspace{5pt}
    \popsemi\fontsize{9.3}{12.6}\selectfont\color{popink}#7
  \end{tcolorbox}
  \noindent\begin{minipage}[t]{0.487\linewidth}
    \begin{tcolorbox}[enhanced,colback=popgreen,colframe=popink,boxrule=2.2pt,arc=10pt,
        drop shadow=popink,left=11pt,right=11pt,top=10pt,bottom=10pt,height=3.1cm,valign=center]
      \tcbox[colback=white,colframe=popink,boxrule=1.3pt,arc=5pt,boxsep=0pt,nobeforeafter,
        left=3pt,right=3pt,top=3pt,bottom=3pt,valign=center]{\qrcode[height=1.9cm]{https://play.google.com/store/apps/details?id=com.luvvix.onbuch&hl=fr}}%
      \hspace{8pt}\begin{minipage}[c]{0.5\linewidth}
        {\popdisplay\fontsize{11}{12.5}\selectfont\color{white}\MakeUppercase{Télécharge OnBuch}}\par\vspace{4pt}
        {\raggedright\popsemi\fontsize{8.4}{11}\selectfont\color{white}Gratuit \textbullet\ Google Play\\Tuteur IA, annales, concours, résultats\par}
      \end{minipage}
    \end{tcolorbox}
  \end{minipage}\hfill
  \begin{minipage}[t]{0.487\linewidth}
    \begin{tcolorbox}[enhanced,colback=poppurple,colframe=popink,boxrule=2.2pt,arc=10pt,
        drop shadow=popink,left=11pt,right=11pt,top=10pt,bottom=10pt,height=3.1cm,valign=center]
      \tikz[baseline=-0.7ex]\node[circle,fill=white,draw=popink,line width=1.3pt,minimum size=1.6cm,inner sep=0]{\popdisplay\fontsize{24}{24}\selectfont\color{poppurple}+};%
      \hspace{8pt}\begin{minipage}[c]{0.6\linewidth}
        {\popdisplay\fontsize{11}{12.5}\selectfont\color{white}\MakeUppercase{Passe à OnBuch+}}\par\vspace{4pt}
        {\raggedright\popsemi\fontsize{8.4}{11}\selectfont\color{white}Tous les cours signés, Tuteur IA illimité, fiches PDF — onglet OnBuch+ de l'app\par}
      \end{minipage}
    \end{tcolorbox}
  \end{minipage}
  \vspace{8pt}
  \popcontenu
  \vfill
  \signatureonbuch
  \restoregeometry
  \newpage
}

% Rangée « Ce cours contient » (page de garde)
\newcommand{\poptile}[3]{% #1 couleur #2 gros texte #3 légende
  \begin{tcolorbox}[enhanced,colback=#1,colframe=popink,boxrule=2pt,arc=9pt,shadow={2.5pt}{-2.5pt}{0pt}{popink},
      left=6pt,right=6pt,top=9pt,bottom=9pt,halign=center,valign=center,height=1.75cm,width=\linewidth,nobeforeafter]
    {\popdisplay\fontsize{12.5}{13}\selectfont\color{white}\MakeUppercase{#2}}\par\vspace{3pt}
    {\centering\popsemi\fontsize{7.8}{9.5}\selectfont\color{white}#3\par}
  \end{tcolorbox}}
\newcommand{\popcontenu}{%
  \par\noindent{\popxboldls\fontsize{7.5}{7.5}\selectfont\color{popmuted}CE COURS SIGNÉ CONTIENT}\par\vspace{7pt}
  \noindent\begin{minipage}[t]{0.235\linewidth}\poptile{popblue}{Cours}{complet, illustré, pas à pas}\end{minipage}\hfill
  \begin{minipage}[t]{0.235\linewidth}\poptile{poporange}{Méthodes}{et exercices résolus}\end{minipage}\hfill
  \begin{minipage}[t]{0.235\linewidth}\poptile{poppink}{Exercices}{type examen + corrigés}\end{minipage}\hfill
  \begin{minipage}[t]{0.235\linewidth}\poptile{popgreen}{Fiche bilan}{l'essentiel à retenir}\end{minipage}\par}

% Sommaire
\newtcolorbox{sommaire}{enhanced,colback=white,colframe=popink,boxrule=2.2pt,arc=10pt,
  drop shadow=popink,left=18pt,right=18pt,top=13pt,bottom=13pt,before skip=6pt,after skip=20pt,
  before upper={\poppill{Sommaire}{poppurple}{white}\par\vspace{9pt}\popsemi\fontsize{10.6}{14.5}\selectfont\color{popink}}}

% Signature de fin de cours — poussée tout en bas de la dernière page
\newcommand{\findecours}{%
  \par\vfill\nopagebreak
  \begin{center}\popsquiggle{poporange}\end{center}\vspace{4pt}
  \signatureonbuch
}
\AtBeginDocument{\def\llbracket{\mathopen{[\mkern-3.5mu[}}\def\rrbracket{\mathclose{]\mkern-3.5mu]}}} % intervalles d'entiers (glyphe ⟦ absent de la police)
% Glyphes absents de Fira Math : redéfinis à partir de glyphes disponibles
\AtBeginDocument{%
  \def\setminus{\mathbin{\backslash}}\let\smallsetminus\setminus
  \def\vdots{\mathord{\vbox{\baselineskip3.2pt\lineskiplimit0pt\kern4pt\hbox{.}\hbox{.}\hbox{.}}}}%
  \def\ddots{\mathinner{\mkern1mu\raise6.5pt\hbox{.}\mkern2mu\raise3.6pt\hbox{.}\mkern2mu\raise0.7pt\hbox{.}\mkern1mu}}%
  \def\triangle{\mathord{\scalebox{0.9}{$\symup{\Delta}$}}}%
  \def\nabla{\mathord{\raisebox{\depth}{\scalebox{1}[-1]{$\symup{\Delta}$}}}}%
  \def\checkmark{\popcheck{popdark}}%
  \def\star{\mathbin{\ast}}\def\bigstar{\mathord{\ast}}%
  \def\diamond{\mathbin{\scalebox{0.6}{$\Diamond$}}}%
  \def\bigcup{\mathop{\mathchoice{\raisebox{-0.3em}{\scalebox{1.7}{$\cup$}}}{\scalebox{1.25}{$\cup$}}{\cup}{\cup}}\displaylimits}%
  \def\bigcap{\mathop{\mathchoice{\raisebox{-0.3em}{\scalebox{1.7}{$\cap$}}}{\scalebox{1.25}{$\cap$}}{\cap}{\cap}}\displaylimits}%
  \def\lesseqqgtr{\mathrel{\lessgtr}}\def\bigoplus{\mathop{\oplus}\displaylimits}%
}
\providecommand{\ssi}{si et seulement si} % abréviation fréquente
\AtBeginDocument{\providecommand{\wideparen}{\overparen}} % arc de cercle

\begin{document}

\pagegarde{Lesson 11 — Vocabulary: Words and expressions related to taking part in recreational activities}{Anglais}{1ère A}{Chapter 1 : Family and Social Life}{ENA11}{2 h}{Cette leçon de vocabulaire permet aux élèves de 1ère A d'acquérir et d'employer le lexique des loisirs et des activités récréatives : sports, jeux, sorties, passe-temps. Elle insiste sur les collocations (play/do/go), les faux amis et la réutilisation du lexique dans des phrases, dialogues et courts textes.
\vspace{4pt}
\begin{multicols}{2}\small
\begin{itemize}
\item À la fin de cette leçon, je sais nommer en anglais au moins vingt activités récréatives (sports, jeux, sorties, passe-temps).
\item À la fin de cette leçon, je sais utiliser correctement les verbes play, do et go avec les activités.
\item À la fin de cette leçon, je sais employer les collocations courantes du thème (take part in, join a club, have fun...).
\item À la fin de cette leçon, je sais reconnaître et éviter les faux amis (assist, actually, eventually...).
\item À la fin de cette leçon, je sais former des mots de la même famille (relax → relaxation, amuse → amusement...).
\item À la fin de cette leçon, je sais prononcer correctement les mots clés du thème et respecter leur orthographe.
\end{itemize}\end{multicols}}

\begin{sommaire}
\begin{multicols}{2}
\begin{enumerate}
\item Champs lexicaux : les activités récréatives
\item Collocations : play, do, go et verbes du thème
\item Expressions idiomatiques et phrases utiles
\item Faux amis et pièges des francophones
\item Familles de mots, prononciation et orthographe
\item Le lexique en contexte : texte, dialogue et production
\item Activité d'intégration
\item Méthodes et exercices résolus
\item Exercices
\item Corrigés des exercices
\item Fiche bilan
\end{enumerate}
\end{multicols}
\end{sommaire}

\begin{objectifs}
\begin{itemize}
\item À la fin de cette leçon, je sais nommer en anglais au moins vingt activités récréatives (sports, jeux, sorties, passe-temps).
\item À la fin de cette leçon, je sais utiliser correctement les verbes play, do et go avec les activités.
\item À la fin de cette leçon, je sais employer les collocations courantes du thème (take part in, join a club, have fun...).
\item À la fin de cette leçon, je sais reconnaître et éviter les faux amis (assist, actually, eventually...).
\item À la fin de cette leçon, je sais former des mots de la même famille (relax → relaxation, amuse → amusement...).
\item À la fin de cette leçon, je sais prononcer correctement les mots clés du thème et respecter leur orthographe.
\item À la fin de cette leçon, je sais parler de mes loisirs et interroger un camarade sur les siens dans un dialogue.
\item À la fin de cette leçon, je sais rédiger un court paragraphe décrivant mes activités de week-end.
\end{itemize}
\end{objectifs}

\begin{prerequis}
\begin{itemize}
\item Vocabulaire de base des sports et loisirs vu au collège et en Seconde (football, music, dance...)
\item Présent simple pour exprimer les habitudes (I play, she goes...)
\item Verbes de goût suivis du gérondif (like, enjoy, love + V-ing)
\item Expressions de fréquence (always, often, twice a week...)
\end{itemize}
\end{prerequis}

\newpage

\coursec{Champs lexicaux : les activités récréatives}

\courssub{Situation d'accroche : le week-end arrive}

\begin{textebox}[It's Friday afternoon!]
It's Friday afternoon at GBHS Bamenda. The weekend is coming and students are excited: some will play football on the school field, others will go swimming, watch a movie, or simply hang out with friends. The teacher asks: “What do you do for fun?” But many students answer “I make sport” or “I assist to a match” — mistakes copied from French! Today, we learn the exact English words and expressions to talk about hobbies, sports, games and free-time activities like a native speaker.
\end{textebox}

\textbf{Problématique.}\par
Comment nommer et organiser, en anglais correct, toutes les activités que nous pratiquons pendant notre temps libre ? Pour répondre à la question \eng{What do you do for fun?} (« Qu'est-ce que tu fais pour t'amuser ? »), il faut d'abord maîtriser le \cle{champ lexical} des activités récréatives : les mots, les verbes qui les accompagnent et les pièges à éviter quand on pense en français.

\begin{definition}[Recreational activities]
\eng{Recreational activities} (prononcé « ré-kri-é-cheu-nal ak-ti-vi-tiz ») = \emph{activities you do for pleasure or fun during your free time}. En français : les \textbf{activités récréatives} sont les activités pratiquées pour le plaisir pendant les loisirs. On parle aussi de \eng{leisure activities} ou \eng{free-time activities}.
\end{definition}

\begin{exemplebox}[Deux exemples pour commencer]
\eng{In Cameroon, many young people play football after school.} — « Au Cameroun, beaucoup de jeunes jouent au football après l'école. »\par
\eng{My grandmother enjoys storytelling in the evening.} — « Ma grand-mère aime raconter des histoires le soir. »
\end{exemplebox}

\courssub{Observer avant de mémoriser : un texte de référence}

Lisons d'abord ce court texte. Souligne mentalement tous les mots qui désignent une activité de loisir.

\begin{textebox}[Weekend plans in Yaoundé]
My name is Larissa and I am a student in Form Five in Yaoundé. For me, the weekend is the best part of the week because it is my free time. On Saturday morning, my brothers usually play football at the neighbourhood field, while my sister goes jogging with her friends. I prefer quieter activities: I love reading novels and doing puzzles. Sometimes our whole family plays board games like chess or draughts, and my little cousins enjoy hide-and-seek in the courtyard.

On Saturday afternoon, we often visit friends or go to a party. Last week, our church organised a picnic near Mfoundi River, and next month the school is planning an excursion to the Mefou National Park. In the evening, my father likes watching a match on television, my mother listens to music, and my aunt enjoys storytelling — she knows wonderful tales about animals.

Sunday is more creative. I sing in the church choir, my friend Brice plays the guitar, and his sister is learning traditional dancing. After the service, we usually hang out together, chatting and laughing. On Sunday evening, I relax and watch a film before the new school week begins. Free time is short, but it makes us happy and healthy!
\end{textebox}

\textbf{Glossaire.} \emph{field} (n.) : terrain ; \emph{draughts} (n.) : jeu de dames ; \emph{hide-and-seek} (n.) : cache-cache ; \emph{courtyard} (n.) : cour ; \emph{choir} (n.) : chorale ; \emph{tale} (n.) : conte ; \emph{to hang out} (v.) : traîner, passer du temps ensemble.

\textbf{Questions de compréhension.}
\begin{enumerate}
\item What does Larissa prefer doing on Saturday morning?
\item Where did the church organise a picnic?
\item Who enjoys storytelling in Larissa's family?
\item What does Brice play?
\item What does Larissa do on Sunday evening?
\end{enumerate}

\courssub{Les cinq champs lexicaux des activités récréatives}

Le vocabulaire du texte se range naturellement en cinq familles. Une carte mentale aide à les visualiser et à les mémoriser.

\begin{popfigure}
\begin{tikzpicture}[mindmap, grow cyclic, every node/.style={concept, text=white, align=center}, root concept/.append style={concept color=popblue, font=\bfseries}, level 1/.append style={level distance=3.4cm, sibling angle=72}, level 2/.append style={level distance=2.2cm, sibling angle=45}]
\node[root concept]{Recreational activities}
  child[concept color=poporange]{ node{Sports} child{ node{football} } child{ node{swimming} } child{ node{cycling} } }
  child[concept color=popgreen]{ node{Games} child{ node{chess} } child{ node{cards} } child{ node{puzzles} } }
  child[concept color=poppurple]{ node{Outings} child{ node{cinema} } child{ node{picnic} } child{ node{party} } }
  child[concept color=poppink]{ node{Creative activities} child{ node{singing} } child{ node{drawing} } child{ node{guitar} } }
  child[concept color=popgold]{ node{Relaxation} child{ node{chatting} } child{ node{music} } child{ node{TV} } };
\end{tikzpicture}
\legende{Carte mentale : les cinq familles du vocabulaire des loisirs}
\end{popfigure}

\textbf{Champ 1 — Sports et activités physiques.}

\begin{center}
\begin{tabularx}{\linewidth}{|X|X|X|}
\hline
\rowcolor{popblueL} English & Français & Remarque \\
\hline
football / soccer & football & \eng{soccer} = anglais américain \\
\hline
basketball & basket-ball & un seul mot en anglais \\
\hline
volleyball & volley-ball & un seul mot en anglais \\
\hline
athletics & athlétisme & toujours avec un \emph{-s} final \\
\hline
swimming & natation & attention au double \emph{m} \\
\hline
cycling & cyclisme & du verbe \eng{to cycle} \\
\hline
jogging & jogging, footing & mot identique au français \\
\hline
wrestling & lutte & prononcé « ré-sling » \\
\hline
boxing & boxe & du verbe \eng{to box} \\
\hline
table tennis & tennis de table, ping-pong & deux mots séparés \\
\hline
\end{tabularx}
\end{center}

\textbf{Champ 2 — Jeux et passe-temps.} \eng{board games} (jeux de société), \eng{chess} (échecs), \eng{draughts} (jeu de dames — prononcé « draafts »), \eng{cards} (les cartes), \eng{video games} (jeux vidéo), \eng{puzzles} (puzzles), \eng{hide-and-seek} (cache-cache), \eng{skipping rope} (corde à sauter).

\textbf{Champ 3 — Sorties et divertissements.} \eng{going to the cinema} (aller au cinéma), \eng{watching a match} (regarder un match), \eng{attending a concert} (assister à un concert), \eng{visiting friends} (rendre visite à des amis), \eng{going to a party} (aller à une fête), \eng{a picnic} (un pique-nique), \eng{an excursion} (une excursion, une sortie scolaire).

\textbf{Champ 4 — Activités culturelles et créatives.} \eng{singing} (chanter), \eng{dancing} (danser), \eng{playing an instrument} (jouer d'un instrument : \eng{the drum} le tambour, \eng{the guitar} la guitare), \eng{drawing} (dessiner), \eng{painting} (peindre), \eng{reading novels} (lire des romans), \eng{storytelling} (l'art de raconter des histoires), \eng{drama / acting} (le théâtre, jouer la comédie).

\textbf{Champ 5 — Détente.} \eng{relaxing} (se détendre), \eng{resting} (se reposer), \eng{chatting with friends} (bavarder avec des amis), \eng{listening to music} (écouter de la musique), \eng{watching TV} (regarder la télévision), \eng{hanging out} (traîner, passer du bon temps ensemble).

\courssub{Expressions idiomatiques et phrases utiles}

Au-delà des noms d'activités, certaines expressions toutes faites (\cle{idioms}) et structures reviennent sans cesse pour parler de ses loisirs. Les connaître permet de passer du vocabulaire « scolaire » à l'anglais naturel.

\begin{definition}[Expressions idiomatiques courantes]
\begin{itemize}
\item \eng{to hang out (with someone)} : passer du temps (avec quelqu'un), traîner ensemble (très courant à l'oral). \eng{We just hung out at the mall.} — « On a juste traîné au centre commercial. »
\item \eng{to chill out} / \eng{to chill} : se relaxer, décompresser. \eng{After the exam, I just want to chill.} — « Après l'examen, je veux juste me poser. »
\item \eng{to have a blast} : s'amuser comme un fou, passer un moment génial. \eng{The party was great, we had a blast!} — « La fête était super, on s'est régalés ! »
\item \eng{to catch a movie / a game / a concert} : aller voir un film / un match / un concert (familier). \eng{Let's catch a movie tonight.} — « On va se faire un film ce soir ? »
\item \eng{to shoot some hoops} : faire un peu de basket (lancer des paniers). \eng{He goes to the court to shoot some hoops.} — « Il va au terrain faire un peu de basket. »
\item \eng{to hit the road} : partir, se mettre en route. \eng{It's late, we must hit the road.} — « Il est tard, il faut qu'on parte. »
\end{itemize}
\end{definition}

\begin{exemplebox}[Phrases-types pour parler de ses goûts]
\begin{itemize}
\item \eng{What do you do for fun?} / \eng{How do you spend your free time?} — Questions standard.
\item \eng{I'm into + \emph{-ing} / noun} : \eng{I'm into photography.} / \eng{I'm into playing chess.} — « Je m'intéresse à… / J'aime bien… »
\item \eng{I'm keen on + \emph{-ing} / noun} (plus soutenu) : \eng{She's keen on hiking.} — « Elle aime beaucoup la randonnée. »
\item \eng{I enjoy + \emph{-ing}} : \eng{I enjoy reading novels.} — « J'aime lire des romans. »
\item \eng{I can't stand + \emph{-ing}} : \eng{I can't stand watching TV all day.} — « Je ne supporte pas regarder la télé toute la journée. »
\item \eng{I'm a fan of ...} : \eng{I'm a big fan of football.} — « Je suis fan de foot. »
\end{itemize}
\end{exemplebox}

\courssub{Familles de mots, prononciation et orthographe}

Construire son lexique, c'est aussi savoir dériver les mots. Voici les principales familles du thème ; attention à l'accentuation (soulignée) et à l'orthographe.

\begin{propriete}[Dérivation nominale et adjectivale]
\begin{center}
\begin{tabular}{|l|l|l|l|l|}
\hline
\rowcolor{popblueL} \textbf{Verbe} & \textbf{Nom (action/état)} & \textbf{Nom (agent)} & \textbf{Adjectif} & \textbf{Adverbe} \\
\hline
recreate & recreation & — & recreational & recreationally \\
\hline
— & leisure & — & leisurely & leisurely \\
\hline
sport & sport & sportsman / sportswoman & sporty & sportily \\
\hline
compete & competition & competitor & competitive & competitively \\
\hline
create & creativity / creation & creator & creative & creatively \\
\hline
relax & relaxation & — & relaxed / relaxing & relaxedly \\
\hline
enjoy & enjoyment & — & enjoyable & enjoyably \\
\hline
entertain & entertainment & entertainer & entertaining & entertainingly \\
\hline
\end{tabular}
\end{center}
\textbf{Points d'orthographe et de prononciation clés :}
\begin{itemize}
\item \eng{recreation} / \eng{recreational} : accent sur \textbf{é} (\eng{recr\textbf{e}ation}, \eng{recr\textbf{e}ational}) — « ré-kri-é-cheun / ré-kri-é-cheu-nal ».
\item \eng{leisure} : prononcé « \textbf{lé}-jeur » (pas « lai-jeur ») ; indénombrable, pas de \emph{-s}.
\item \eng{competitive} : accent sur \textbf{é} (\eng{comp\textbf{e}titive}) — « com-\textbf{pé}-ti-tiv ».
\item \eng{creative} : accent sur \textbf{é} (\eng{cr\textbf{e}ative}) — « \textbf{cré}-é-tiv ».
\item \eng{swimming} : double \emph{m} (comme \eng{running}, \eng{beginning}).
\item \eng{basketball}, \eng{volleyball}, \eng{football} : mots composés \textbf{solidaires} (un seul mot).
\item \eng{athletics}, \eng{draughts}, \eng{cards} : noms en \emph{-s} mais souvent singuliers dans l'usage (\eng{Athletics is my favourite sport}).
\end{itemize}
\end{propriete}

\begin{aretenir}[Le point de départ : play, do, go]
Beaucoup de noms d'activités s'emploient avec un verbe particulier : \eng{play football}, \eng{play chess}, \eng{play the guitar} ; \eng{do puzzles}, \eng{do athletics} ; \eng{go swimming}, \eng{go jogging}, \eng{go cycling}, \eng{go to the cinema}. Nous étudierons ces \cle{collocations} en détail dans la prochaine section ; retiens déjà que le choix du verbe ne se traduit jamais mot à mot.
\end{aretenir}

\courssub{Leisure, free time, spare time : trois synonymes de « temps libre »}

\begin{definition}[Trois façons de dire « temps libre »]
\begin{itemize}
\item \eng{free time} : l'expression la plus courante et la plus neutre. \eng{What do you do in your free time?} — « Que fais-tu pendant ton temps libre ? »
\item \eng{spare time} : synonyme de \eng{free time}, très fréquent à l'oral. \eng{I read novels in my spare time.} — « Je lis des romans pendant mes loisirs. »
\item \eng{leisure} (prononcé « lé-jeur ») : plus soutenu ; souvent dans \eng{leisure activities} (activités de loisirs) ou \eng{at leisure} (à loisir, sans se presser).
\end{itemize}
\end{definition}

\begin{attention}[« Loisirs » ne se dit pas « leisures »]
\eng{Leisure} est \textbf{indénombrable} : on ne dit jamais \emph{leisures}. Pour parler des activités, dites \eng{leisure activities} ou \eng{hobbies} (passe-temps). De même, évitez le calque \emph{I make sport} : l'anglais dit \eng{I do sport} ou \eng{I play sports}. Et \eng{to attend a match} signifie « assister à un match » (être présent), jamais « attendre un match » — \eng{to wait for a match}.
\end{attention}

\begin{attention}[Watch ou see ?]
On dit \eng{watch TV} et \eng{watch a match} (regarder attentivement, suivre quelque chose qui bouge), mais \eng{see a film} au cinéma. Ne confondez pas \eng{watch} (regarder avec attention) et \eng{see} (apercevoir, voir le résultat) : \eng{I watched the match and I saw three goals.} — « J'ai regardé le match et j'ai vu trois buts. »
\end{attention}

\begin{savaistu}[Football ou soccer ?]
En Grande-Bretagne, on dit \eng{football} ; aux États-Unis, on dit \eng{soccer}, car le mot \eng{football} y désigne le football américain. Au Cameroun anglophone, comme en Grande-Bretagne, on dit \eng{football} — et le pays entier connaît les Lions Indomptables !
\end{savaistu}

\begin{exoresolu}[Classer le vocabulaire par champ lexical]
Énoncé : range les mots suivants dans les cinq champs lexicaux (Sports / Games / Outings / Creative activities / Relaxation) : \eng{wrestling, chess, a picnic, painting, chatting, table tennis, cards, a concert, storytelling, hanging out}.
\tcblower
Solution détaillée :
\begin{itemize}
\item \textbf{Sports} : \eng{wrestling} (lutte), \eng{table tennis} (tennis de table) — ce sont des activités physiques codifiées.
\item \textbf{Games} : \eng{chess} (échecs), \eng{cards} (cartes) — ce sont des jeux avec des règles.
\item \textbf{Outings} : \eng{a picnic} (pique-nique), \eng{a concert} (concert) — ce sont des sorties, des événements auxquels on se rend.
\item \textbf{Creative activities} : \eng{painting} (peinture), \eng{storytelling} (raconter des histoires) — elles produisent ou expriment quelque chose.
\item \textbf{Relaxation} : \eng{chatting} (bavarder), \eng{hanging out} (traîner entre amis) — ce sont des moments de détente sans règles ni effort.
\end{itemize}
Astuce : demande-toi toujours « est-ce un sport, un jeu, une sortie, une création ou un moment de repos ? »
\end{exoresolu}

\begin{exercice}[À toi de jouer]
\begin{enumerate}
\item Traduis en anglais : « Le samedi, je joue au basket-ball et le soir je regarde un film avec mes amis. »
\item Complète avec \eng{watch} ou \eng{see} : (a) We \ldots\ a film at the cinema last night. (b) My father is \ldots\ a football match on TV.
\item Trouve le mot du champ lexical qui correspond : un jeu de société avec des pions noirs et blancs sur un damier (\ldots) ; une sortie scolaire organisée (\ldots) ; l'activité qui consiste à raconter des contes (\ldots).
\item Réponds à la question du professeur en trois phrases complètes : \eng{What do you do in your free time?}
\end{enumerate}
\end{exercice}

\begin{corrige}[Exercice « À toi de jouer »]
\begin{enumerate}
\item \eng{On Saturday, I play basketball and in the evening I watch a film with my friends.} (ou \eng{see a film} si c'est au cinéma).
\item (a) \eng{saw} — au cinéma, on emploie \eng{see} ; (b) \eng{watching} — on suit le match avec attention.
\item \eng{draughts} (jeu de dames) ; \eng{an excursion} ; \eng{storytelling}.
\item Modèle de réponse : \eng{In my free time, I play football with my friends. I also enjoy listening to music and reading novels. At the weekend, I sometimes visit my grandmother in the village.}
\end{enumerate}
\end{corrige}

\begin{aretenir}[Ce qu'il faut retenir de cette section]
\begin{itemize}
\item Les \cle{activités récréatives} se classent en cinq champs : sports, jeux, sorties, activités créatives, détente.
\item \eng{Free time}, \eng{spare time} et \eng{leisure} signifient « temps libre », mais \eng{leisure} est indénombrable.
\item On \eng{watches} la télévision et les matchs, mais on \eng{sees} un film au cinéma.
\item Méfie-toi des calques du français : jamais \emph{make sport}, jamais \emph{assist to a match}.
\item Utilise \eng{I'm into / keen on / enjoy + -ing} pour exprimer tes goûts naturellement.
\item Maîtrise les familles de mots : \eng{create $\to$ creative}, \eng{compete $\to$ competitive}, \eng{leisure $\to$ leisurely}.
\end{itemize}
\end{aretenir}

\coursec{Vocabulary : Taking part in recreational activities}

\courssub{Observer d'abord : trois amis parlent de leurs loisirs}

\begin{textebox}[After school in Yaoundé]
Every afternoon, the students of Government High School leave their classrooms and rush to their favourite activities. Alain and his friends \textbf{play football} on the school pitch, while the girls' team \textbf{plays handball} near the gym. In the music room, Clarisse \textbf{plays the guitar} and her brother \textbf{plays the drums}. On Saturday mornings, many families \textbf{go swimming} at the municipal pool, and some joggers \textbf{go jogging} around the Ahmadou Ahidjo Stadium before sunrise. During the holidays, young people \textbf{go camping} near Buea or \textbf{go fishing} on the Wouri River. Meanwhile, Sandra \textbf{does judo} at a club in Nlongkak, and her mother \textbf{does yoga} every evening to relax. Under the big mango tree, two old men \textbf{play draughts} and laugh all afternoon. Whatever they choose, all these people \textbf{take part in} recreational activities and \textbf{have fun} together.
\end{textebox}

\textbf{Glossaire :} \emph{pitch} (n., terrain de sport) ; \emph{drums} (n., batterie) ; \emph{joggers} (n., joggeurs) ; \emph{draughts} (n., jeu de dames) ; \emph{whatever} (conj., quoi que) ; \emph{municipal} (adj., communal).

Observez les mots en gras : devant \eng{football}, \eng{handball}, \eng{the guitar}, \eng{draughts}, on trouve \eng{play} ; devant \eng{swimming}, \eng{jogging}, \eng{camping}, \eng{fishing}, on trouve \eng{go} ; devant \eng{judo}, \eng{yoga}, on trouve \eng{do}. Ce n'est pas un hasard : c'est une règle fondamentale de l'anglais.

\courssub{La règle des trois verbes : play, go, do}

\begin{propriete}[PLAY + jeux, sports avec ballon et instruments]
On utilise \cle{play} devant :
\begin{itemize}
\item les sports qui se jouent avec un ballon ou une balle : \eng{play football}, \eng{play basketball}, \eng{play volleyball}, \eng{play tennis}, \eng{play handball}, \eng{play rugby} ;
\item les jeux de société ou d'esprit : \eng{play chess} (jouer aux échecs), \eng{play cards} (jouer aux cartes), \eng{play draughts} (jouer aux dames), \eng{play video games}, \eng{play Scrabble} ;
\item les instruments de musique, \textbf{avec l'article \eng{the}} : \eng{play the guitar}, \eng{play the piano}, \eng{play the drums}, \eng{play the violin}, \eng{play the flute}.
\end{itemize}
\textbf{Attention :} pas d'article devant les noms de sports et de jeux (\eng{play football}, jamais \eng{play the football}), mais l'article \eng{the} est \textbf{obligatoire} devant les instruments de musique.
\end{propriete}

\begin{propriete}[GO + activité en -ing (mouvement, plein air)]
On utilise \cle{go} suivi d'un verbe en \cle{-ing} pour les activités de déplacement ou de plein air que l'on « va » pratiquer : \eng{go swimming} (aller nager), \eng{go jogging}, \eng{go cycling}, \eng{go fishing}, \eng{go shopping}, \eng{go dancing}, \eng{go camping}, \eng{go hiking}, \eng{go hunting}, \eng{go skiing}, \eng{go sailing}.
\\Forme : \eng{go} + verbe + \eng{-ing}. Le verbe en \eng{-ing} ne change jamais ; c'est \eng{go} qui se conjugue : \eng{I go}, \eng{she goes}, \eng{they went}, \eng{we have gone}, \eng{we are going}.
\end{propriete}

\begin{propriete}[DO + sports individuels, arts martiaux et exercices]
On utilise \cle{do} devant les disciplines individuelles, les arts martiaux, la gymnastique et les activités de remise en forme : \eng{do athletics} (faire de l'athlétisme), \eng{do gymnastics}, \eng{do judo}, \cle{do karate} (ou \eng{karate-do}), \eng{do yoga}, \eng{do exercise} (faire de l'exercice, indénombrable), \eng{do aerobics}, \eng{do weightlifting} (haltérophilie), \eng{do sport} (faire du sport, en général).
\\Idée directrice : pas de ballon, pas de notion de « aller quelque part » : c'est une discipline que l'on pratique sur place.
\end{propriete}

\begin{exemplebox}[Exemples traduits]
\eng{We go swimming every Saturday.} « Nous allons nager chaque samedi. »\\
\eng{She does judo at the club.} « Elle fait du judo au club. »\\
\eng{They play draughts under the tree.} « Ils jouent aux dames sous l'arbre. »\\
\eng{My father plays the guitar in the evening.} « Mon père joue de la guitare le soir. »\\
\eng{The students play basketball after classes.} « Les élèves jouent au basket après les cours. »\\
\eng{Did you go fishing last Sunday?} « Es-tu allé à la pêche dimanche dernier ? »\\
\eng{I do yoga to relax.} « Je fais du yoga pour me détendre. »\\
\eng{He does exercise every morning before work.} « Il fait de l'exercice chaque matin avant le travail. »
\end{exemplebox}

\begin{methode}[Comment choisir le bon verbe ?]
Face à une activité, posez-vous trois questions dans l'ordre :
\begin{enumerate}
\item Y a-t-il un \textbf{ballon} / une balle, ou est-ce un \textbf{jeu} ou un \textbf{instrument de musique} ? $\rightarrow$ utilisez \cle{play}.
\item L'activité est-elle un \textbf{mouvement} exprimé par un verbe en \textbf{-ing} (nager, courir, camper, skier) ? $\rightarrow$ utilisez \cle{go} + \eng{-ing}.
\item Est-ce une \textbf{discipline individuelle} (art martial, gymnastique, yoga, athlétisme) ou le sport en général ? $\rightarrow$ utilisez \cle{do}.
\end{enumerate}
Exemple : \eng{karate} $\rightarrow$ pas de ballon, pas de \eng{-ing} $\rightarrow$ \eng{do karate}. \quad \eng{tennis} $\rightarrow$ balle $\rightarrow$ \cle{play tennis}.
\end{methode}

\begin{popfigure}
\begin{center}
\begin{tikzpicture}[
  box/.style={draw=popline, rounded corners, align=center, minimum width=4.5cm, text width=4.2cm, font=\small, inner sep=4pt},
  head/.style={font=\bfseries\large, align=center}
]
\node[head, popblue] (p) at (0,0) {PLAY};
\node[box, fill=popblueL, below=0.2cm of p, align=center]{
  \textbf{Ball games:} football, basketball, volleyball, tennis, handball, rugby\\[4pt]
  \textbf{Games:} chess, cards, draughts, Scrabble, video games\\[4pt]
  \textbf{Instruments (+ the):} the guitar, the piano, the drums, the violin
};
\node[head, poporange] (g) at (5.5,0) {GO};
\node[box, fill=poporangeL, below=0.2cm of g, align=center]{
  \textbf{Verb + -ing (movement, outdoors):}\\[2pt]
  swimming, jogging, cycling, fishing, shopping,\\
  dancing, camping, hiking, hunting, skiing, sailing
};
\node[head, popgreen] (d) at (11,0) {DO};
\node[box, fill=popgreenL, below=0.2cm of d, align=center]{
  \textbf{Individual sports, martial arts, fitness:}\\[2pt]
  athletics, gymnastics, judo, karate, yoga,\\
  exercise, aerobics, weightlifting, sport (general)
};
\end{tikzpicture}
\end{center}
\legende{Les trois verbes des activités récréatives : play, go, do.}
\end{popfigure}

\courssub{Tableau récapitulatif}

\begin{center}
\begin{tabularx}{\linewidth}{|l|X|X|}\hline
\rowcolor{popblueL} \textbf{Verbe} & \textbf{Emploi} & \textbf{Exemples traduits} \\ \hline
PLAY & sports avec ballon, jeux, instruments (+ \eng{the}) & \eng{They play football.} « Ils jouent au football. » \\ 
& & \eng{She plays the piano.} « Elle joue du piano. » \\ \hline
GO & activités en -ing (mouvement, plein air) & \eng{We go jogging on Sundays.} « Nous allons courir le dimanche. » \\
& & \eng{He went camping in Buea.} « Il est allé camper à Buea. » \\ \hline
DO & sports individuels, arts martiaux, exercice & \eng{I do karate twice a week.} « Je fais du karaté deux fois par semaine. » \\
& & \eng{She does exercise every morning.} « Elle fait de l'exercice chaque matin. » \\ \hline
\end{tabularx}
\end{center}

Comparaison avec le français : le français utilise presque toujours « faire » ou « jouer » ; l'anglais, lui, impose trois verbes différents. C'est pourquoi la traduction mot à mot est dangereuse.

\begin{center}
\begin{tabularx}{\linewidth}{|X|X|X|}\hline
\rowcolor{popblueL} Français & English & Remarque \\ \hline
faire du sport & do sport / play sports & jamais \eng{make sport} \\ \hline
jouer au football & play football & pas d'article \\ \hline
jouer du piano & play the piano & article \eng{the} obligatoire \\ \hline
faire du judo / du karaté & do judo / do karate & discipline individuelle \\ \hline
aller nager / faire de la natation & go swimming & \eng{go} + -ing \\ \hline
faire les courses / du shopping & go shopping & \eng{go} + -ing \\ \hline
faire de la gymnastique & do gymnastics & pluriel en anglais \\ \hline
faire de l'athlétisme & do athletics & pluriel en anglais \\ \hline
\end{tabularx}
\end{center}

\courssub{Collocations verbales clés du thème}

Au-delà de \eng{play}/\eng{go}/\eng{do}, certaines combinaisons verbe + nom reviennent constamment dans les sujets d'examen (Bac, contrôles) :

\begin{itemize}
\item \cle{take part in} + activité / nom : participer à. \eng{All the students take part in the tournament.} « Tous les élèves participent au tournoi. »
\item \cle{join a club} / \cle{join a team} : adhérer à un club, s'inscrire dans une équipe. \eng{I joined the music club last year.} « Je me suis inscrit au club de musique l'an dernier. »
\item \cle{train} / \cle{practise} (GB) / \cle{practice} (US) : s'entraîner. \eng{The team trains every Tuesday.} « L'équipe s'entraîne chaque mardi. »
\item \cle{win} / \cle{lose} / \cle{draw} a match : gagner / perdre / faire match nul. \eng{Our school won the match 2--0.} « Notre école a gagné le match 2 à 0. »
\item \cle{score a goal} / \cle{score a point} : marquer un but / un point. \eng{The striker scored two goals.} « L'attaquant a marqué deux buts. »
\item \cle{support} / \cle{cheer for} a team : soutenir une équipe. \eng{Many Cameroonians support the Indomitable Lions.} « Beaucoup de Camerounais soutiennent les Lions Indomptables. »
\item \cle{have fun} : s'amuser. \eng{We had fun at the picnic.} « Nous nous sommes amusés au pique-nique. »
\item \cle{enjoy oneself} (pronom réfléchi obligatoire) : bien s'amuser. \eng{They enjoyed themselves at the party.} « Ils se sont bien amusés à la fête. »
\item \cle{spend time} (on / \eng{-ing}) : passer du temps. \eng{She spends her free time reading.} « Elle passe son temps libre à lire. »
\item \cle{waste time} (on / \eng{-ing}) : perdre son temps. \eng{Don't waste time on video games all day.} « Ne perds pas ton temps sur les jeux vidéo toute la journée. »
\end{itemize}

\courssub{Collocations nominales}

\begin{definition}[Les noms à connaître]
\cle{a leisure activity} : une activité de loisir ; \cle{a hobby} : un passe-temps favori ; \cle{a pastime} : un passe-temps (plus littéraire) ; \cle{a football pitch} (GB) / \cle{a football field} (US) : un terrain de football ; \cle{a swimming pool} : une piscine ; \cle{a gym} / \cle{a gymnasium} : un gymnase ; \cle{a match} : un match ; \cle{a tournament} : un tournoi ; \cle{a fan} / \cle{a supporter} : un supporter ; \cle{a teammate} : un coéquipier ; \cle{an opponent} : un adversaire ; \cle{a referee} : un arbitre ; \cle{a spectator} : un spectateur.
\end{definition}

\begin{exemplebox}[En contexte]
\eng{Reading is my favourite pastime.} « La lecture est mon passe-temps préféré. »\\
\eng{The tournament takes place on the school pitch.} « Le tournoi a lieu sur le terrain de l'école. »\\
\eng{Thousands of fans watched the match.} « Des milliers de supporters ont regardé le match. »\\
\eng{Our teammates celebrated the victory.} « Nos coéquipiers ont fêté la victoire. »
\end{exemplebox}

\courssub{Faux amis et pièges des francophones}

\begin{attention}[Faux amis : \eng{sport} / \eng{exercise} / \eng{training}]
\begin{itemize}
\item \textbf{\eng{sport}} (n.) = le sport (en général) ou un sport. \eng{Football is a popular sport.} Ne pas confondre avec \eng{sportswear} (vêtements de sport).
\item \textbf{\eng{exercise}} (n. indénombrable) = l'exercice physique. \eng{Do exercise.} Mais \textbf{\eng{exercises}} (n. pluriel) = des exercices (scolaires, de musculation). \eng{Do your homework exercises.}
\item \textbf{\eng{training}} (n.) = l'entraînement (processus). \eng{Training is hard.} Le nom de la séance est \eng{a training session} ou \eng{a practice} (US).
\item \textbf{\eng{stage}} (n.) = un stage (théâtre, formation) \textbf{MAIS} \eng{internship} (stage professionnel en entreprise) / \eng{work placement} (GB). Ne pas dire \eng{I do a stage in a company} $\rightarrow$ \eng{I do an internship.}
\end{itemize}
\end{attention}

\begin{attention}[Erreurs fréquentes : verbes et structures]
\begin{itemize}
\item \eng{I make sport.} $\times$ $\rightarrow$ \eng{I do sport.} ou \eng{I play sports.} $\checkmark$ (\eng{make} = fabriquer).
\item \eng{I play judo.} $\times$ $\rightarrow$ \eng{I do judo.} $\checkmark$ (pas de ballon !).
\item \eng{I go to swim every day.} possible mais lourd $\rightarrow$ \eng{I go swimming every day.} $\checkmark$ (forme idiomatique).
\item \eng{She makes gymnastics.} $\times$ $\rightarrow$ \eng{She does gymnastics.} $\checkmark$.
\item \textbf{\eng{practise} (verbe, GB) vs \eng{practice} (nom, GB).} \eng{I practise every day.} / \eng{Practice makes perfect.} En anglais américain : \eng{practice} pour les deux.
\item \eng{Play the football.} $\times$ $\rightarrow$ \eng{Play football.} $\checkmark$ (pas d'article pour les sports). Mais \eng{Play the guitar.} $\checkmark$ (article obligatoire pour instruments).
\item \eng{Do a match.} $\times$ $\rightarrow$ \eng{Play a match.} $\checkmark$.
\end{itemize}
\end{attention}

\courssub{Familles de mots, prononciation et orthographe}

\begin{propriete}[Dérivation : noms, verbes, adjectifs]
Maîtriser les familles de mots permet d'enrichir ses rédactions (Writing) et de comprendre des textes complexes (Reading).
\end{propriete}

\begin{center}
\begin{tabular}{|l|l|l|l|}\hline
\rowcolor{popblueL} \textbf{Verbe} & \textbf{Nom (action/état)} & \textbf{Nom (agent)} & \textbf{Adjectif / Participe} \\ \hline
\eng{recreate} & \eng{recreation} & -- & \eng{recreational} \\ \hline
\eng{compete} & \eng{competition} & \eng{competitor} & \eng{competitive} \\ \hline
\eng{entertain} & \eng{entertainment} & \eng{entertainer} & \eng{entertaining} / \eng{entertained} \\ \hline
\eng{relax} & \eng{relaxation} & -- & \eng{relaxed} / \eng{relaxing} \\ \hline
\eng{act} & \eng{activity} & \eng{actor} / \eng{actress} & \eng{active} / \eng{inactive} \\ \hline
\eng{exercise} & \eng{exercise} & \eng{exerciser} (rare) & \eng{exercising} \\ \hline
\eng{support} & \eng{support} & \eng{supporter} / \eng{fan} & \eng{supportive} \\ \hline
\eng{train} & \eng{training} & \eng{trainer} / \eng{coach} & \eng{trained} \\ \hline
\eng{win} & \eng{victory} / \eng{win} & \eng{winner} & \eng{winning} \\ \hline
\eng{lose} & \eng{loss} / \eng{defeat} & \eng{loser} & \eng{lost} \\ \hline
\end{tabular}
\end{center}

\begin{aretenir}[Point d'orthographe : \eng{-ise} vs \eng{-ize}]
En anglais britannique (norme au Cameroun), on écrit \textbf{\eng{-ise}} pour la plupart des verbes : \eng{organise}, \eng{realise}, \eng{recognise}, \eng{practise} (verbe). L'anglais américain utilise \textbf{\eng{-ize}} (\eng{organize}, \eng{practice}). \textbf{Consigne :} restez cohérent dans votre copie (britannique recommandé).
\end{aretenir}

\begin{propriete}[Prononciation : pièges sonores]
\begin{itemize}
\item \textbf{\eng{leisure}} /ˈleʒə/ (GB) « lé-jeur » — \emph{pas} « lai-jour ». Le \eng{ei} se prononce /ʒ/ (comme \eng{measure}).
\item \textbf{\eng{recreational}} /ˌrekriˈeɪʃənl/ « ré-kri-é-ï-cho-nal » — accent sur le \textbf{3\textsuperscript{e}} syllabe (\eng{-eɪ-}).
\item \textbf{\eng{guitar}} /ɡɪˈtɑː/ « guittare » — le \eng{u} est \textbf{muet} après \eng{g}.
\item \textbf{\eng{draughts}} /drɑːfts/ « draafts » — \eng{augh} = /ɑːf/ (comme \eng{laugh}, \eng{draft}).
\item \textbf{\eng{yoga}} /ˈjəʊɡə/ « yo-ga » — \eng{y} initial = /j/ (yod).
\item \textbf{\eng{judo}} /ˈdʒuːdəʊ/ « djou-do » — \eng{j} = /dʒ/.
\item \textbf{\eng{exercise}} /ˈeksəsaɪz/ « èk-se-sa-ïz » — le \eng{c} final se prononce /s/ (doux), le \eng{ise} = /aɪz/.
\item \textbf{\eng{athlete}} /ˈæθliːt/ « a-thlite » — \eng{th} = /θ/ ; \textbf{\eng{athletics}} /æθˈletɪks/ — accent sur \textbf{2\textsuperscript{e}} syllabe.
\end{itemize}
\end{propriete}

\courssub{Le lexique en contexte : dialogue et production}

\begin{experience}[Jeu de rôle : « Planning the Weekend »]
\textbf{Consigne :} En binôme. L'élève A propose des activités pour le week-end en utilisant \eng{play}, \eng{go}, \eng{do} et les collocations vues. L'élève B accepte, refuse ou propose une alternative. Utilisez le vocabulaire de la leçon (\eng{join}, \eng{train}, \eng{supporter}, \eng{have fun}, etc.).
\begin{itemize}
\item \textbf{A :} \eng{What do you want to do this weekend?}
\item \textbf{B :} \eng{I'd like to go swimming at the pool. Do you want to join me?}
\item \textbf{A :} \eng{I can't swim well. I prefer to play football with the team. We have a match on Sunday.}
\item \textbf{B :} \eng{Oh, I'll come and support you then!}
\end{itemize}
\textbf{Variante :} Échangez les rôles. Ajoutez des contraintes (météo, argent, devoirs).
\end{experience}

\begin{textebox}[Dialogue : At the Sports Club Reception]
\textbf{Receptionist :} Good morning! How can I help you?
\textbf{Paul :} Good morning. I'd like to \textbf{join the club}. I want to \textbf{do karate} and \textbf{go swimming}.
\textbf{Receptionist :} Great. The karate class is on Tuesdays and Thursdays at 6 p.m. The swimming pool is open every day until 8 p.m.
\textbf{Paul :} Perfect. I also \textbf{play the guitar}. Is there a music room?
\textbf{Receptionist :} Yes, there is. You can \textbf{practise} there on Saturdays. Do you \textbf{play any sports}?
\textbf{Paul :} I \textbf{play basketball} with my friends sometimes. We \textbf{train} together on weekends.
\textbf{Receptionist :} Excellent. \textbf{Take part in} our tournament next month!
\textbf{Paul :} I will. I want to \textbf{win a match} for once!
\textbf{Receptionist :} \textbf{Have fun} and \textbf{enjoy yourself}!
\end{textebox}

\courssub{Modèle de production écrite : « My favourite recreational activity »}

\begin{methode}[Rédiger un paragraphe descriptif (100--120 mots)]
\begin{enumerate}
\item \textbf{Accroche (Topic sentence) :} Présentez l'activité et le verbe correct. \eng{My favourite recreational activity is playing basketball.}
\item \textbf{Fréquence et lieu :} \eng{I play basketball three times a week at the school pitch.}
\item \textbf{Compagnons / Contexte :} \eng{I usually play with my classmates after school. We are a good team.}
\item \textbf{Pourquoi ? (Bienfaits) :} \eng{It helps me relax and stay fit. I enjoy the competition and the team spirit.}
\item \textbf{Projet / Objectif :} \eng{Next year, I want to join the regional tournament and win a trophy.}
\item \textbf{Conclusion :} \eng{In short, basketball is more than a game for me; it is a passion.}
\end{enumerate}
\end{methode}

\begin{exemplebox}[Modèle de rédaction]
\eng{My favourite recreational activity is playing basketball. I play basketball three times a week at the school pitch with my classmates. We are a good team and we train hard every Saturday morning. I chose this sport because it helps me relax after classes and stay fit. I enjoy the competition and the team spirit; we always support each other. Last year, we took part in the inter-school tournament but we lost the final. This year, my goal is to win the trophy. In short, basketball is more than a game for me; it is a passion that teaches me discipline and friendship.}
\end{exemplebox}

\courssub{Exercice résolu et pièges à éviter}

\begin{exoresolu}[Complétez avec la forme correcte de \eng{play}, \eng{go} ou \eng{do}]
Énoncé — Complétez avec la forme correcte de \eng{play}, \eng{go} ou \eng{do} :\\
1. My sister \ldots\ judo at the club in Douala.\\
2. They \ldots\ swimming every Saturday morning.\\
3. The boys \ldots\ football after school.\\
4. My grandfather \ldots\ chess with his friends.\\
5. We \ldots\ camping near Kribi during the last holidays.\\
6. She \ldots\ the piano very well.\\
7. Look! That man \ldots\ jogging in the street.\\
8. We must \ldots\ exercise to stay healthy.
\tcblower
Solution détaillée :\\
1. \eng{does judo} — art martial, discipline individuelle $\rightarrow$ \eng{do} ; 3\textsuperscript{e} personne sg : \eng{does}.\\
2. \eng{go swimming} — activité en \eng{-ing} $\rightarrow$ \eng{go} (présent habituel).\\
3. \eng{play football} — sport avec ballon $\rightarrow$ \eng{play} (présent pluriel).\\
4. \eng{plays chess} — jeu $\rightarrow$ \eng{play} ; 3\textsuperscript{e} personne sg : \eng{plays}.\\
5. \eng{went camping} — activité en \eng{-ing} au passé (\eng{during the last holidays}) $\rightarrow$ passé de \eng{go} : \eng{went}.\\
6. \eng{plays the piano} — instrument $\rightarrow$ \eng{play} + \eng{the} ; 3\textsuperscript{e} personne sg : \eng{plays}.\\
7. \eng{is jogging} — action en cours (\eng{Look!}) $\rightarrow$ \eng{be} + \eng{-ing} (présent continu). \textbf{Piège :} ce n'est pas \eng{go} ici mais \eng{be} + \eng{-ing}.\\
8. \eng{do exercise} — activité physique générale $\rightarrow$ \eng{do} (infinitif après \eng{must}).
\end{exoresolu}

\begin{savaistu}[Au Bac : la question des loisirs]
À l'épreuve orale comme à l'écrit, la question classique est : \eng{What do you do in your free time?} « Que fais-tu pendant ton temps libre ? ». Remarquez la structure : auxiliaire \eng{do/does} + sujet + verbe de base, \textbf{sans -ing} sur le verbe principal. Répondez avec les collocations de cette leçon : \eng{In my free time, I play football with my friends and I go swimming on Saturdays.} Une réponse riche en collocations correctes (\eng{play}, \eng{go}, \eng{do}, \eng{join}, \eng{train}, \eng{support}) est toujours mieux notée !
\end{savaistu}

\begin{exercice}[À vous de jouer]
1. Traduisez en anglais : « Le week-end, mon frère fait du karaté et moi je vais à la pêche. »\\
2. Trouvez l'erreur et corrigez la phrase : \eng{She makes gymnastics and plays swimming.}\\
3. Complétez avec le bon verbe (\eng{play}, \eng{go}, \eng{do}, \eng{join}, \eng{train}, \eng{support}) à la bonne forme :\\
   a) The Indomitable Lions \ldots\ hard for the next CAN.\\
   b) I want to \ldots\ the chess club next term.\\
   c) My father \ldots\ cycling every Sunday morning.\\
   d) We \ldots\ our local team, Canon of Yaoundé.\\
4. Rédigez un paragraphe de 80--100 mots sur votre activité récréative préférée en suivant le modèle de la méthode (accroche, fréquence, lieu, compagnons, bienfaits, projet).
\end{exercice}

\begin{corrige}[Exercice « À vous de jouer »]
1. \eng{At the weekend, my brother does karate and I go fishing.} (BrE : \eng{At the weekend} ; \eng{do karate}, \eng{go fishing}).\\
2. \eng{She does gymnastics and goes swimming.} (\eng{make} $\rightarrow$ \eng{do} ; \eng{play} + \eng{-ing} impossible $\rightarrow$ \eng{go swimming}).\\
3. a) \eng{train} (présent pluriel) ou \eng{are training} (en ce moment).\\
   b) \eng{join} (infinitif après \eng{want to}).\\
   c) \eng{goes cycling} (3\textsuperscript{e} pers. sg + \eng{go} + \eng{-ing}).\\
   d) \eng{support} (présent pluriel).\\
4. \textbf{Production attendue :} Vérifiez l'emploi correct de \eng{play}/\eng{go}/\eng{do}, la cohérence des temps (présent habituel), le vocabulaire précis (\eng{pitch}, \eng{teammates}, \eng{train}, \eng{win}, \eng{enjoy}) et la structure paragraphe (connecteurs : \eng{usually}, \eng{because}, \eng{also}, \eng{in short}).
\end{corrige}

\coursec{Expressions idiomatiques et phrases utiles}

Dans la section précédente, nous avons appris à construire correctement le nom des activités avec les verbes \eng{play}, \eng{do} et \eng{go} : \eng{play football}, \eng{do gymnastics}, \eng{go swimming}. Savoir nommer une activité, c'est bien ; mais dans la vraie vie, on ne se contente pas de nommer : on \cle{parle de ses goûts}, on \cle{invite ses amis}, on \cle{accepte} ou on \cle{refuse}, et on raconte ensuite le bon moment passé ensemble. C'est exactement l'objet de cette section : vous donner les phrases toutes faites et les expressions idiomatiques qui font vivre le vocabulaire des loisirs dans la conversation quotidienne.

\courssub{Observer avant d'apprendre : un dialogue entre amis}

\begin{textebox}[At the school gate — a dialogue]
\textbf{Mary:} Hi, John! Are you doing anything special this weekend?\par
\textbf{John:} Not really. I'll probably just stay at home and kill time watching television.\par
\textbf{Mary:} That sounds boring! Do you fancy a game of draughts at my place on Saturday afternoon? My cousins from Bamenda are visiting, and they are crazy about board games.\par
\textbf{John:} Sorry, I'm busy on Saturday. I have football training twice a week, and Saturday is one of the days. You know I'm keen on football!\par
\textbf{Mary:} What about Sunday, then? How about a picnic by the river instead?\par
\textbf{John:} Sunday? That sounds great! I'd love to. What time shall we meet?\par
\textbf{Mary:} Let's meet at ten o'clock, after church. Bring your guitar — I know you enjoy playing music.\par
\textbf{John:} No problem. I'm sure we'll have a good time. Time flies when you're having fun!\par
\textbf{Mary:} True! See you on Sunday, then. Take it easy until then!\par
\textbf{John:} You too. Bye!
\end{textebox}

\begin{definition}[Glossaire du dialogue]
\begin{itemize}
\item \eng{to kill time} (expression) : tuer le temps, occuper le temps.
\item \eng{draughts} (nom, britannique) : le jeu de dames.
\item \eng{board games} (nom pluriel) : les jeux de société.
\item \eng{training} (nom) : l'entraînement.
\item \eng{shall we meet?} : on se retrouve ? / est-ce qu'on se retrouve ? (proposition).
\item \eng{take it easy} (expression) : détends-toi, ne te fatigue pas, prends ton temps.
\end{itemize}
\end{definition}

\textbf{Questions d'observation.} 1) How does Mary invite John? 2) Why does John refuse the first invitation? 3) Which expression shows that John accepts the second invitation? 4) Find two expressions that talk about free time. — Réponses : 1) \eng{Do you fancy a game of draughts?} ; 2) Because he has football training ; 3) \eng{I'd love to} ; 4) \eng{kill time}, \eng{have a good time}.

\courssub{Exprimer ses goûts : likes and dislikes}

\begin{propriete}[Les structures pour exprimer les goûts]
Toutes ces expressions sont suivies d'un \textbf{nom} ou d'un \textbf{verbe en -ing} (jamais de l'infinitif avec \eng{to} après \eng{enjoy}, \eng{fancy}, \eng{be fond of}, \eng{be keen on}) :
\begin{itemize}
\item \eng{I enjoy reading.} — J'aime lire. (\eng{enjoy} + V-ing)
\item \eng{I'm fond of dancing.} — J'aime beaucoup danser. (\eng{be fond of} + nom/V-ing)
\item \eng{I'm keen on football.} — Je suis passionné de foot. (\eng{be keen on} + nom/V-ing)
\item \eng{I'm crazy about music.} — Je suis fou de musique. (familier ; \eng{be crazy about} + nom/V-ing)
\item \eng{I'm a fan of Makossa.} — Je suis fan de makossa. (\eng{be a fan of} + nom)
\item Pour le dégoût : \eng{I don't like...}, \eng{I hate...}, \eng{I can't stand video games.} — Je ne supporte pas les jeux vidéo. (\eng{can't stand} + nom/V-ing)
\end{itemize}
\end{propriete}

\begin{exemplebox}[De l'aversion à la passion : l'échelle des goûts]
\begin{itemize}
\item \eng{I can't stand video games.} — Je ne supporte pas les jeux vidéo.
\item \eng{I don't really like cooking.} — Je n'aime pas trop cuisiner.
\item \eng{I like playing cards.} — J'aime jouer aux cartes.
\item \eng{I enjoy reading novels.} — J'aime bien lire des romans.
\item \eng{I'm fond of dancing.} — J'aime beaucoup danser.
\item \eng{I'm keen on basketball.} — Je suis passionné de basket.
\item \eng{I'm crazy about music.} — Je suis fou de musique.
\end{itemize}
\end{exemplebox}

\begin{attention}[Le piège de « enjoy »]
\eng{Enjoy} n'est \textbf{jamais} suivi de l'infinitif : on dit \eng{I enjoy swimming} et non \emph{I enjoy to swim} ✗. Autre piège : \eng{enjoy} ne prend jamais de complément introduit par \eng{of} : \eng{I enjoyed the party} (j'ai aimé la fête), pas \emph{I enjoyed of the party} ✗. Enfin, pour dire « amuse-toi bien », l'anglais utilise le réfléchi : \eng{Enjoy yourself!} (singulier), \eng{Enjoy yourselves!} (pluriel).
\end{attention}

\courssub{Inviter et proposer une activité}

\begin{propriete}[Les formules d'invitation, du plus formel au plus familier]
\begin{itemize}
\item \eng{Would you like to join us?} — Voudrais-tu te joindre à nous ? (poli, suivi de l'infinitif : \eng{Would you like to come?})
\item \eng{Do you fancy a game of cards?} — Ça te dit une partie de cartes ? (familier, britannique ; \eng{fancy} + nom ou V-ing : \eng{Do you fancy going out?})
\item \eng{How about a picnic?} / \eng{What about a picnic?} — Et si on faisait un pique-nique ? (suivi d'un nom ou d'un V-ing : \eng{How about going to the cinema?})
\item \eng{Let's go to the cinema!} — Allons au cinéma ! (\eng{Let's} + base verbale, pour proposer en s'incluant soi-même)
\item \eng{Why don't we play draughts?} — Pourquoi ne jouerions-nous pas aux dames ?
\end{itemize}
\end{propriete}

\begin{center}
\begin{tabularx}{\linewidth}{|X|X|X|}
\hline
\rowcolor{popblueL} Français & English & Remarque \\
\hline
Ça te dit une partie de dames ? & Do you fancy a game of draughts? & familier, britannique \\
\hline
Tu veux sortir ? & Do you fancy going out? & \eng{fancy} + V-ing \\
\hline
Voudrais-tu venir avec nous ? & Would you like to come with us? & poli, + infinitif \\
\hline
Et si on allait au cinéma ? & How about going to the cinema? & \eng{how about} + V-ing \\
\hline
Faisons un pique-nique ! & Let's have a picnic! & \eng{let's} + base verbale \\
\hline
\end{tabularx}
\end{center}

\begin{savaistu}[« Fancy », un mot très britannique]
Le verbe \eng{fancy} suivi d'un nom ou d'un verbe en \textbf{-ing} est typique de l'anglais britannique familier. \eng{Do you fancy going out?} signifie simplement « Tu veux sortir ? » et \eng{I fancy a drink} veut dire « J'ai envie d'un verre ». À l'écrit ou dans un contexte formel (examen, lettre), préférez \eng{Would you like...?}. Comme adjectif, \eng{fancy} signifie « élégant, sophistiqué » (ex. \eng{a fancy restaurant}, un restaurant chic). Attention : \eng{a fancy dress party} est une « soirée déguisée » (costume party), et non une fête de robes chiques. Ne pas confondre avec le français « fantaisie » !
\end{savaistu}

\courssub{Accepter et refuser une invitation}

\begin{propriete}[Accepter et refuser poliment]
\textbf{Accepter :} \eng{That sounds great!} (Ça a l'air génial !) ; \eng{I'd love to.} (J'aimerais beaucoup / Avec plaisir.) ; \eng{Good idea!} ; \eng{Why not?} ; \eng{Great, let's go!}\par
\textbf{Refuser :} on refuse toujours \textbf{poliment}, en s'excusant et souvent en donnant une raison : \eng{Sorry, I'm busy.} (Désolé, je suis occupé.) ; \eng{I'm afraid I can't.} (Je crains de ne pas pouvoir.) ; \eng{I'd love to, but I have to study.} (J'aimerais bien, mais je dois étudier.) ; \eng{Maybe another time.} (Une autre fois, peut-être.)
\end{propriete}

\begin{attention}[« I am agree » ✗ et « I love » ✗]
Deux erreurs très fréquentes chez les francophones :
\begin{itemize}
\item Pour dire « je suis d'accord », on dit \eng{I agree} ✓ et jamais \emph{I am agree} ✗ : \eng{agree} est un \textbf{verbe}, pas un adjectif ; il n'a donc pas besoin du verbe \eng{be}. De même : \eng{I disagree} (je ne suis pas d'accord).
\item Pour accepter une invitation, on dit \eng{I'd love to} ✓ (= \eng{I would love to}) et non \emph{I love} ✗ seul. Le petit mot \eng{to} est obligatoire : il remplace toute la proposition (\eng{I'd love to come} → \eng{I'd love to}).
\end{itemize}
\end{attention}

\begin{popfigure}
\begin{center}
\begin{tikzpicture}[
node distance=1.1cm,
inv/.style={rectangle, rounded corners, draw=popblue, fill=popblueL, thick, align=center, minimum width=4.6cm, minimum height=1.1cm, font=\small},
acc/.style={rectangle, rounded corners, draw=popgreen, fill=popgreenL, thick, align=center, minimum width=4.2cm, minimum height=1.1cm, font=\small},
ref/.style={rectangle, rounded corners, draw=poporange, fill=poporangeL, thick, align=center, minimum width=4.2cm, minimum height=1.1cm, font=\small},
fleche/.style={-{Stealth[length=3mm]}, thick, popdark}]
\node[inv, align=center] (invite){\textbf{Invitation}\\ “Do you fancy a game of cards?”};
\node[acc, below left=1.2cm and -1.2cm of invite, align=center] (accept){\textbf{Accept}\\ “I'd love to!”\\ “That sounds great!”};
\node[ref, below right=1.2cm and -1.2cm of invite, align=center] (refuse){\textbf{Refuse politely}\\ “Sorry, I'm busy.”\\ “I'm afraid I can't.”};
\draw[fleche] (invite.south) -- (accept.north) node[midway, left, font=\small\itshape, popgreen]{yes};
\draw[fleche] (invite.south) -- (refuse.north) node[midway, right, font=\small\itshape, poporange]{no};
\end{tikzpicture}
\end{center}
\legende{Réagir à une invitation : accepter avec enthousiasme ou refuser poliment.}
\end{popfigure}

\courssub{Les expressions idiomatiques du temps libre}

Une \cle{expression idiomatique} est un groupe de mots dont le sens global ne se devine pas mot à mot. Il faut l'apprendre comme un tout, sans la traduire littéralement.

\begin{definition}[Cinq expressions idiomatiques à connaître par cœur]
\begin{itemize}
\item \eng{to have a good time} : passer un bon moment. \eng{We had a good time at the party.} — Nous avons passé un bon moment à la fête. (Attention au passé irrégulier : \eng{have} → \eng{had} → \eng{had}.)
\item \eng{to kill time} : tuer le temps. \eng{I played cards to kill time at the station.} — J'ai joué aux cartes pour tuer le temps à la gare.
\item \eng{to take it easy} : se détendre, y aller doucement. \eng{After the exam, just take it easy!} — Après l'examen, détends-toi !
\item \eng{to be a fan of} : être fan de. \eng{My father is a fan of Eto'o.} — Mon père est fan d'Eto'o.
\item \eng{Time flies when you're having fun.} — Le temps passe vite quand on s'amuse. (Proverbe : \eng{flies} est le présent de \eng{fly}, voler — le temps « vole ».)
\end{itemize}
\end{definition}

\begin{exemplebox}[Les expressions idiomatiques en contexte]
\begin{itemize}
\item \eng{We had a good time at the party in Bonabéri.} — Nous avons passé un bon moment à la fête à Bonabéri.
\item \eng{During the holidays, I just want to take it easy.} — Pendant les vacances, je veux juste me détendre.
\item \eng{She listened to music to kill time on the bus to Douala.} — Elle a écouté de la musique pour tuer le temps dans le bus pour Douala.
\item \eng{Are you a fan of basketball?} — Tu es fan de basket ?
\item \eng{The match ended quickly — time flies when you're having fun!} — Le match a fini vite : le temps passe vite quand on s'amuse !
\end{itemize}
\end{exemplebox}

\courssub{Parler de la fréquence de ses activités}

\begin{propriete}[Les expressions de fréquence et de moment]
\begin{itemize}
\item \eng{once a week} : une fois par semaine ; \eng{twice a week} : deux fois par semaine ; \eng{three times a month} : trois fois par mois. (On ne dit jamais \emph{two times a week} ✗ : \eng{once} et \eng{twice} sont des mots spéciaux.)
\item \eng{at weekends} (britannique) / \eng{on weekends} (américain) : le week-end.
\item \eng{in my spare time} / \eng{in my free time} : pendant mon temps libre.
\item \eng{after school} : après l'école ; \eng{during the holidays} : pendant les vacances.
\end{itemize}
Ces expressions se placent généralement \textbf{à la fin} de la phrase : \eng{I play football twice a week.} — Je joue au foot deux fois par semaine.
\end{propriete}

\begin{exemplebox}[Fréquence : exemples traduits]
\begin{itemize}
\item \eng{We go swimming once a week, on Saturdays.} — Nous allons nager une fois par semaine, le samedi.
\item \eng{What do you do in your spare time?} — Que fais-tu pendant ton temps libre ?
\item \eng{She visits her grandmother at weekends.} — Elle rend visite à sa grand-mère le week-end.
\item \eng{I read novels during the holidays.} — Je lis des romans pendant les vacances.
\end{itemize}
\end{exemplebox}

\courssub{Méthode : réussir un mini-dialogue d'invitation}

\begin{methode}[Construire un mini-dialogue en cinq étapes]
Pour produire un dialogue d'invitation naturel (à l'oral comme à l'écrit), suivez toujours ce plan :
\begin{enumerate}
\item \textbf{Saluer} et prendre des nouvelles : \eng{Hi, Mary! How are you?}
\item \textbf{Proposer une activité} avec une formule d'invitation : \eng{Do you fancy a game of draughts?} / \eng{How about a picnic?}
\item \textbf{Réagir} : accepter (\eng{I'd love to!}) ou refuser poliment avec une excuse (\eng{Sorry, I'm busy. I have training.}).
\item \textbf{Fixer le jour, l'heure et le lieu} : \eng{Let's meet on Sunday at ten, at my place.}
\item \textbf{Conclure} chaleureusement : \eng{Great! See you on Sunday. Take it easy!}
\end{enumerate}
Astuce : pimentez votre dialogue avec une expression idiomatique (\eng{have a good time}, \eng{kill time}) et une expression de fréquence (\eng{twice a week}, \eng{in my spare time}) : c'est ce qui fait la différence entre un dialogue scolaire et un dialogue authentique.
\end{methode}

\begin{experience}[Jeu de rôle : the weekend plan]
Par groupes de deux. L'élève A invite l'élève B à une activité de son choix (pique-nique à Mvog-Betsi, partie de dames, match de foot, cinéma à Douala). L'élève B refuse la première proposition avec une excuse crédible, puis accepte la deuxième. Les deux élèves fixent ensemble le jour, l'heure et le lieu, puis concluent. Chaque dialogue doit contenir au moins : une formule d'invitation, une expression de goût, une expression idiomatique et une expression de fréquence. Les meilleurs dialogues sont joués devant la classe.
\end{experience}

\begin{exoresolu}[Compléter un dialogue avec les expressions utiles]
Énoncé — Complete the dialogue with: \eng{I'd love to} / \eng{Do you fancy} / \eng{I'm afraid I can't} / \eng{have a good time} / \eng{twice a week}.\par
Amina: Hi, Paul! \_\_\_\_\_\_\_\_ going to the stadium on Saturday? There's a big match.\par
Paul: Hmm, \_\_\_\_\_\_\_\_. I have my piano lesson on Saturdays.\par
Amina: Oh, what a pity! Do you play the piano often?\par
Paul: Yes, I practise \_\_\_\_\_\_\_\_, on Wednesdays and Saturdays. But what about Sunday afternoon?\par
Amina: Sunday? \_\_\_\_\_\_\_\_! Let's meet at two o'clock.\par
Paul: Perfect. I'm sure we'll \_\_\_\_\_\_\_\_ at the match!
\tcblower
Solution détaillée :
\begin{enumerate}
\item \eng{Do you fancy} : c'est la formule d'invitation, suivie du V-ing \eng{going} — exactement la structure \eng{fancy} + V-ing.
\item \eng{I'm afraid I can't} : Paul refuse poliment et donne une raison (son cours de piano). \eng{Sorry, I'm busy} aurait aussi été possible.
\item \eng{twice a week} : la précision « on Wednesdays and Saturdays » confirme la fréquence de deux fois par semaine ; on n'écrit jamais \emph{two times a week}.
\item \eng{I'd love to} : Amina accepte avec enthousiasme ; le \eng{to} final est obligatoire.
\item \eng{have a good time} : expression idiomatique « passer un bon moment » ; après \eng{we'll} (will), le verbe reste à la base verbale \eng{have}.
\end{enumerate}
\end{exoresolu}

\begin{exercice}[À vous de jouer]
1) Traduisez en anglais : a) Je suis passionné de danse. b) Ça te dit une partie de cartes ? c) J'aimerais bien, mais je suis occupé. d) Nous avons passé un bon moment au pique-nique. e) Je joue au foot deux fois par semaine.\par
2) Corrigez les erreurs : a) \emph{I am agree with you.} b) \emph{— Would you like to come? — Yes, I love.} c) \emph{I enjoy to swim.} d) \emph{I go swimming two times a week.}\par
3) Rédigez un mini-dialogue de six répliques : vous invitez un(e) camarade à une activité du week-end ; il/elle refuse d'abord, puis accepte une autre proposition.
\end{exercice}

\begin{corrige}[Exercice « À vous de jouer »]
1) a) \eng{I'm keen on dancing.} b) \eng{Do you fancy a game of cards?} c) \eng{I'd love to, but I'm busy.} d) \eng{We had a good time at the picnic.} e) \eng{I play football twice a week.}\par
2) a) \eng{I agree with you.} (\eng{agree} est un verbe, sans \eng{be}.) b) \eng{Yes, I'd love to.} (le \eng{to} est obligatoire.) c) \eng{I enjoy swimming.} (\eng{enjoy} + V-ing.) d) \eng{I go swimming twice a week.}\par
3) Modèle : \eng{— Hi, Sandra! Do you fancy a picnic at the park on Saturday? — Sorry, I'm afraid I can't. I visit my aunt on Saturdays. — What about Sunday, then? — Sunday? I'd love to! — Great! Let's meet at ten at my place. — Perfect. I'm sure we'll have a good time. See you!}
\end{corrige}

\begin{aretenir}[L'essentiel de la section]
\begin{itemize}
\item Goûts : \eng{enjoy / be fond of / be keen on / be crazy about / be a fan of / can't stand} + nom ou V-ing.
\item Inviter : \eng{Would you like to...?} (poli), \eng{Do you fancy...?} (familier), \eng{How about...?}, \eng{Let's...!}
\item Accepter : \eng{I'd love to!} / \eng{That sounds great!} — Refuser : \eng{Sorry, I'm busy.} / \eng{I'm afraid I can't.}
\item Idiomatismes : \eng{have a good time}, \eng{kill time}, \eng{take it easy}, \eng{be a fan of}, \eng{time flies when you're having fun}.
\item Fréquence : \eng{once/twice a week}, \eng{at weekends}, \eng{in my spare time}, \eng{after school}, \eng{during the holidays}.
\item Pièges : \eng{I agree} (pas \emph{I am agree}), \eng{I'd love to} (pas \emph{I love}), \eng{enjoy} + V-ing, \eng{twice} (pas \emph{two times}).
\end{itemize}
\end{aretenir}

\coursec{Faux amis et pièges des francophones}

Dans la section précédente, nous avons appris à manier les expressions idiomatiques et les phrases utiles pour parler de nos loisirs (\eng{to have a blast}, \eng{to kill time}, \eng{to be into something}...). Mais attention : le vocabulaire des activités récréatives cache un danger redoutable pour les francophones. Beaucoup de mots anglais \emph{ressemblent} à des mots français sans avoir le même sens. Ces pièges s'appellent les \cle{faux amis} (\eng{false friends} ou \eng{false cognates}). Cette section vous apprend à les repérer, à les corriger et à ne plus jamais tomber dans leurs filets — ni à l'écrit, ni à l'oral, ni au Bac.

\courssub{Qu'est-ce qu'un faux ami ? Observation d'abord}

\begin{textebox}[A letter full of traps — read carefully]
Dear Peter,

Last Saturday, I \textbf{assisted to} a football match at the Ahmadou Ahidjo Stadium in Yaoundé. \textbf{Actually}, the Cameroon championship is very exciting this season. My cousin, who is very \textbf{sensible}, was so emotional when his team scored that he cried! During the half-time break, we attended a \textbf{lecture} about the history of the Indomitable Lions. The match was difficult, but our team \textbf{eventually} scored in the last minute. Next week, I will \textbf{pass} my English exam, so I must revise hard.

See you soon,
Alain
\end{textebox}

\textbf{Glossaire :} half-time (n.m.) : mi-temps ; to score (verbe) : marquer (un but) ; to revise (verbe) : réviser.

\textbf{Questions de compréhension :}
\begin{enumerate}
\item Did Alain help someone at the stadium, or did he simply go to the match?
\item Is the championship exciting \emph{now}, or is the speaker correcting an idea?
\item Did the team score \emph{possibly}, or \emph{in the end}?
\item Will Alain succeed in his exam, or simply sit for it?
\end{enumerate}

Toutes les phrases en gras de cette lettre contiennent une erreur ou une ambiguïté typique des francophones. Observons-les une à une avant d'énoncer les règles.

\begin{definition}[Faux ami (false friend)]
Un \cle{faux ami} est un mot anglais qui ressemble par sa forme à un mot français mais dont le sens est différent. Exemple : \eng{actually} ressemble à « actuellement » mais signifie « en fait ». Il existe aussi des \cle{semi-faux amis} : mots qui ont un sens commun \emph{et} un sens différent (ex. \eng{to pass} = « réussir » mais aussi « passer » au sens de « dépasser »).
\end{definition}

\courssub{Les six grands faux amis des loisirs et de la vie scolaire}

\textbf{1. \eng{to assist} $\neq$ « assister (à) »}

\begin{propriete}[Assist = aider]
\eng{To assist somebody} signifie \cle{aider quelqu'un} (soutien, aide active). Pour dire « assister à » un match, un concert, une réunion (être présent, spectateur), on utilise \eng{to attend} (formel) ou \eng{to watch} (regarder un spectacle, un match), ou encore \eng{to go to}.
\end{propriete}

\begin{exemplebox}[Assist, attend, watch : ne pas confondre]
\eng{The nurse assisted the doctor during the operation.} « L'infirmière a aidé le médecin pendant l'opération. »

\eng{I attended the match at the stadium in Bafoussam.} « J'ai assisté au match au stade de Bafoussam. »

\eng{We watched the final on television.} « Nous avons regardé la finale à la télévision. »

\eng{Thousands of fans attended the concert in Douala.} « Des milliers de fans ont assisté au concert à Douala. »
\end{exemplebox}

\begin{attention}[I assisted to the match ✗]
La phrase \eng{I assisted to the match} est doublement fausse : \eng{assist} signifie « aider » et ne se construit jamais avec \eng{to}. Dites : \eng{I attended the match} ou \eng{I watched the match}. Réservez \eng{assist} à l'aide : \eng{She assisted the old man in crossing the road.} « Elle a aidé le vieil homme à traverser la route. »
\end{attention}

\textbf{2. \eng{actually} $\neq$ « actuellement »}

\begin{propriete}[Actually = en fait, en réalité]
\eng{Actually} sert à corriger une idée, à préciser la réalité ou à introduire une information surprenante : « en fait, en réalité, à vrai dire ». Pour « actuellement », dites \eng{currently}, \eng{at present}, \eng{at the moment} ou \eng{now}.
\end{propriete}

\begin{exemplebox}[Actually et currently en contexte]
\eng{Actually, I prefer chess to video games.} « En fait, je préfère les échecs aux jeux vidéo. »

\eng{You think he is lazy? Actually, he trains every morning.} « Tu crois qu'il est paresseux ? En réalité, il s'entraîne tous les matins. »

\eng{I am currently learning to play the guitar.} « J'apprends actuellement à jouer de la guitare. »

\eng{At present, she is the best swimmer in the school.} « Actuellement, c'est la meilleure nageuse de l'école. »
\end{exemplebox}

\textbf{3. \eng{eventually} $\neq$ « éventuellement »}

\begin{propriete}[Eventually = finalement, à la fin]
\eng{Eventually} exprime le résultat final après une attente ou des difficultés : « finalement, en fin de compte ». Le français « éventuellement » (= peut-être, si nécessaire) se traduit par \eng{possibly}, \eng{if necessary}, \eng{if need be} ou \eng{maybe}.
\end{propriete}

\begin{exemplebox}[Eventually et possibly en contexte]
\eng{He eventually won the tournament after three difficult matches.} « Il a finalement gagné le tournoi après trois matchs difficiles. »

\eng{After many attempts, she eventually learnt to ride a bike.} « Après de nombreux essais, elle a finalement appris à faire du vélo. »

\eng{We could possibly go swimming this weekend.} « Nous pourrions éventuellement aller nager ce week-end. »

\eng{If necessary, we can cancel the picnic.} « Éventuellement / si nécessaire, nous pouvons annuler le pique-nique. »
\end{exemplebox}

\textbf{4. \eng{to pass an exam} $\neq$ « passer un examen » (s'y présenter)}

\begin{propriete}[Pass = réussir ; take/sit = se présenter]
\eng{To pass an exam} signifie \cle{réussir} un examen. « Passer un examen » au sens de « s'y présenter, le subir » se dit \eng{to take an exam} ou \eng{to sit (for) an exam} (britannique). L'échec se dit \eng{to fail an exam}.
\end{propriete}

\begin{exemplebox}[Pass, take, fail]
\eng{I will take my English exam next Monday.} « Je passerai mon examen d'anglais lundi prochain. »

\eng{She worked hard and passed all her exams.} « Elle a travaillé dur et a réussi tous ses examens. »

\eng{He failed his driving test twice before passing it.} « Il a échoué deux fois à son permis avant de le réussir. »
\end{exemplebox}

\textbf{5. \eng{a lecture} $\neq$ « une lecture »}

\begin{propriete}[Lecture = cours magistral, conférence]
\eng{A lecture} est un cours magistral ou une conférence. « Une lecture » (action de lire) se dit \eng{reading} ; « un livre/une lecture agréable » se dit \eng{a good read}. Celui qui donne la conférence est \eng{a lecturer} (un chargé de cours), pas un « lecteur » au sens de « liseur » (\eng{a reader}).
\end{propriete}

\begin{exemplebox}[Lecture et reading]
\eng{We attended a lecture on traditional dances at the university.} « Nous avons assisté à une conférence sur les danses traditionnelles à l'université. »

\eng{Reading is my favourite hobby.} « La lecture est mon passe-temps préféré. »

\eng{This novel is a very good read.} « Ce roman est une lecture très agréable. »
\end{exemplebox}

\textbf{6. \eng{sensible} $\neq$ « sensible »}

\begin{propriete}[Sensible = raisonnable ; sensitive = sensible]
\eng{Sensible} qualifie une personne ou une décision \cle{raisonnable, sensée}. Le français « sensible » (émotif, fragile) se traduit par \eng{sensitive}. Attention aussi : \eng{sensible shoes} = « des chaussures pratiques » !
\end{propriete}

\begin{exemplebox}[Sensible et sensitive]
\eng{It is sensible to warm up before doing sport.} « Il est raisonnable de s'échauffer avant de faire du sport. »

\eng{My sister is very sensitive; she cries during sad films.} « Ma sœur est très sensible ; elle pleure devant les films tristes. »

\eng{Be sensible and take a bottle of water for the hike.} « Sois raisonnable et prends une bouteille d'eau pour la randonnée. »
\end{exemplebox}

\courssub{Tableaux récapitulatifs et carte mentale}

\begin{center}
\begin{tabularx}{\linewidth}{|X|X|X|}
\hline
\rowcolor{popblueL} Faux ami (anglais) & Vrai sens en français & Mot français — traduction correcte \\
\hline
to assist & aider & assister à $\rightarrow$ to attend / to watch \\
\hline
actually & en fait, en réalité & actuellement $\rightarrow$ currently, at present \\
\hline
eventually & finalement & éventuellement $\rightarrow$ possibly, if necessary \\
\hline
to pass an exam & réussir un examen & passer (se présenter) $\rightarrow$ to take / to sit an exam \\
\hline
a lecture & une conférence, un cours magistral & une lecture $\rightarrow$ reading \\
\hline
sensible & raisonnable, sensé & sensible $\rightarrow$ sensitive \\
\hline
\end{tabularx}
\end{center}

\begin{popfigure}
\begin{tikzpicture}[
  box/.style={draw=popblue, fill=popblueL, rounded corners=2pt, align=center, font=\small, inner sep=4pt},
  frbox/.style={draw=poporange, fill=poporangeL, rounded corners=2pt, align=center, font=\small, inner sep=4pt},
  lbl/.style={font=\footnotesize\itshape, text=popmuted}
]
\node[box, align=center] (a1) at (0,4.4){\textbf{to assist}\\ = aider};
\node[frbox, align=center] (a2) at (6.5,4.4){« assister à »\\ \textbf{to attend / to watch}};
\node[box, align=center] (b1) at (0,3.2){\textbf{actually}\\ = en fait};
\node[frbox, align=center] (b2) at (6.5,3.2){« actuellement »\\ \textbf{currently / at present}};
\node[box, align=center] (c1) at (0,2.0){\textbf{eventually}\\ = finalement};
\node[frbox, align=center] (c2) at (6.5,2.0){« éventuellement »\\ \textbf{possibly / if necessary}};
\node[box, align=center] (d1) at (0,0.8){\textbf{to pass an exam}\\ = réussir};
\node[frbox, align=center] (d2) at (6.5,0.8){« passer (se présenter) »\\ \textbf{to take / to sit}};
\node[box, align=center] (e1) at (0,-0.4){\textbf{a lecture}\\ = une conférence};
\node[frbox, align=center] (e2) at (6.5,-0.4){« une lecture »\\ \textbf{reading}};
\node[box, align=center] (f1) at (0,-1.6){\textbf{sensible}\\ = raisonnable};
\node[frbox, align=center] (f2) at (6.5,-1.6){« sensible »\\ \textbf{sensitive}};
\draw[<->, thick, popdark] (a1) -- (a2);
\draw[<->, thick, popdark] (b1) -- (b2);
\draw[<->, thick, popdark] (c1) -- (c2);
\draw[<->, thick, popdark] (d1) -- (d2);
\draw[<->, thick, popdark] (e1) -- (e2);
\draw[<->, thick, popdark] (f1) -- (f2);
\node[lbl] at (0,5.2) {Faux ami anglais — vrai sens};
\node[lbl] at (6.5,5.2) {Mot français — traduction correcte};
\end{tikzpicture}
\legende{Tableau des faux amis : à gauche, le mot anglais et son vrai sens ; à droite, le mot français piégé et sa traduction correcte.}
\end{popfigure}

\courssub{Les calques : traduire mot à mot, le piège invisible}

Au-delà des faux amis, il existe les \cle{calques} : on traduit une expression française mot à mot et on obtient un anglais incorrect. En voici trois à bannir dans le thème des loisirs.

\begin{propriete}[Les calques à corriger]
\begin{itemize}
\item « faire du sport » : on dit \eng{to do sport} ou \eng{to play sport}, jamais \eng{to make sport} ✗. (\eng{I do sport twice a week.} « Je fais du sport deux fois par semaine. »)
\item « j'ai 17 ans » : l'anglais utilise le verbe \emph{être} : \eng{I am 17 (years old)}, jamais \eng{I have 17 years} ✗.
\item « se promener » : \eng{to take a walk} ✓ est correct, mais \eng{to make a walk} ✗ est un calque. On peut aussi dire \eng{to go for a walk}.
\end{itemize}
\end{propriete}

\begin{methode}[Comment éviter les calques et les faux amis]
\begin{enumerate}
\item \textbf{Méfiez-vous} de tout mot anglais qui ressemble à un mot français : la ressemblance est un signal d'alarme, pas une aide.
\item \textbf{Vérifiez dans le dictionnaire} (bilingue puis monolingue) le sens exact et un exemple d'emploi avant d'utiliser le mot.
\item \textbf{Apprenez les mots en contexte}, dans des collocations complètes : mémorisez \eng{to attend a match}, \eng{to take an exam}, \eng{to do sport} comme des blocs, jamais le mot isolé.
\item \textbf{Traduisez à l'envers} : demandez-vous « comment un Anglais dirait-il cela ? » plutôt que de traduire mot à mot depuis le français.
\item \textbf{Relisez vos productions} en traquant spécifiquement les six faux amis de cette leçon.
\end{enumerate}
\end{methode}

\begin{exoresolu}[Corrige les erreurs]
Énoncé — Chaque phrase contient un faux ami ou un calque. Corrige-la et traduis la phrase corrigée.

1. \eng{I assisted to a basketball match last Friday.}

2. \eng{Actually, my brother is training for the marathon in Bamenda.} (sens voulu : « actuellement »)

3. \eng{We will eventually play football if it rains.} (sens voulu : « éventuellement »)

4. \eng{She passed her exam yesterday, so she is very stressed.} (sens voulu : elle s'est présentée, résultat inconnu)

5. \eng{He makes sport every weekend.}

6. \eng{My friend has 16 years and he is very sensible.} (sens voulu : « sensible » = émotif)

\tcblower
Solution détaillée :

1. \eng{assist} = aider ; « assister à » un match = \eng{attend/watch}. Correction : \eng{I attended (watched) a basketball match last Friday.} « J'ai assisté à (regardé) un match de basket vendredi dernier. »

2. « Actuellement » ne se dit jamais \eng{actually}. Correction : \eng{Currently (At present), my brother is training for the marathon in Bamenda.} « Actuellement, mon frère s'entraîne pour le marathon de Bamenda. »

3. « Éventuellement » = \eng{possibly} ; \eng{eventually} signifie « finalement ». Correction : \eng{We will possibly play football if it rains.} « Nous jouerons éventuellement au football s'il pleut. »

4. Si le résultat est inconnu, elle s'est seulement \emph{présentée} : \eng{She took (sat) her exam yesterday, so she is very stressed.} « Elle a passé son examen hier, donc elle est très stressée. » \eng{Passed} voudrait dire qu'elle l'a réussi — elle serait plutôt soulagée !

5. Calque de « faire du sport » : \eng{He does (plays) sport every weekend.} « Il fait du sport chaque week-end. »

6. Double erreur : l'âge se dit avec \emph{être}, et « sensible » (émotif) = \eng{sensitive}. Correction : \eng{My friend is 16 (years old) and he is very sensitive.} « Mon ami a 16 ans et il est très sensible. »
\end{exoresolu}

\begin{attention}[Le piège du contexte au Bac]
Dans un texte de compréhension écrite, ne traduisez jamais un mot par son « jumeau » français sans vérifier le contexte. Si le texte dit \eng{The team eventually qualified}, demandez-vous : l'équipe s'est-elle qualifiée \emph{finalement} (après des efforts) ou \emph{peut-être} ? Le contexte tranche toujours.
\end{attention}

\begin{savaistu}[Les faux amis, stars des épreuves du Bac]
Les faux amis sont très fréquents dans les épreuves de compréhension du Bac : le mot \eng{actually} dans un texte ne signifie \emph{jamais} « actuellement ». Les concepteurs d'épreuves adorent placer \eng{actually}, \eng{eventually} ou \eng{to assist} dans les textes pour tester votre vigilance. Les faux amis existent parce que le français et l'anglais partagent des milliers de mots d'origine latine — mais au fil des siècles, certains ont dérivé vers des sens différents. \eng{To assist} vient du latin \emph{assistere} (« se tenir auprès de ») : le français a gardé l'idée de « présence », l'anglais celle d'« aide ». Bonne nouvelle : la plupart des mots semblables (\eng{activity}, \eng{family}, \eng{sport}) sont de vrais amis !
\end{savaistu}

\begin{exercice}[À toi de jouer]
1. Traduis en anglais : « En fait, je n'aime pas la natation ; je préfère la randonnée. »

2. Traduis en anglais : « Elle a finalement réussi son examen après l'avoir passé deux fois. »

3. Choisis le bon mot : \eng{We (attended / assisted) a lecture about Cameroonian music.}

4. Corrige : \eng{I have 17 years and I make sport every day.}

5. Vrai ou faux ? \eng{My sensible friend always thinks before acting.} — Ici, \eng{sensible} signifie « raisonnable ».
\end{exercice}

\begin{corrige}[Exercice « À toi de jouer »]
1. \eng{Actually, I don't like swimming; I prefer hiking.} (« En fait » = \eng{actually}, bien placé en tête de phrase.)

2. \eng{She eventually passed her exam after taking it twice.} (\eng{eventually} = finalement ; \eng{passed} = réussi ; \eng{taking} = passée au sens de « présentée ».)

3. \eng{attended} : on « assiste à » une conférence ; \eng{assisted} signifierait que nous avons aidé le conférencier.

4. \eng{I am 17 (years old) and I do (play) sport every day.} (âge avec \emph{être} ; « faire du sport » = \eng{do/play sport}.)

5. Vrai. \eng{Sensible} = raisonnable, sensé : l'ami réfléchit avant d'agir, c'est cohérent. Pour « sensible » (émotif), il faudrait \eng{sensitive}.
\end{corrige}

\begin{aretenir}[L'essentiel de la section]
\begin{itemize}
\item \eng{To assist} = aider ; « assister à » = \eng{to attend / to watch}.
\item \eng{Actually} = en fait ; « actuellement » = \eng{currently, at present}.
\item \eng{Eventually} = finalement ; « éventuellement » = \eng{possibly, if necessary}.
\item \eng{To pass an exam} = réussir ; « passer (se présenter) » = \eng{to take / to sit an exam}.
\item \eng{A lecture} = une conférence ; « une lecture » = \eng{reading}.
\item \eng{Sensible} = raisonnable ; « sensible » = \eng{sensitive}.
\item Calques à bannir : \eng{make sport} ✗ $\rightarrow$ \eng{do/play sport} ; \eng{I have 17 years} ✗ $\rightarrow$ \eng{I am 17} ; \eng{make a walk} ✗ $\rightarrow$ \eng{take a walk / go for a walk}.
\item Réflexe d'or : tout mot anglais qui ressemble à un mot français doit être vérifié dans le dictionnaire avant emploi.
\end{itemize}
\end{aretenir}

\coursec{Familles de mots, prononciation et orthographe}

Après avoir démasqué les faux amis qui guettent les francophones, nous allons maintenant apprendre à \emph{construire} le vocabulaire des loisirs à partir d'un seul mot. En anglais comme en français, les mots vivent en famille : un verbe peut donner naissance à un nom, à un adjectif, parfois à un nom de personne. Connaître ces familles, c'est multiplier son vocabulaire par trois ou quatre sans effort supplémentaire. Nous verrons ensuite comment prononcer correctement quelques mots pièges du thème, et comment les orthographier sans faute.

\courssub{Observer avant de mémoriser : un texte pour découvrir les familles de mots}

\begin{textebox}[A weekend of leisure in Bamenda]
Last weekend was pure relaxation for my family. On Saturday morning, my father went jogging while my mother practised her guitar — she finds playing very relaxing. In the afternoon, we watched a football competition at the municipal stadium. The local team played well, and their best athlete, a young runner, won the race easily. She is a very competitive person, but also a fair competitor: after her victory, she congratulated all the other runners.

In the evening, we looked for some entertainment in town. A famous comedian gave a show at the community hall. His jokes were so amusing that everyone laughed for an hour. My little brother, who is usually bored, was really amused. For me, the most enjoyable moment of the weekend was Sunday afternoon: we had a picnic by the river, and I finally finished my novel. Reading is my favourite hobby, and it gives me great enjoyment. My sister prefers more active leisure activities like swimming and dancing. “There is something for everyone in this family,” my father said with a smile. And he was right!
\end{textebox}

\textbf{Glossaire :}

\begin{center}
\begin{tabularx}{\linewidth}{|l|l|X|}
\hline
\rowcolor{popblueL} Mot anglais & Nature & Traduction française \\
\hline
relaxation & nom & détente, relaxation \\
\hline
to practise & verbe & pratiquer, s'entraîner à \\
\hline
competition & nom & compétition \\
\hline
competitor & nom & compétiteur, concurrent \\
\hline
entertainment & nom & divertissement, animation \\
\hline
amusing & adjectif & amusant, drôle \\
\hline
enjoyment & nom & plaisir \\
\hline
leisure & nom & loisir(s) \\
\hline
\end{tabularx}
\end{center}

\textbf{Questions de compréhension :}
\begin{enumerate}
\item What relaxing activity does the mother practise?
\item Who won the race at the stadium, and what quality does she show?
\item Why was the little brother amused in the evening?
\item What is the narrator's favourite hobby?
\end{enumerate}

\begin{corrige}[Questions de compréhension]
\begin{enumerate}
\item She practises the guitar, which she finds very relaxing.
\item The best athlete of the local team, a young runner, won the race. She is competitive but fair: she congratulated all the other runners.
\item He was amused because the comedian's jokes were very amusing.
\item The narrator's favourite hobby is reading, which gives him great enjoyment.
\end{enumerate}
\end{corrige}

\courssub{Les cinq grandes familles de mots des loisirs}

Dans le texte ci-dessus, observez : \eng{relaxation}, \eng{relaxing}, \eng{competition}, \eng{competitive}, \eng{competitor}, \eng{entertainment}, \eng{amusing}, \eng{amused}, \eng{enjoyable}, \eng{enjoyment}. Tous ces mots dérivent de cinq verbes de base. Voici les familles complètes.

\begin{definition}[Famille 1 : relax]
\begin{itemize}
\item \cle{relax} (verbe) : se détendre. \eng{I relax by listening to music.} « Je me détends en écoutant de la musique. »
\item \cle{relaxation} (nom) : la détente, la relaxation. \eng{Yoga is good for relaxation.} « Le yoga est bon pour la relaxation. »
\item \cle{relaxing} (adjectif) : relaxant, reposant (qualité de l'activité). \eng{a relaxing evening} « une soirée relaxante ».
\item \cle{relaxed} (adjectif) : détendu (sentiment de la personne). \eng{I feel relaxed after the holidays.} « Je me sens détendu après les vacances. »
\end{itemize}
\end{definition}

\begin{definition}[Famille 2 : amuse]
\begin{itemize}
\item \cle{amuse} (verbe) : amuser, divertir. \eng{The clown amuses the children.} « Le clown amuse les enfants. »
\item \cle{amusement} (nom) : l'amusement. \eng{an amusement park} « un parc d'attractions ».
\item \cle{amusing} (adjectif) : amusant, drôle. \eng{an amusing story} « une histoire amusante ».
\item \cle{amused} (adjectif) : amusé (sentiment). \eng{We were amused by his jokes.} « Nous étions amusés par ses blagues. »
\end{itemize}
\end{definition}

\begin{definition}[Famille 3 : enjoy]
\begin{itemize}
\item \cle{enjoy} (verbe) : apprécier, prendre plaisir à. Attention : \eng{enjoy} est suivi d'un nom ou d'un verbe en \emph{-ing}. \eng{I enjoy swimming.} « J'aime nager. »
\item \cle{enjoyment} (nom) : le plaisir. \eng{She sings for enjoyment, not for money.} « Elle chante pour le plaisir, pas pour l'argent. »
\item \cle{enjoyable} (adjectif) : agréable, plaisant. \eng{It was a very enjoyable trip.} « C'était un voyage très agréable. »
\end{itemize}
\end{definition}

\begin{definition}[Famille 4 : entertain]
\begin{itemize}
\item \cle{entertain} (verbe) : divertir, recevoir (des invités). \eng{The singer entertained the crowd.} « Le chanteur a diverti la foule. »
\item \cle{entertainment} (nom) : le divertissement, l'animation. \eng{There is little entertainment in the village.} « Il y a peu de divertissements au village. »
\item \cle{entertaining} (adjectif) : divertissant. \eng{an entertaining film} « un film divertissant ».
\end{itemize}
\end{definition}

\begin{definition}[Famille 5 : compete]
\begin{itemize}
\item \cle{compete} (verbe) : concourir, rivaliser. \eng{Eight teams compete in the tournament.} « Huit équipes participent au tournoi. »
\item \cle{competition} (nom) : la compétition. \eng{a singing competition} « un concours de chant ».
\item \cle{competitor} (nom de personne) : le/la compétiteur(-trice), le concurrent. \eng{Each competitor receives a number.} « Chaque concurrent reçoit un numéro. »
\item \cle{competitive} (adjectif) : compétitif, qui a l'esprit de compétition. \eng{My brother is very competitive.} « Mon frère a un grand esprit de compétition. »
\end{itemize}
\end{definition}

\begin{popfigure}
\begin{center}
\begin{tikzpicture}[
  grow=right,
  level distance=3.6cm,
  level 1/.style={sibling distance=2.6cm},
  every node/.style={draw, rounded corners, align=center, font=\small},
  edge from parent/.style={draw=popmuted, -{Stealth}}
]
\node[fill=popblue, text=white, font=\bfseries, align=center]{compete\\(verbe)}
  child { node[fill=poppurpleL, align=center]{competitive\\(adjectif : compétitif)} }
  child { node[fill=popgreenL, align=center]{competitor\\(nom : le concurrent)} }
  child { node[fill=poporangeL, align=center]{competition\\(nom : la compétition)} };
\end{tikzpicture}
\end{center}
\legende{La famille du verbe \eng{compete} : un verbe, deux noms (l'activité et la personne), un adjectif.}
\end{popfigure}

\begin{propriete}[Adjectifs en -ing et adjectifs en -ed : la règle d'or]
De nombreux verbes de sentiment forment deux adjectifs :
\begin{itemize}
\item L'adjectif en \cle{-ing} décrit la \textbf{chose ou l'activité} qui provoque le sentiment : \eng{a relaxing evening} « une soirée relaxante » (c'est la soirée qui détend), \eng{an amusing story} « une histoire amusante », \eng{an exciting match} « un match passionnant ».
\item L'adjectif en \cle{-ed} décrit le \textbf{sentiment ressenti par la personne} : \eng{I feel relaxed} « je me sens détendu », \eng{we were amused} « nous étions amusés », \eng{the fans were excited} « les supporters étaient enthousiastes ».
\end{itemize}
Le français rend souvent les deux par des mots proches, et la paire « intéressant » / « intéressé » prête à confusion chez les élèves. En anglais, la distinction est obligatoire et sans exception.
\end{propriete}

\begin{exemplebox}[La distinction -ing / -ed en action]
\begin{itemize}
\item \eng{The match was exciting.} « Le match était passionnant. » (le match provoque l'excitation)
\item \eng{The fans were excited.} « Les supporters étaient enthousiastes. » (les supporters ressentent l'excitation)
\item \eng{The film is boring.} « Le film est ennuyeux. » / \eng{The students are bored.} « Les élèves s'ennuient. »
\item \eng{The lesson was interesting.} « La leçon était intéressante. » / \eng{The pupils were interested.} « Les élèves étaient intéressés. »
\item \eng{This game is tiring.} « Ce jeu est fatigant. » / \eng{After the game, I am tired.} « Après le match, je suis fatigué. »
\end{itemize}
\end{exemplebox}

\begin{attention}[Ne dites pas « I am boring » pour « je m'ennuie » !]
Erreur classique du francophone : \eng{I am boring} signifie « je suis ennuyeux » (c'est-à-dire que les autres s'ennuient avec moi !). Pour dire « je m'ennuie », il faut l'adjectif en \emph{-ed} : \eng{I am bored}. De même : \eng{I am interested in sports} « je m'intéresse au sport », et non \emph{I am interesting in sports} (« je suis intéressant »).
\end{attention}

\begin{center}
\begin{tabularx}{\linewidth}{|X|X|X|}
\hline
\rowcolor{popblueL} Verbe & Adjectif en -ing (l'activité) & Adjectif en -ed (le sentiment) \\
\hline
relax & relaxing (relaxant) & relaxed (détendu) \\
\hline
amuse & amusing (amusant) & amused (amusé) \\
\hline
excite & exciting (passionnant) & excited (enthousiaste) \\
\hline
bore & boring (ennuyeux) & bored (ennuyé) \\
\hline
interest & interesting (intéressant) & interested (intéressé) \\
\hline
tire & tiring (fatigant) & tired (fatigué) \\
\hline
\end{tabularx}
\end{center}

\courssub{Prononciation : les mots pièges du thème}

L'orthographe anglaise ne dit pas toujours comment prononcer. Voici les mots du vocabulaire des loisirs que les francophones prononcent le plus souvent de travers.

\begin{center}
\begin{tabularx}{\linewidth}{|l|X|X|}
\hline
\rowcolor{popblueL} Mot & Prononciation (lettres françaises) & Piège à éviter \\
\hline
leisure & « lè-jeur » & ne pas dire « lé-sure » ni « lé-zeur » \\
\hline
choir & « kouaïeur » & le \emph{ch} se prononce « k », pas « ch » français \\
\hline
guitar & « gui-tar » & accent sur la 2\textsuperscript{e} syllabe ; le \emph{u} ne se prononce pas \\
\hline
tournament & « tour-ne-ment » & ne pas dire « tour-na-ment » à la française \\
\hline
athlete & « az-lite » (2 syllabes) & ne jamais dire « az-lète » en 3 syllabes \\
\hline
\end{tabularx}
\end{center}

\begin{attention}[Athlete : deux syllabes seulement !]
Sous l'influence du français « athlète » (3 syllabes), beaucoup d'élèves disent « az-lète » ou « az-a-lite ». Faux ! \eng{Athlete} se prononce « az-lite », en \textbf{deux syllabes}, avec un \emph{th} doux comme dans \eng{think}. Entraînez-vous : \eng{The athlete won the tournament.} « L'athlète a gagné le tournoi » — « di az-lite ouane de tour-ne-ment ».
\end{attention}

\courssub{Orthographe : pluriels, doublements et variantes britanniques}

\begin{propriete}[Le pluriel de hobby]
Les noms terminés par \textbf{consonne + y} forment leur pluriel en \emph{-ies} : \cle{hobby} $\rightarrow$ \cle{hobbies} « passe-temps ». \eng{My hobbies are reading and dancing.} « Mes passe-temps sont la lecture et la danse. » De même : \eng{activity} $\rightarrow$ \eng{activities}, \eng{story} $\rightarrow$ \eng{stories}. Mais après une voyelle, le \emph{y} reste : \eng{boy} $\rightarrow$ \eng{boys}, \eng{day} $\rightarrow$ \eng{days}.
\end{propriete}

\begin{propriete}[Le doublement de la consonne finale]
Un verbe ou un adjectif d'\textbf{une syllabe} terminé par \textbf{une voyelle + une consonne} double cette consonne devant \emph{-ing}, \emph{-er}, \emph{-ed} :
\begin{itemize}
\item \cle{swim} $\rightarrow$ \cle{swimming} (deux \emph{m}) : \eng{Swimming is my favourite sport.} « La natation est mon sport préféré. »
\item \cle{jog} $\rightarrow$ \cle{jogging} (deux \emph{g}) : \eng{He goes jogging every morning.} « Il fait du jogging chaque matin. »
\item \cle{win} $\rightarrow$ \cle{winner} (deux \emph{n}) : \eng{The winner receives a trophy.} « Le gagnant reçoit un trophée. »
\item \cle{run} $\rightarrow$ \cle{running}, \cle{runner} ; \cle{begin} $\rightarrow$ \cle{beginning} (le doublement s'explique ici par l'accent tonique sur la dernière syllabe).
\end{itemize}
En revanche, pas de doublement si la consonne finale suit deux voyelles : \eng{read} $\rightarrow$ \eng{reading}, \eng{play} $\rightarrow$ \eng{playing}.
\end{propriete}

\begin{savaistu}[Practise ou practice ?]
En anglais \textbf{britannique}, le verbe et le nom s'écrivent différemment : le verbe s'écrit \eng{practise} (\eng{I practise the guitar every day.} « Je m'entraîne à la guitare tous les jours. ») et le nom s'écrit \eng{practice} (\eng{Practice makes perfect.} « C'est en forgeant qu'on devient forgeron. »). En anglais \textbf{américain}, les deux s'écrivent \eng{practice}. Le même phénomène existe pour \eng{licence} (nom) et \eng{license} (verbe). Le français connaît une distinction voisine avec « conseiller » (verbe) et « conseil » (nom) : la langue anglaise britannique fonctionne de la même manière avec \emph{-se} pour le verbe et \emph{-ce} pour le nom.
\end{savaistu}

\begin{attention}[Les trois fautes d'orthographe les plus fréquentes]
\begin{enumerate}
\item \eng{swimming} prend \textbf{deux m} (jamais \emph{swiming}).
\item \eng{winner} prend \textbf{deux n} (jamais \emph{winer}).
\item \eng{hobbies} s'écrit avec \textbf{-ies} au pluriel (jamais \emph{hobbys}).
\end{enumerate}
Ces trois mots reviennent constamment dans les rédactions sur les loisirs : apprenez-les une fois pour toutes !
\end{attention}

\begin{exoresolu}[Compléter avec la bonne forme du mot entre parenthèses]
Énoncé — Complete the sentences with the correct form of the word in brackets.
\begin{enumerate}
\item The \dots\ of the race received a gold medal. (win)
\item Watching comedies is very \dots\ . (relax)
\item We were all \dots\ by the magician's tricks. (amuse)
\item My \dots\ are swimming and playing the guitar. (hobby)
\item She is a great \dots\ ; she has won many races. (compete)
\end{enumerate}
\tcblower
Solution détaillée :
\begin{enumerate}
\item \textbf{winner} : il faut un nom de personne ; \eng{win} double le \emph{n} devant \emph{-er}. « Le gagnant de la course a reçu une médaille d'or. »
\item \textbf{relaxing} : adjectif en \emph{-ing}, car on décrit l'activité (regarder des comédies). « Regarder des comédies est très relaxant. »
\item \textbf{amused} : adjectif en \emph{-ed}, car on décrit le sentiment des personnes. « Nous avons tous été amusés par les tours du magicien. »
\item \textbf{hobbies} : pluriel en \emph{-ies} (consonne + y), et le verbe \eng{are} confirme le pluriel. « Mes passe-temps sont la natation et la guitare. »
\item \textbf{competitor} : il faut le nom de personne de la famille de \eng{compete}. « C'est une grande compétitrice ; elle a gagné de nombreuses courses. »
\end{enumerate}
\end{exoresolu}

\begin{exercice}[À vous de jouer]
\begin{enumerate}
\item Mettez le mot entre parenthèses à la bonne forme : \eng{The film was so \dots\ (bore) that half the audience left.} / \eng{The children were \dots\ (excite) before the show.} / \eng{Reading gives me great \dots\ (enjoy).}
\item Écrivez le pluriel : \eng{hobby}, \eng{activity}, \eng{story}, \eng{boy}.
\item Écrivez la forme en \emph{-ing} : \eng{swim}, \eng{jog}, \eng{run}, \eng{read}, \eng{play}.
\item Traduisez : « Je m'ennuie pendant les longs matchs ennuyeux. » (attention à \emph{-ed} et \emph{-ing} !)
\end{enumerate}
\end{exercice}

\begin{corrige}[Exercice « À vous de jouer »]
\begin{enumerate}
\item \eng{boring} (le film provoque l'ennui) ; \eng{excited} (les enfants ressentent l'excitation) ; \eng{enjoyment} (nom après l'adjectif \eng{great}).
\item \eng{hobbies}, \eng{activities}, \eng{stories}, \eng{boys} (voyelle + y : le \emph{y} reste).
\item \eng{swimming}, \eng{jogging}, \eng{running} (doublement) ; \eng{reading}, \eng{playing} (pas de doublement : deux voyelles avant la consonne).
\item \eng{I am bored during long boring matches.} — \eng{bored} pour mon sentiment, \eng{boring} pour les matchs.
\end{enumerate}
\end{corrige}

\begin{aretenir}[L'essentiel de la section]
\begin{itemize}
\item Cinq familles à connaître par cœur : \eng{relax / relaxation / relaxing / relaxed}, \eng{amuse / amusement / amusing / amused}, \eng{enjoy / enjoyment / enjoyable}, \eng{entertain / entertainment / entertaining}, \eng{compete / competition / competitor / competitive}.
\item Adjectif en \emph{-ing} = l'activité qui provoque le sentiment ; adjectif en \emph{-ed} = le sentiment de la personne.
\item Prononciation : \eng{leisure} « lè-jeur », \eng{choir} « kouaïeur », \eng{guitar} « gui-tar », \eng{tournament} « tour-ne-ment », \eng{athlete} « az-lite » (2 syllabes).
\item Orthographe : \eng{hobbies}, \eng{swimming}, \eng{jogging}, \eng{winner} ; en britannique, verbe \eng{practise} et nom \eng{practice}.
\end{itemize}
\end{aretenir}

Ces familles de mots et ces règles de forme sont maintenant des outils solides. Dans la prochaine section, nous les remettrons en vie dans un texte, un dialogue et des productions personnelles : c'est en employant le lexique en contexte qu'on le rend vraiment sien.

\coursec{Le lexique en contexte : texte, dialogue et production}

Dans la section précédente, nous avons étudié les familles de mots (\emph{relax / relaxation / relaxing}), la prononciation et l'orthographe du vocabulaire des loisirs. Connaître des mots isolément ne suffit pas : un bon élève de Première doit savoir les \cle{réemployer en contexte}, c'est-à-dire les reconnaître dans un texte (\emph{reading}), dans un dialogue (\emph{listening / speaking}) et les utiliser dans sa propre production écrite (\emph{writing}). C'est l'objectif de cette dernière section : lire, comprendre, parler et écrire avec le lexique des \emph{recreational activities}.

\courssub{Lecture : \emph{Weekend Fun in Yaoundé}}

\begin{methode}[Comment lire un texte sur les loisirs]
\begin{enumerate}
\item \textbf{Lire les questions d'abord} : elles vous indiquent quelles informations chercher.
\item \textbf{Repérer les noms d'activités} (\emph{football match, choir, draughts, cinema}) et les \cle{marqueurs de temps} (\emph{on Saturday mornings, in the afternoon, in the evening}) : ils structurent le texte.
\item \textbf{Souligner les mots inconnus} et deviner leur sens grâce au contexte avant de consulter le glossaire.
\item \textbf{Répondre avec des phrases courtes}, au présent simple, sans recopier tout le texte.
\end{enumerate}
\end{methode}

\begin{textebox}[Weekend Fun in Yaoundé]
On Saturday mornings, the Ngo family goes to the pitch to watch a football match. Mr Ngo is a great fan of the local team: he never misses a game and he always wears the team's colours. His wife, Mrs Ngo, prefers to sit under a tree and chat with the other mothers while the players train.

In the afternoon, his daughter Amina sings in the church choir. She has a beautiful voice and she practises twice a week. Meanwhile, his son Junior plays draughts with his friends under the mango tree near their house. Junior is the champion of the neighbourhood: nobody can beat him easily.

In the evening, the whole family sometimes goes to the cinema in town to watch a Nigerian film, or simply relaxes at home, listening to music and telling stories. On Sundays, after the church service, the children often hang out with their friends in the playground.

Mr Ngo believes that recreational activities are very important. “Work without play makes life boring,” he often says. His children agree: thanks to sport, music and games, they stay healthy, make new friends and forget the stress of school. They always have a good time together, and weekends are their favourite days of the week.
\end{textebox}

\begin{definition}[Glossaire du texte]
\begin{itemize}
\item \eng{pitch} (n.) : terrain (de sport) — « le terrain »
\item \eng{choir} (n.) : chorale — se prononce « kouaïeur » ; le \emph{ch} se prononce « k » !
\item \eng{draughts} (n.) : jeu de dames — se prononce « draafts » ; aux États-Unis on dit \eng{checkers}.
\item \eng{to hang out (with friends)} (v.) : traîner, passer du temps avec des amis
\item \eng{fan} (n.) : supporter, admirateur
\item \eng{to relax} (v.) : se détendre ; \eng{to have a good time} : passer un bon moment
\end{itemize}
\end{definition}

\begin{exemplebox}[Deux phrases clés du texte]
\eng{Amina sings in the church choir.} « Amina chante dans la chorale de l'église. »

\eng{They always have a good time together.} « Ils passent toujours un bon moment ensemble. »

Observez : \eng{to sing} prend un \textbf{-s} à la troisième personne du présent simple (\emph{she sings}), et l'adverbe de fréquence \eng{always} se place \textbf{avant} le verbe principal : \eng{They always have}, jamais \emph{They have always a good time} dans ce type de phrase simple.
\end{exemplebox}

\courssub{Compréhension du texte (reading)}

\begin{exercice}[Comprehension check]
Répondez en anglais, par des phrases courtes.

\textbf{A. True or False?} (justifiez par un mot ou une courte phrase du texte)
\begin{enumerate}
\item Mr Ngo never watches football matches.
\item Junior plays draughts under the mango tree.
\end{enumerate}

\textbf{B. Choose the correct answer.}
\begin{enumerate}
\item In the evening, the Ngo family sometimes \ldots
\begin{enumerate}
\item goes to the stadium \quad b) goes to the cinema \quad c) goes to the church
\end{enumerate}
\end{enumerate}

\textbf{C. Answer with a short sentence.}
\begin{enumerate}
\item Where does Amina sing?
\item Why does Mr Ngo think recreational activities are important?
\end{enumerate}
\end{exercice}

\begin{corrige}[Comprehension check]
\textbf{A.} 1. \textbf{False.} He never misses a game. \quad 2. \textbf{True.} He plays draughts under the mango tree.

\textbf{B.} 1. \textbf{b)} goes to the cinema.

\textbf{C.} 1. \eng{She sings in the church choir.} \quad 2. \eng{He thinks they are important because they make life interesting / because work without play makes life boring.}
\end{corrige}

\begin{attention}[Ne recopiez pas tout le texte !]
Dans les réponses de compréhension, beaucoup d'élèves francophones recopient trois lignes du texte « pour être sûrs ». C'est une erreur de méthode : l'examinateur veut une \cle{phrase courte et correcte}, au \textbf{présent simple}, qui reprend les mots de la question.

Question : \eng{Where does Amina sing?}

Mauvaise réponse : \emph{In the afternoon, his daughter Amina sings in the church choir. She has a beautiful voice and she practises twice a week.} (trop long, recopié)

Bonne réponse : \eng{She sings in the church choir.}

Attention aussi à l'accord : \eng{She sing} est faux ; il faut \eng{She \textbf{sings}}.
\end{attention}

\begin{exoresolu}[Short-answer question]
Énoncé : répondez à la question suivante par une phrase complète et courte : \eng{What does Junior do in the afternoon?}
\tcblower
\textbf{Étape 1} : repérer le mot-clé de la question (\emph{Junior, afternoon}) et le retrouver dans le texte : « his son Junior plays draughts with his friends under the mango tree ».

\textbf{Étape 2} : transformer en réponse personnelle au présent simple, avec le verbe conjugué correctement.

\textbf{Réponse} : \eng{He plays draughts with his friends.} « Il joue aux dames avec ses amis. »

Remarque : on remplace \emph{Junior} par le pronom \eng{he}, et le verbe prend le \textbf{-s} de la troisième personne : \emph{plays}.
\end{exoresolu}

\courssub{Dialogue modèle : planifier son samedi (listening / speaking)}

Le lexique des loisirs sert d'abord à \cle{communiquer}. Voici un dialogue typique entre deux élèves de Première qui organisent leur samedi. Observez les formules d'\textbf{invitation}, d'\textbf{acceptation} et de \textbf{choix de l'activité, de l'heure et du lieu}.

\begin{textebox}[Planning Saturday — a dialogue]
\textbf{Sandrine:} Hi, Boris! Do you have any plans for Saturday?

\textbf{Boris:} Not really. Why do you ask?

\textbf{Sandrine:} Would you like to go swimming with me at the municipal pool?

\textbf{Boris:} Hmm, I can't stand cold water in the morning. What about playing table tennis instead?

\textbf{Sandrine:} Good idea! I am fond of table tennis. Where shall we play?

\textbf{Boris:} At the youth centre near the post office. They have two tables.

\textbf{Sandrine:} Perfect. What time shall we meet?

\textbf{Boris:} Let's meet at ten o'clock, in front of the centre.

\textbf{Sandrine:} All right. And after the game, we can hang out at my place and listen to music.

\textbf{Boris:} Great! I'm looking forward to it. See you on Saturday!

\textbf{Sandrine:} See you! Don't forget your racket!
\end{textebox}

\begin{propriete}[Les formules du dialogue à mémoriser]
\begin{center}
\begin{tabularx}{\linewidth}{|X|X|X|}\hline
\rowcolor{popblueL} Fonction & English & Français \\ \hline
Inviter & \eng{Would you like to go swimming?} & Tu voudrais aller nager ? \\ \hline
Proposer & \eng{What about playing table tennis?} & Et si on jouait au ping-pong ? \\ \hline
Accepter & \eng{Good idea! / All right. / Great!} & Bonne idée ! / D'accord. \\ \hline
Refuser poliment & \eng{I can't stand cold water.} & Je ne supporte pas l'eau froide. \\ \hline
Fixer le lieu & \eng{Where shall we meet?} & Où est-ce qu'on se retrouve ? \\ \hline
Fixer l'heure & \eng{Let's meet at ten o'clock.} & Retrouvons-nous à dix heures. \\ \hline
Conclure & \eng{I'm looking forward to it.} & J'ai hâte d'y être. \\ \hline
\end{tabularx}
\end{center}
Remarque de grammaire : après \eng{would you like to} on met l'\textbf{infinitif} (\emph{to go}), mais après \eng{what about} et \eng{can't stand} on met le verbe en \textbf{-ing} (\emph{playing}, \emph{swimming}).
\end{propriete}

\begin{attention}[Pièges des francophones dans le dialogue]
\begin{itemize}
\item Ne dites pas \emph{Do you want to go swimming with me?} pour inviter un ami : c'est trop direct. Préférez \eng{Would you like to\ldots?} ou \eng{What about\ldots?}
\item \eng{What about} est suivi du verbe en \textbf{-ing} : \eng{What about \textbf{playing} chess?} et non \emph{What about to play chess?}
\item Pour l'heure, on utilise la préposition \eng{at} : \eng{at ten o'clock} ; pour le lieu de rendez-vous : \eng{in front of the centre}, \eng{at the youth centre}.
\item \eng{I'm looking forward \textbf{to} it} : ici \emph{to} est une préposition ; si un verbe suit, il prend \textbf{-ing} : \eng{I'm looking forward to \textbf{seeing} you.}
\end{itemize}
\end{attention}

\courssub{Production écrite guidée : \emph{My favourite recreational activity}}

Consigne type : \eng{Write a paragraph of 6 to 8 lines about your favourite recreational activity. Say what the activity is, who you do it with, where and when, and explain why you like it using \emph{because}.}

\begin{popfigure}
\begin{center}
\begin{tikzpicture}[node distance=0.5cm, every node/.style={font=\small}]
\node[draw=popblue, fill=popblueL, rounded corners, align=center, minimum width=4.2cm] (a) {\textbf{1. Activity}\\ \eng{My favourite activity is\ldots}};
\node[draw=popgreen, fill=popgreenL, rounded corners, align=center, minimum width=4.2cm, below=of a] (b) {\textbf{2. Who / Where / When}\\ \eng{I play it with\ldots\ at\ldots\ every\ldots}};
\node[draw=poporange, fill=poporangeL, rounded corners, align=center, minimum width=4.2cm, below=of b] (c) {\textbf{3. Why}\\ \eng{I like it \textbf{because}\ldots}};
\node[draw=poppurple, fill=poppurpleL, rounded corners, align=center, minimum width=4.2cm, below=of c] (d) {\textbf{4. Conclusion}\\ \eng{That is why I always\ldots}};
\draw[-{Stealth}, thick, popdark] (a) -- (b);
\draw[-{Stealth}, thick, popdark] (b) -- (c);
\draw[-{Stealth}, thick, popdark] (c) -- (d);
\end{tikzpicture}
\end{center}
\legende{Plan du paragraphe : Activity $\rightarrow$ Who/Where/When $\rightarrow$ Why (because\ldots) $\rightarrow$ Conclusion}
\end{popfigure}

\begin{methode}[Rédiger le paragraphe en quatre étapes]
\begin{enumerate}
\item \textbf{Nommer l'activité} : \eng{My favourite recreational activity is playing football.}
\item \textbf{Dire avec qui, où et quand} : \eng{I play it with my classmates on the school pitch every Friday afternoon.}
\item \textbf{Expliquer pourquoi avec} \eng{because} : \eng{I like it because it keeps me fit and because I make many friends.}
\item \textbf{Conclure} : \eng{That is why I never miss a training session.}
\end{enumerate}
\end{methode}

\begin{exemplebox}[Modèle de paragraphe rédigé]
My favourite recreational activity is singing in the school choir. I sing with twenty other students in the school hall every Wednesday and Saturday afternoon. Our conductor, Mr Etoundi, teaches us English and French songs. I like this activity because music helps me to relax after class, and because I have made many new friends in the choir. Last month, we even sang at the youth festival in Yaoundé. That is why I always look forward to our rehearsals.

« Mon activité de loisir préférée est de chanter dans la chorale de l'école. Je chante avec vingt autres élèves dans la salle de l'école chaque mercredi et samedi après-midi. Notre chef de chœur, M. Etoundi, nous apprend des chansons en anglais et en français. J'aime cette activité parce que la musique m'aide à me détendre après les cours, et parce que je me suis fait beaucoup de nouveaux amis à la chorale. Le mois dernier, nous avons même chanté au festival de la jeunesse à Yaoundé. C'est pourquoi j'attends toujours nos répétitions avec impatience. »
\end{exemplebox}

\begin{exoresolu}[Compléter un paragraphe]
Énoncé : complétez ce paragraphe avec les mots suivants : \emph{because, pitch, fan, every, with}.

My favourite activity is watching football. I am a great \ldots\ of Canon of Yaoundé. \ldots\ Saturday, I go to the \ldots\ \ldots\ my father. I love this activity \ldots\ it is exciting and we have a good time together.
\tcblower
My favourite activity is watching football. I am a great \textbf{fan} of Canon of Yaoundé. \textbf{Every} Saturday, I go to the \textbf{pitch} \textbf{with} my father. I love this activity \textbf{because} it is exciting and we have a good time together.

Points de méthode : \eng{fan} demande l'article indéfini (\emph{a great fan}) ; \eng{every} + jour exprime l'habitude (présent simple) ; \eng{because} introduit la justification, jamais \emph{because of} devant une phrase complète avec verbe conjugué.
\end{exoresolu}

\courssub{Réemploi oral : le sondage \emph{Find someone who\ldots}}

\begin{experience}[Class survey — Find someone who\ldots]
\textbf{Durée :} 10 minutes. \textbf{Objectif :} réemployer le lexique et le présent simple à l'oral.

\textbf{Déroulement :}
\begin{enumerate}
\item Chaque élève reçoit une grille : \emph{Find someone who\ldots\ plays chess / goes swimming / can't stand video games / sings in a choir / is a football fan / plays draughts / hangs out with friends on Saturdays.}
\item Les élèves circulent dans la classe et posent des questions complètes : \eng{Do you play chess?} — \eng{Do you go swimming?} — \eng{Can you stand video games?}
\item Si le camarade répond \eng{Yes, I do.}, on note son nom ; sinon on passe à un autre élève.
\item \textbf{Bilan collectif} : chaque élève rapporte une phrase à la troisième personne : \eng{Amina plays chess.} / \eng{Boris can't stand video games.} / \eng{Junior goes swimming every Sunday.}
\end{enumerate}
\textbf{Rappel} : question avec \eng{Do you\ldots?}, réponse courte \eng{Yes, I do. / No, I don't.}, et \textbf{-s} à la troisième personne au bilan.
\end{experience}

\begin{attention}[Les trois erreurs à éviter dans le sondage]
\begin{itemize}
\item \emph{You play chess?} $\rightarrow$ il faut l'auxiliaire : \eng{\textbf{Do} you play chess?}
\item \emph{Yes, I play.} $\rightarrow$ réponse courte avec l'auxiliaire : \eng{Yes, I \textbf{do}.}
\item \emph{Amina play chess.} $\rightarrow$ troisième personne : \eng{Amina \textbf{plays} chess.}
\end{itemize}
\end{attention}

\begin{savaistu}[Les loisirs au Baccalauréat]
Au Bac A, les textes de compréhension écrite parlent très souvent de la \cle{vie quotidienne} et des \cle{loisirs} : sport, musique, jeux, sorties en famille ou entre amis. Maîtriser ce champ lexical (\emph{pitch, choir, fan, to hang out, to relax}) vous aide donc \textbf{deux fois} : d'abord à \textbf{comprendre} le texte de l'épreuve de reading, ensuite à \textbf{rédiger} votre propre composition, car les sujets d'expression écrite demandent fréquemment de décrire vos activités préférées ou un week-end mémorable. Le vocabulaire appris aujourd'hui est un investissement direct pour l'examen !
\end{savaistu}

\begin{aretenir}[Bilan de la section]
\begin{itemize}
\item Lire un texte de loisirs : questions d'abord, mots-clés ensuite, réponses courtes au présent simple.
\item Inviter et planifier : \eng{Would you like to\ldots?}, \eng{What about\ldots -ing?}, \eng{Let's meet at\ldots}
\item Rédiger un paragraphe en quatre étapes : activité $\rightarrow$ avec qui / où / quand $\rightarrow$ pourquoi (\eng{because}) $\rightarrow$ conclusion.
\item À l'oral : \eng{Do you\ldots?} / \eng{Yes, I do.} / bilan avec le \textbf{-s} de la troisième personne.
\end{itemize}
\end{aretenir}

\newpage

\coursec{Activité d'intégration}

\courssub{Objectifs de l'activité}

À l'issue de cette activité d'intégration, tu seras capable de :
\begin{itemize}
\item mener un mini-sondage en anglais sur les loisirs, en posant des questions correctes au présent simple (\eng{What do you do in your free time? Do you fancy playing chess?}) ;
\item réemployer à l'oral et à l'écrit le vocabulaire des activités récréatives et les collocations avec \cle{play}, \cle{do} et \cle{go} ;
\item utiliser les expressions de goût (\eng{enjoy, be fond of, be keen on, can't stand}) et au moins deux expressions idiomatiques de la leçon (\eng{to have a good time, to kill time}) ;
\item compter et rapporter des résultats en phrases complètes (\eng{Five students play football at the weekend.}) ;
\item produire collectivement une affiche claire et la présenter à l'oral en une minute.
\end{itemize}

\courssub{Situation}

Le club d'anglais de ton lycée de Yaoundé prépare la « \eng{Languages Open Day} ». Le proviseur a demandé à chaque classe de Première de présenter un panneau en anglais sur la vie des élèves. Ta classe a choisi le thème des loisirs. Pour que le panneau soit fondé sur des faits réels, vous devez d'abord interroger vos camarades sur leurs activités du week-end, puis rédiger une affiche intitulée « \eng{Our Class at the Weekend} ».

\begin{textebox}[Message du président du club d'anglais]
Dear Form Five students,

Our school is getting ready for the Languages Open Day next month. Every class must display a poster in English. Your class will present “Our Class at the Weekend”.

First, carry out a short survey: ask your schoolmates what they do in their free time. Do they play football or draughts? Do they go swimming or jogging? Do they do judo or aerobics? Do they listen to music or watch films?

Then, count the answers and write your poster. Use full sentences, correct collocations with \emph{play}, \emph{do} and \emph{go}, and expressions of likes and dislikes. Finally, present your poster to the class in one minute. The class will vote for the clearest poster with the richest vocabulary.

Good luck, and have a good time!

The English Club President
\end{textebox}

\textbf{Glossaire :} \emph{to display} (v., afficher) ; \emph{to carry out a survey} (exp., mener une enquête) ; \emph{schoolmate} (n., camarade de classe) ; \emph{draughts} (n., jeu de dames) ; \emph{jogging} (n., footing).

\courssub{Consignes}

Travaille en groupe de quatre. Suis les étapes dans l'ordre.

\begin{enumerate}
\item \textbf{Prépare la grille du sondage.} Recopie cette liste de huit activités et vérifie que chaque verbe est correct (\eng{play / do / go}) : \eng{play football, go swimming, do judo, play draughts, listen to music, go jogging, do aerobics, watch films}. Ajoute deux activités de ton choix avec la bonne collocation.
\item \textbf{Mène l'enquête.} Chaque membre du groupe interroge au moins six camarades en anglais, avec les questions du modèle : \eng{What do you do in your free time? Do you fancy playing draughts? Are you keen on swimming?} Note les réponses dans la grille (un bâton par réponse).
\item \textbf{Compte les résultats.} Rédige une phrase complète par activité, au présent simple : \eng{Five students play football at the weekend. Two students go swimming. Nobody does judo.} Attention à l'accord du verbe avec \eng{nobody} et \eng{everybody}.
\item \textbf{Rédige l'affiche.} Sur une feuille, écris le titre « \eng{Our Class at the Weekend} », puis huit à dix phrases présentant les résultats. Ton affiche doit contenir : au moins trois collocations correctes avec \eng{play / do / go}, deux expressions de goût (\eng{enjoy, be fond of, be keen on, can't stand}) et deux expressions idiomatiques de la leçon (\eng{to have a good time, to kill time, to get together}).
\item \textbf{Présente l'affiche à l'oral.} En une minute, un ou deux membres du groupe lisent les résultats à la classe, avec une prononciation soignée (par exemple \eng{leisure} « \textbf{lé-jeur} », \eng{swimming} « \textbf{soui-ming} »).
\item \textbf{Vote.} La classe vote pour l'affiche la plus claire et la plus riche en vocabulaire correct. Justifie ton vote en une phrase en anglais : \eng{I vote for Group 2 because they use five correct collocations.}
\end{enumerate}

\courssub{Ressources à mobiliser}

\begin{aretenir}[Collocations play / do / go]
\begin{center}
\begin{tabularx}{\linewidth}{|l|X|}\hline
\rowcolor{popblueL} \textbf{Verbe} & \textbf{Emploi et exemples} \\ \hline
\cle{play} & jeux et sports avec ballon ou adversaire : \eng{play football, play chess, play the guitar} \\ \hline
\cle{do} & activités individuelles, arts martiaux : \eng{do judo, do aerobics, do gymnastics} \\ \hline
\cle{go} & activités en \emph{-ing} : \eng{go swimming, go jogging, go dancing} \\ \hline
\end{tabularx}
\end{center}
\end{aretenir}

\begin{aretenir}[Expressions utiles pour le sondage et l'affiche]
\begin{itemize}
\item Poser la question : \eng{What do you do in your free time? Do you fancy going swimming? Are you keen on music?}
\item Exprimer le goût : \eng{enjoy + -ing} (\eng{She enjoys listening to music.}), \eng{be fond of, be keen on, can't stand + -ing}.
\item Idiomes de la leçon : \eng{to have a good time} (passer un bon moment), \eng{to kill time} (tuer le temps), \eng{to get together} (se retrouver).
\item Rapporter les résultats : \eng{Five students play football. Most students enjoy music. Nobody does gymnastics.}
\end{itemize}
\end{aretenir}

\begin{attention}[Pièges à éviter dans ton affiche]
\begin{itemize}
\item Pas de \eng{to} après \eng{go} devant un verbe en \emph{-ing} : \eng{go swimming}, jamais \emph{go to swimming}.
\item \eng{enjoy} est toujours suivi d'un verbe en \emph{-ing} : \eng{enjoy playing}, jamais \emph{enjoy to play}.
\item \eng{nobody} et \eng{everybody} prennent le verbe au singulier : \eng{Nobody does judo.}
\item Faux ami : \eng{actually} signifie « en fait », pas « actuellement » ; pour « actuellement », dis \eng{at the moment}.
\end{itemize}
\end{attention}

\courssub{Proposition de production et corrigé commenté}

\begin{exoresolu}[Affiche modèle : « Our Class at the Weekend »]
Rédige l'affiche de ton groupe (huit à dix phrases) à partir des résultats de ton sondage, avec au moins trois collocations \eng{play / do / go}, deux expressions de goût et deux expressions idiomatiques.
\tcblower
\textbf{Production modèle (poster text) :}

\eng{OUR CLASS AT THE WEEKEND}

\eng{We asked twenty-four students about their free time. Here are the results.}

\eng{Ten students play football at the weekend, and six play draughts with their friends. Four students go swimming at the municipal pool, and three go jogging on Sunday morning. Two students do judo at the sports centre, but nobody does gymnastics. Most students enjoy listening to music, and many are fond of watching films at home.}

\eng{Our classmates are keen on outdoor games, but they can't stand doing nothing. At the weekend, we get together, we relax and we always have a good time. Free time is never wasted: it is the best way to kill time happily!}

\textbf{Commentaire des choix de langue :}
\begin{itemize}
\item Collocations correctes : \eng{play football, play draughts, go swimming, go jogging, do judo, do gymnastics} — le verbe change selon le type d'activité.
\item Expressions de goût : \eng{enjoy listening} (\eng{enjoy + -ing}), \eng{are fond of watching}, \eng{are keen on}, \eng{can't stand doing} — quatre structures variées, toutes suivies du gérondif.
\item Expressions idiomatiques réemployées : \eng{get together} et \eng{have a good time} ; \eng{kill time} clôt l'affiche avec humour.
\item Grammaire : présent simple partout (habitudes), verbe au singulier après \eng{nobody} (\eng{nobody does}), adverbes de fréquence bien placés (\eng{always have}).
\item Cohésion : \eng{and, but, here are the results} relient les phrases et donnent un rythme d'affiche.
\end{itemize}
\end{exoresolu}

\courssub{Bilan de l'activité}

Grâce à cette enquête et à cette affiche, tu as réemployé dans une situation réelle tout le lexique de la leçon : les noms d'activités récréatives, les collocations avec \cle{play}, \cle{do} et \cle{go}, les expressions de goût et les idiomes comme \eng{have a good time} et \eng{kill time}. Tu as aussi pratiqué le présent simple des habitudes, l'accord avec \eng{nobody} et \eng{most students}, et la prononciation des mots difficiles. Vérifie maintenant que tu sais : poser une question sur les loisirs sans calquer le français, choisir le bon verbe devant chaque activité, et rapporter un résultat chiffré en phrase complète. Ces compétences te serviront pour l'épreuve d'expression orale comme pour la rédaction d'un court texte sur ta vie quotidienne.

\newpage

\coursec{Méthodes et exercices résolus}

\coursec{Lesson 11 — Vocabulary: Words and expressions related to taking part in recreational activities}

\courssub{Champs lexicaux : les activités récréatives}

\begin{methode}[Identifier et employer le vocabulaire des loisirs]
Pour employer correctement le vocabulaire des loisirs, suis ces étapes :
\begin{enumerate}
\item \textbf{Repère le type d'activité} : sport d'équipe (\eng{football, volleyball, basketball, handball}), sport individuel (\eng{swimming, jogging, athletics, boxing}), jeu d'esprit ou de société (\eng{chess, cards, ludo, scrabble}), activité culturelle (\eng{dancing, singing, reading, acting}), activité de plein air (\eng{camping, hiking, fishing, cycling}).
\item \textbf{Vérifie la nature du mot} : nom (\eng{a game, a match, a hobby, a tournament, a spectator}), verbe (\eng{to relax, to compete, to train, to cheer, to coach}), adjectif (\eng{relaxing, competitive, entertaining, amateur, professional}), adverbe (\eng{competitively, regularly}).
\item \textbf{Associe le mot à son contexte} : on joue \eng{a match}, on gagne \eng{a trophy} ou \eng{a medal}, on pratique \eng{a hobby}, on assiste à \eng{a concert} ou \eng{a sporting event}, on encourage \eng{a team}.
\item \textbf{Enrichis ta phrase} avec un lieu, un moment et une fréquence : \eng{I play football with my friends at the school field every Saturday.}
\item \textbf{Relis} pour vérifier l'orthographe, l'accord du verbe et la collocation.
\end{enumerate}
\end{methode}

\begin{exemplebox}[Vocabulaire de base — Recreational Activities]
\begin{center}
\begin{tabularx}{\linewidth}{|X|X|X|}
\hline
\rowcolor{popblueL} \textbf{Catégorie} & \textbf{English} & \textbf{Français} \\
\hline
Sports d'équipe & football, volleyball, basketball, handball, rugby & football, volley-ball, basket-ball, handball, rugby \\
\hline
Sports individuels & swimming, jogging, running, athletics, boxing, cycling & natation, footing, course à pied, athlétisme, boxe, cyclisme \\
\hline
Jeux & chess, cards, ludo, scrabble, draughts & échecs, cartes, ludo, scrabble, dames \\
\hline
Activités culturelles & dancing, singing, acting, reading, playing music & danse, chant, théâtre, lecture, musique \\
\hline
Plein air & camping, hiking, fishing, trekking & camping, randonnée, pêche, trekking \\
\hline
Noms clés & match, tournament, championship, training, coach, referee, spectator, fan, supporter, trophy, medal, prize & match, tournoi, championnat, entraînement, entraîneur, arbitre, spectateur, fan, supporter, trophée, médaille, prix \\
\hline
Verbes clés & play, do, go, train, compete, win, lose, beat, cheer, support, relax, unwind & jouer, pratiquer, aller (faire), s'entraîner, concourir, gagner, perdre, battre, encourager, soutenir, se détendre, décompresser \\
\hline
Adjectifs clés & relaxing, competitive, entertaining, exciting, boring, tiring, amateur, professional, fit, exhausted & relaxant, compétitif, divertissant, passionnant, ennuyeux, fatigant, amateur, professionnel, en forme, épuisé \\
\hline
\end{tabularx}
\end{center}
\end{exemplebox}

\begin{exoresolu}[Compléter avec le vocabulaire des loisirs]
\textbf{Énoncé} — Complete the sentences with a suitable word from the box: \emph{hobby, match, tournament, training, spectators, chess, trophy, coach, cheer, fit}.
\begin{enumerate}
\item My favourite \ldots\ldots\ldots is reading novels by African writers.
\item Our school won the football \ldots\ldots\ldots organised in Yaoundé last term.
\item Thousands of \ldots\ldots\ldots watched the cup final at the Ahmadou Ahidjo stadium.
\item The team has \ldots\ldots\ldots every Tuesday and Thursday afternoon.
\item My grandfather taught me how to play \ldots\ldots\ldots when I was ten.
\item The \ldots\ldots\ldots gave the players tactical advice at half-time.
\item The fans started to \ldots\ldots\ldots loudly when the goal was scored.
\item She does yoga three times a week to keep \ldots\ldots\ldots
\end{enumerate}
\tcblower
\textbf{Solution détaillée}
\begin{enumerate}
\item \textbf{hobby} : le contexte « favourite » + une activité régulière (reading) désigne un passe-temps. \eng{My favourite hobby is reading novels by African writers.} « Mon passe-temps préféré est de lire des romans d'écrivains africains. »
\item \textbf{tournament} : une compétition organisée avec plusieurs équipes ; « organised » et « won » orientent vers ce mot. \eng{Our school won the football tournament organised in Yaoundé last term.} « Notre école a remporté le tournoi de football organisé à Yaoundé le trimestre dernier. »
\item \textbf{spectators} : ce sont les personnes qui regardent (\eng{watched}) un match au stade. Attention : \eng{spectators} (public neutre), pas \eng{supporters} (ceux qui encouragent une équipe précise).
\item \textbf{training} : la fréquence « every Tuesday and Thursday » indique des séances d'entraînement. On dit \eng{to have training} ou \eng{to train}.
\item \textbf{chess} : « play » + un jeu de société ; on dit \eng{play chess} sans article.
\item \textbf{coach} : personne qui donne des conseils tactiques (\eng{tactical advice}). \eng{The coach gave the players tactical advice at half-time.}
\item \textbf{cheer} : verbe signifiant « acclamer, encourager ». \eng{The fans started to cheer loudly when the goal was scored.}
\item \textbf{fit} : adjectif signifiant « en forme » ; collocation \eng{keep fit}. \eng{She does yoga three times a week to keep fit.}
\end{enumerate}
\textbf{Conclusion} : toujours s'appuyer sur les indices du contexte (verbe, fréquence, lieu, fonction) pour choisir le mot exact du champ lexical.
\end{exoresolu}

\courssub{Collocations : play, do, go et verbes du thème}

\begin{methode}[Choisir entre play, do et go]
\begin{enumerate}
\item \textbf{\eng{play} + sport avec ballon/ballon ou jeu} : \eng{play football, play basketball, play volleyball, play cards, play chess, play drums, play the guitar}. Règle : jeux et sports où l'on « joue » (opposition, balle, dés, cartes).
\item \textbf{\eng{do} + activité individuelle, discipline, art martial ou exercice} : \eng{do athletics, do karate, do judo, do gymnastics, do yoga, do exercise, do sport}. Règle : activités qu'on « pratique » comme une discipline.
\item \textbf{\eng{go} + verbe en \emph{-ing}} : \eng{go swimming, go dancing, go hiking, go fishing, go jogging, go cycling, go camping}. Règle : activités de déplacement ou de loisir exprimées par un gérondif.
\item \textbf{Mémorise les exceptions courantes} : \eng{do yoga}, \eng{go for a walk}, \eng{go for a run}, \eng{go for a swim}, \eng{play music} (général), \eng{do a puzzle}.
\item \textbf{Interdiction formelle} : jamais \eng{make sport} (calque du français « faire du sport ») ; dis \eng{do sport} (général) ou \eng{play sports} (pluriel pour plusieurs disciplines).
\end{enumerate}
\end{methode}

\begin{exemplebox}[Tableau récapitulatif des collocations]
\begin{center}
\begin{tabularx}{\linewidth}{|X|X|X|}
\hline
\rowcolor{popgreenL} \textbf{\eng{play} (+ nom)} & \textbf{\eng{do} (+ nom)} & \textbf{\eng{go} (+ -ing)} \\
\hline
football, basketball, volleyball, handball, rugby, tennis, badminton, table tennis, cards, chess, ludo, scrabble, draughts, drums, guitar, piano & athletics, gymnastics, karate, judo, yoga, exercise, sport, a puzzle, crosswords & swimming, dancing, hiking, fishing, jogging, running, cycling, camping, shopping, skiing \\
\hline
\end{tabularx}
\end{center}
\end{exemplebox}

\begin{exoresolu}[Play, do ou go ?]
\textbf{Énoncé} — Fill in the blanks with \emph{play}, \emph{do} or \emph{go} in the correct form.
\begin{enumerate}
\item Every Sunday, my friends and I \ldots\ldots\ldots swimming at the municipal pool in Douala.
\item My sister \ldots\ldots\ldots karate; she has a brown belt.
\item During the holidays, we \ldots\ldots\ldots volleyball on the beach in Kribi.
\item Would you like to \ldots\ldots\ldots dancing with us tonight?
\item My father \ldots\ldots\ldots not \ldots\ldots\ldots any sport because of his knee. (negative)
\item Look! The children \ldots\ldots\ldots chess under the mango tree. (present continuous)
\item We usually \ldots\ldots\ldots for a walk after dinner. (exception)
\end{enumerate}
\tcblower
\textbf{Solution détaillée}
\begin{enumerate}
\item \textbf{go} : « swimming » est un gérondif en \emph{-ing} ; avec « every Sunday » (habitude), present simple : \eng{Every Sunday, my friends and I go swimming at the municipal pool in Douala.}
\item \textbf{does} : le karaté est un art martial, donc \eng{do} ; sujet « my sister » = 3e personne du singulier au present simple : \eng{My sister does karate; she has a brown belt.}
\item \textbf{play} : le volleyball est un sport de ballon : \eng{During the holidays, we play volleyball on the beach in Kribi.}
\item \textbf{go} : « dancing » est un gérondif ; après « would you like to », base verbale : \eng{Would you like to go dancing with us tonight?}
\item \textbf{does ... do} : négation au present simple avec l'auxiliaire \eng{does not} + base verbale \eng{do} : \eng{My father does not do any sport because of his knee.} Remarque : « any » s'emploie à la forme négative.
\item \textbf{are playing} : present continuous (action en cours \eng{Look!}) : \eng{Look! The children are playing chess under the mango tree.}
\item \textbf{go} : exception \eng{go for a walk} : \eng{We usually go for a walk after dinner.}
\end{enumerate}
\textbf{Conclusion} : la collocation dépend de la nature de l'activité (ballon/jeu = \eng{play} ; discipline = \eng{do} ; \emph{-ing} = \eng{go}), puis on conjugue selon le sujet et le temps.
\end{exoresolu}

\courssub{Expressions idiomatiques et phrases utiles}

\begin{definition}[Expressions figées du thème « loisirs »]
\begin{itemize}
\item \eng{take part in} / \eng{participate in} : participer à (\eng{Many students took part in the tournament.})
\item \eng{join a club / a team} : s'inscrire dans un club / une équipe (\eng{He joined the drama club.})
\item \eng{keep fit} / \eng{stay in shape} : rester en forme (\eng{She jogs to keep fit.})
\item \eng{have fun} / \eng{have a good time} : s'amuser, passer un bon moment (\eng{We had a lot of fun at the party.})
\item \eng{spend time + -ing} : passer du temps à (\eng{I spend my weekends reading.})
\item \eng{unwind} / \eng{relax} : se détendre, décompresser (\eng{Music helps me unwind after school.})
\item \eng{cheer for} / \eng{support} : encourager (une équipe) (\eng{The crowd cheered for the local team.})
\item \eng{win a trophy / a medal / a prize} : gagner un trophée / une médaille / un prix
\item \eng{beat an opponent / a team} : battre un adversaire / une équipe (\eng{We beat our rivals 2-1.})
\item \eng{break a record} : battre un record
\item \eng{train hard} : s'entraîner dur
\item \eng{be a fan of} : être fan de (\eng{I am a big fan of Canon de Yaoundé.})
\end{itemize}
\end{definition}

\begin{exemplebox}[Dialogue : Organiser un match amical]
\textbf{Context} : Two classmates, \eng{Alice} and \eng{Brian}, discuss weekend plans.
\begin{verbatim}
Alice:  Hey Brian, what are you up to this Saturday?
Brian:  I'm going to play football with the neighbourhood team. Want to join us?
Alice:  I'd love to, but I have a dance rehearsal in the morning.
Brian:  No problem. The match starts at 3 p.m. at the college field.
Alice:  Great! I'll come and cheer for you. Which team do you support?
Brian:  We're playing against the "Lions". It's always a tough match.
Alice:  Well, good luck! I hope you win the trophy.
Brian:  Thanks! We train hard every evening. See you there!
\end{verbatim}
\textbf{Traduction} :
Alice : Salut Brian, tu fais quoi ce samedi ?
Brian : Je vais jouer au foot avec l'équipe du quartier. Tu veux te joindre à nous ?
Alice : J'adorerais, mais j'ai une répétition de danse le matin.
Brian : Pas de souci. Le match commence à 15h au terrain du collège.
Alice : Super ! Je viendrai t'encourager. Quelle équipe tu supportes ?
Brian : On joue contre les « Lions ». C'est toujours un match difficile.
Alice : Eh bien, bonne chance ! J'espère que vous gagnerez le trophée.
Brian : Merci ! On s'entraîne dur tous les soirs. À tout à l'heure !
\end{exemplebox}

\courssub{Faux amis et pièges des francophones}

\begin{methode}[Déjouer les pièges de traduction]
\begin{enumerate}
\item \textbf{Méfie-toi des mots identiques (faux amis)} :
\begin{itemize}
\item \eng{actually} = « en fait, en réalité » (pas « actuellement » $\rightarrow$ \eng{currently}, \eng{at present}).
\item \eng{event} = « événement » (pas « éventuel » $\rightarrow$ \eng{possible}, \eng{potential}).
\item \eng{assist} = « aider, secourir » (pas « assister à » $\rightarrow$ \eng{attend}, \eng{watch}).
\item \eng{repose} (verbe) = « reposer (un objet), reposer en paix » (pas « se reposer » $\rightarrow$ \eng{rest}, \eng{relax}).
\item \eng{pass} (verbe) = « passer (le temps, un examen, devant) » ; « passer du temps » = \eng{spend time}.
\end{itemize}
\item \textbf{Attention aux calques structuraux} :
\begin{itemize}
\item « Faire du sport » $\neq$ \eng{make sport} $\rightarrow$ \cle{do sport} ou \cle{play sports}.
\item « Faire une promenade » $\neq$ \eng{make a walk} $\rightarrow$ \cle{go for a walk} / \cle{take a walk}.
\item « Passer un bon moment » $\neq$ \eng{pass a good time} $\rightarrow$ \cle{have a good time} / \cle{have fun}.
\item « S'amuser » $\neq$ \eng{amuse oneself} (souvent réflexif, soutenu) $\rightarrow$ \cle{have fun} / \cle{enjoy oneself}.
\item « Se détendre » $\neq$ \eng{detend oneself} $\rightarrow$ \cle{relax} / \cle{unwind}.
\end{itemize}
\item \textbf{Vérifie toujours} le verbe anglais dans un dictionnaire (collocation) avant de traduire mot à mot.
\end{enumerate}
\end{methode}

\begin{exoresolu}[Corriger les calques du français]
\textbf{Énoncé} — Each sentence contains a mistake caused by French interference. Correct it.
\begin{enumerate}
\item I make sport every weekend to stay healthy.
\item We passed a good time at the school party.
\item Actually, I am training for the athletics competition.
\item Many people assisted the match at the stadium.
\item I like to repose myself after the games.
\item The players are very competitives.
\item She practises the piano every day. (Attention : \eng{practise} v. / \eng{practice} n.)
\end{enumerate}
\tcblower
\textbf{Solution détaillée}
\begin{enumerate}
\item \eng{make sport} est un calque de « faire du sport ». Correction : \eng{I do sport every weekend to stay healthy.} (ou \eng{play sports}).
\item « passer un bon moment » se traduit par \eng{have a good time} (past simple \eng{had}) : \eng{We had a good time at the school party.} Le verbe \eng{pass} ne convient pas ici ; on dit \eng{spend time} pour « passer du temps ».
\item \eng{actually} = « en fait », faux ami d'« actuellement ». Correction : \eng{Currently, I am training for the athletics competition.} (ou \eng{At present, I am training...}).
\item \eng{assist} signifie « aider » ; « assister à » (être présent) = \eng{attend} ou \eng{watch} : \eng{Many people attended the match at the stadium.}
\item « se reposer » ne se dit pas \eng{repose oneself} (registre soutenu, « reposer en paix ») ; correction : \eng{I like to rest after the games.} ou \eng{I like to relax after the games.}
\item \eng{competitive} est un adjectif invariable (pas de \emph{s} au pluriel) : \eng{The players are very competitive.}
\item \eng{practise} (verbe, orthographe britannique) / \eng{practice} (nom). Phrase correcte : \eng{She practises the piano every day.} (US : \eng{practice} pour les deux).
\end{enumerate}
\textbf{Conclusion} : un mot anglais qui ressemble à un mot français n'a presque jamais exactement le même sens ; vérifie et privilégie les verbes simples et fréquents (\eng{do, play, go, have, take}).
\end{exoresolu}

\courssub{Familles de mots, prononciation et orthographe}

\begin{methode}[Former les mots de la même famille (Word Formation)]
\begin{enumerate}
\item \textbf{Identifie la nature demandée} dans la phrase (nom, verbe, adjectif, adverbe) grâce à la place du mot : après \eng{a/the/my} = nom ; après le sujet / auxiliaire = verbe ; avant un nom / après \eng{be/look/feel} = adjectif ; après un verbe / en fin de phrase = adverbe.
\item \textbf{Connais les suffixes productifs} :
\begin{itemize}
\item Nom (personne) : \emph{-er}, \emph{-or}, \emph{-ist}, \emph{-ant} (\eng{play} $\rightarrow$ \eng{player}, \eng{compete} $\rightarrow$ \eng{competitor}, \eng{cycle} $\rightarrow$ \eng{cyclist}, \eng{assist} $\rightarrow$ \eng{assistant}).
\item Nom (action/état) : \emph{-tion/-sion}, \emph{-ment}, \emph{-ance/-ence}, \emph{-al}, \emph{-ing} (\eng{compete} $\rightarrow$ \eng{competition}, \eng{entertain} $\rightarrow$ \eng{entertainment}, \eng{perform} $\rightarrow$ \eng{performance}, \eng{recreate} $\rightarrow$ \eng{recreation}, \eng{train} $\rightarrow$ \eng{training}).
\item Adjectif : \emph{-ive}, \emph{-al}, \emph{-ous}, \emph{-ful}, \emph{-less}, \emph{-ing}, \emph{-ed} (\eng{compete} $\rightarrow$ \eng{competitive}, \eng{recreate} $\rightarrow$ \eng{recreational}, \eng{excite} $\rightarrow$ \eng{exciting}, \eng{relax} $\rightarrow$ \eng{relaxing} / \eng{relaxed}).
\item Adverbe : \emph{-ly} (\eng{competitive} $\rightarrow$ \eng{competitively}).
\end{itemize}
\item \textbf{Attention aux modifications orthographiques} :
\begin{itemize}
\item Doublement de la consonne finale (voyelle courte + consonne) : \eng{win} $\rightarrow$ \eng{winner}, \eng{run} $\rightarrow$ \eng{runner}, \eng{swim} $\rightarrow$ \eng{swimmer}, \eng{begin} $\rightarrow$ \eng{beginning}.
\item Chute du \emph{e} muet : \eng{compete} $\rightarrow$ \eng{competition}, \eng{create} $\rightarrow$ \eng{creation}, \eng{recreate} $\rightarrow$ \eng{recreation}.
\item Changement \emph{y} $\rightarrow$ \emph{i} : \eng{happy} $\rightarrow$ \eng{happiness}, \eng{easy} $\rightarrow$ \eng{easily}.
\end{itemize}
\item \textbf{Relis la phrase} pour vérifier que le mot formé s'accorde (singulier/pluriel, temps verbal).
\end{enumerate}
\end{methode}

\begin{exemplebox}[Tableau des familles de mots — Recreational Activities]
\begin{center}
\begin{tabularx}{\linewidth}{|c|X|X|X|X|}
\hline
\rowcolor{poppurpleL} \textbf{Verbe} & \textbf{Nom (action)} & \textbf{Nom (personne)} & \textbf{Adjectif} & \textbf{Adverbe} \\
\hline
compete & competition & competitor & competitive & competitively \\
\hline
entertain & entertainment & entertainer & entertaining / entertained & entertainingly \\
\hline
relax & relaxation & relaxer (rare) & relaxing / relaxed & relaxedly \\
\hline
win & win / victory & winner & winning & --- \\
\hline
lose & loss & loser & losing & --- \\
\hline
train & training & trainer / coach & trained & --- \\
\hline
perform & performance & performer & performing & --- \\
\hline
recreate & recreation & --- & recreational & recreationally \\
\hline
support & support / supporter & supporter & supportive & supportively \\
\hline
cheer & cheer & cheerleader & cheerful & cheerfully \\
\hline
\end{tabularx}
\end{center}
\end{exemplebox}

\begin{attention}[Prononciation — Pièges fréquents (notation française approximative)]
\begin{itemize}
\item \eng{leisure} : « lé-jeur » (UK) / « li-jeur » (US) — pas « lai-jour ».
\item \eng{chaos} : « ké-os » — pas « tcha-os ».
\item \eng{muscle} : « masseul » — le \emph{c} est muet.
\item \eng{rhythm} : « ri-dme » — pas de voyelle après \emph{h}.
\item \eng{championship} : « tcha-mpi-eun-chip » — \emph{ch} = /tʃ/.
\item \eng{tournament} : « tour-na-ment » — \emph{ou} = /ɔː/.
\item \eng{athlete} : « a-tlèt » — deux syllabes, pas « a-thé-lète ».
\item \eng{medal} : « mé-dle » — pas « mé-dal ».
\item \eng{victory} : « vic-to-ri » — accent sur 1ère syllabe.
\item \eng{exercise} : « ex-er-saïz » — \eng{ex} + \eng{er} + \eng{cise}.
\end{itemize}
\end{attention}

\begin{exoresolu}[Word formation]
\textbf{Énoncé} — Complete each sentence with the correct form of the word in brackets.
\begin{enumerate}
\item The \ldots\ldots\ldots of the singing competition received a trophy. (win)
\item Reading is a very \ldots\ldots\ldots activity for me. (relax)
\item The school organised an evening of music and \ldots\ldots\ldots. (entertain)
\item Our players are very \ldots\ldots\ldots; they always want to win. (compete)
\item Many young \ldots\ldots\ldots dream of joining the national team. (play)
\item The \ldots\ldots\ldots lasted for three hours. (perform)
\item He is a very \ldots\ldots\ldots cyclist; he trains every day. (compete)
\item We need more \ldots\ldots\ldots for the school choir. (support)
\end{enumerate}
\tcblower
\textbf{Solution détaillée}
\begin{enumerate}
\item \textbf{winner} : après l'article \eng{the}, il faut un nom de personne ; \eng{win} $\rightarrow$ \eng{winner} avec doublement du \emph{n} (voyelle courte + consonne finale). \eng{The winner of the singing competition received a trophy.}
\item \textbf{relaxing} : avant le nom \eng{activity}, il faut un adjectif ; \eng{relaxing} décrit l'activité (alors que \eng{relaxed} décrirait une personne). \eng{Reading is a very relaxing activity for me.}
\item \textbf{entertainment} : après « of music and », il faut un nom ; suffixe \emph{-ment}. \eng{The school organised an evening of music and entertainment.}
\item \textbf{competitive} : après \eng{are}, adjectif ; suffixe \emph{-itive}. \eng{Our players are very competitive; they always want to win.}
\item \textbf{players} : après l'adjectif \eng{young}, nom de personne ; « many » impose le pluriel. \eng{Many young players dream of joining the national team.}
\item \textbf{performance} : sujet de \eng{lasted}, nom d'action ; \eng{perform} $\rightarrow$ \eng{performance}. \eng{The performance lasted for three hours.}
\item \textbf{competitive} : adjectif qualificatif après \eng{very} ; \eng{competitive} (pas \eng{competitives}, adjectif invariable). \eng{He is a very competitive cyclist...}
\item \textbf{supporters} : nom de personne au pluriel (contexte « we need more »). \eng{We need more supporters for the school choir.}
\end{enumerate}
\textbf{Conclusion} : repère d'abord la nature grammaticale attendue, applique le bon suffixe, surveille l'orthographe (doublement, chute \emph{e}) et le nombre.
\end{exoresolu}

\courssub{Le lexique en contexte : texte, dialogue et production}

\begin{textebox}[Blog post : The Mount Cameroon Race of Hope]
\eng{Every year in February, the city of Buea comes alive with the Mount Cameroon Race of Hope. It is one of the most exciting sporting events in Africa. Athletes from all over the world take part in this gruelling competition. They run up the steep slopes of the mountain and down again, covering about 38 kilometres.}

\eng{Last year, I went to Buea with my friends to watch the race. We arrived early in the morning to get a good spot near the finish line. The atmosphere was electric. Supporters were cheering, waving flags and blowing vuvuzelas. When the winner, a young Cameroonian athlete, crossed the line, the crowd went wild. He broke the course record!}

\eng{After the race, we spent the afternoon relaxing at the Limbe beach. It was the perfect way to unwind after such an intense day. This event is not just about sport; it is a celebration of endurance, culture and national pride. I am already looking forward to next year's edition.}
\textbf{Glossaire} : \eng{gruelling} (adj.) = éprouvant, exténuant ; \eng{steep slopes} = pentes raides ; \eng{spot} = place, endroit ; \eng{electric} (atmosphere) = électrique, survoltée ; \eng{vuvuzela} = vuvuzela (corne de supporter) ; \eng{crossed the line} = a franchi la ligne ; \eng{went wild} = s'est déchaînée ; \eng{unwind} = se détendre, décompresser ; \eng{endurance} = endurance ; \eng{looking forward to} = avoir hâte de.
\end{textebox}

\begin{experience}[Expression orale : Interview d'un sportif]
\textbf{Consigne} : En binôme. L'élève A est journaliste pour \eng{The Cameroon Tribune}, l'élève B est un athlète local (footballeur, coureur, lutteur, etc.).
\textbf{Étapes} :
\begin{enumerate}
\item Préparer 5 questions en utilisant le vocabulaire de la leçon (\eng{training, tournament, win, lose, coach, team, supporter, keep fit, hobby, relax}).
\item Exemples : \eng{How often do you train?} / \eng{What is your favourite tournament?} / \eng{How do you relax after a match?}
\item Jouer l'interview (3 minutes). Échanger les rôles.
\item Restitution orale devant la classe : « \eng{He/She told me that...} » (discours rapporté optionnel).
\end{enumerate}
\end{experience}

\courssub{Méthode 4 : Réemployer le lexique à l'écrit (paragraphe guidé)}

\begin{methode}[Rédiger un paragraphe sur ses loisirs]
\begin{enumerate}
\item \textbf{Phrase d'introduction} : annonce ton activité préférée. Structure : \eng{My favourite recreational activity is + nom/-ing.}
\item \textbf{Développement} : dis avec qui, où, quand et pourquoi. Utilise les connecteurs \eng{because, moreover, in addition, for example, furthermore}.
\item \textbf{Lexique riche} : emploie au moins cinq mots du thème (\eng{team, training, tournament, relaxing, competitive, hobby, supporter, cheer, keep fit, unwind}).
\item \textbf{Conclusion} : donne ton opinion ou un bénéfice : \eng{It helps me relax and keep fit.}
\item \textbf{Relecture} : vérifie les collocations (\eng{play/do/go}), les temps (present simple pour les habitudes, past simple pour l'événement ponctuel) et l'orthographe.
\end{enumerate}
\end{methode}

\begin{exoresolu}[Rédaction guidée : « My favourite recreational activity »]
\textbf{Énoncé} — Write a paragraph of about 80--100 words on your favourite recreational activity. Say what it is, who you do it with, when and where, and why you enjoy it.
\tcblower
\textbf{Solution détaillée — modèle de production}

\eng{My favourite recreational activity is playing football. Every Saturday afternoon, I play with my classmates on the school field in Yaoundé. We usually train for one hour and then organise a short match between two teams. I enjoy football because it is exciting and it teaches me discipline and teamwork. Moreover, it helps me keep fit and unwind after a long week of studies. Last term, our team even won the inter-class tournament, which made us very proud. Football is more than a hobby for me; it is a passion.}

« Mon activité récréative préférée est le football. Chaque samedi après-midi, je joue avec mes camarades sur le terrain de l'école à Yaoundé. Nous nous entraînons généralement une heure, puis nous organisons un petit match entre deux équipes. J'aime le football parce que c'est passionnant et que cela m'apprend la discipline et l'esprit d'équipe. De plus, cela m'aide à rester en forme et à me détendre après une longue semaine d'études. Le trimestre dernier, notre équipe a même remporté le tournoi inter-classes, ce qui nous a rendus très fiers. Le football est plus qu'un passe-temps pour moi ; c'est une passion. »

\textbf{Justifications} : present simple pour les habitudes (\eng{I play, we train, it helps, it teaches}) ; past simple pour l'événement ponctuel (\eng{our team won, made}) ; collocation correcte \eng{play football} ; connecteurs \eng{because, moreover} ; lexique du thème : \eng{train, match, team, tournament, keep fit, unwind, hobby, discipline, teamwork}.
\end{exoresolu}

\courssub{Entraînement et évaluation}

\begin{exercice}[Exercices de consolidation]
\textbf{A. Vocabulary — Choose the correct word.}
\begin{enumerate}
\item The \ldots\ldots\ldots (spectators / supporters / assistants) clapped when the national anthem was played.
\item My brother does \ldots\ldots\ldots (karate / football / chess) three times a week at the dojo.
\item We \ldots\ldots\ldots (made / did / played) a lot of sport during the holidays.
\item She \ldots\ldots\ldots (passed / spent / had) her free time reading novels.
\item The team \ldots\ldots\ldots (won / beat / gained) their opponents 3-0 in the final.
\end{enumerate}

\textbf{B. Collocations — Complete with play, do or go (correct form).}
\begin{enumerate}
\item My parents often \ldots\ldots\ldots hiking in the mountains near Dschang.
\item \ldots\ldots\ldots you \ldots\ldots\ldots (play) the guitar? — Yes, I \ldots\ldots\ldots
\item She hates \ldots\ldots\ldots (do) gymnastics because it is too difficult.
\item We decided to \ldots\ldots\ldots (go) for a run early in the morning.
\item They \ldots\ldots\ldots (not / play) cards on Sundays; they prefer ludo.
\end{enumerate}

\textbf{C. Word Formation — Complete with the correct form.}
\begin{enumerate}
\item The \ldots\ldots\ldots (organise) of the event did a great job.
\item He is a very \ldots\ldots\ldots (talent) young \ldots\ldots\ldots (swim).
\item \ldots\ldots\ldots (participate) in sports builds character.
\item The match was very \ldots\ldots\ldots (excite) until the last minute.
\item She received a \ldots\ldots\ldots (victory) medal.
\end{enumerate}

\textbf{D. Translation — Translate into English.}
\begin{enumerate}
\item Actuellement, je me prépare pour le tournoi de basket-ball.
\item Nous avons passé un bon moment au concert hier soir.
\item Mon père ne fait pas de sport, il préfère lire.
\item Les supporters ont encouragé leur équipe jusqu'à la fin.
\item Jouer aux échecs est une activité relaxante.
\end{enumerate}

\textbf{E. Writing — Composition (100--120 words).}
\textit{« Write an article for your school magazine about the importance of recreational activities for students. Mention examples of activities, benefits (health, social, academic) and give advice. »}
\end{exercice}

\begin{corrige}[Exercices de consolidation]
\textbf{A.}
\begin{enumerate}
\item \textbf{spectators} (public neutre). \eng{The spectators clapped...}
\item \textbf{karate} (art martial $\rightarrow$ \eng{do}). \eng{My brother does karate...}
\item \textbf{did} (calque \eng{make sport} interdit). \eng{We did a lot of sport...}
\item \textbf{spent} (\eng{spend time}, pas \eng{pass time}). \eng{She spent her free time reading novels.}
\item \textbf{beat} (transitif direct : battre qqn ; \eng{win} prend \eng{a match/trophy}, \eng{gain} est rare pour le sport). \eng{The team beat their opponents 3-0...}
\end{enumerate}

\textbf{B.}
\begin{enumerate}
\item \textbf{go} (hiking = -ing). \eng{My parents often go hiking...}
\item \textbf{Do ... play} / \textbf{do} (present simple, question/short answer). \eng{Do you play the guitar? — Yes, I do.}
\item \textbf{doing} (gérondif après \eng{hate}). \eng{She hates doing gymnastics...}
\item \textbf{go} (exception \eng{go for a run}). \eng{We decided to go for a run...}
\item \textbf{don't play} (négation present simple). \eng{They don't play cards on Sundays...}
\end{enumerate}

\textbf{C.}
\begin{enumerate}
\item \textbf{organisers} (nom personne, pluriel après \eng{The}). \eng{The organisers of the event...}
\item \textbf{talented} (adj. après \eng{very}) / \textbf{swimmer} (nom personne, doublement \emph{m}). \eng{He is a very talented young swimmer.}
\item \textbf{Participating} / \textbf{Participation} (sujet : gérondif ou nom action). \eng{Participating in sports builds character.} / \eng{Participation in sports builds character.}
\item \textbf{exciting} (adjectif en \emph{-ing} car le match \emph{procure} l'émotion). \eng{The match was very exciting...}
\item \textbf{winner's} / \textbf{winning} (adjectif \eng{winning medal} ou génitif \eng{winner's medal}). \eng{She received a winning medal.} (ou \eng{gold medal}).
\end{enumerate}

\textbf{D.}
\begin{enumerate}
\item \eng{Currently, I am training for the basketball tournament.} (pas \eng{Actually})
\item \eng{We had a good time at the concert last night.} (pas \eng{passed})
\item \eng{My father does not do sport; he prefers reading.} (pas \eng{make sport})
\item \eng{The supporters cheered for their team until the end.} (ou \eng{cheered their team on})
\item \eng{Playing chess is a relaxing activity.} (gérondif sujet, adjectif \eng{relaxing})
\end{enumerate}

\textbf{E. — Modèle de rédaction}
\eng{The Importance of Recreational Activities for Students}

\eng{Recreational activities play a vital role in a student's life. They are not just a way to pass time; they are essential for our physical and mental well-being. Activities like football, dancing, swimming or even playing chess help us keep fit and unwind after long hours in class.}

\eng{Moreover, taking part in team sports teaches us discipline, teamwork and how to handle both victory and defeat. For example, when our school team won the inter-school tournament last year, we learned the value of hard training and mutual support. Socially, clubs and teams allow us to make new friends and develop leadership skills.}

\eng{In conclusion, I strongly advise every student to join a club or practise a sport regularly. It improves our health, boosts our academic performance by reducing stress, and makes school life more enjoyable. A healthy mind in a healthy body!}
\end{corrige}

\begin{attention}[Les 5 erreurs qui coûtent des points au Bac]
\begin{enumerate}
\item \textbf{Le calque \eng{make sport}} : « faire du sport » se dit \eng{do sport} ou \eng{play sports}. \eng{Make} ne se combine jamais avec \eng{sport}.
\item \textbf{Le faux ami \eng{actually}} : il signifie « en fait », jamais « actuellement ». Pour « actuellement », écris \eng{currently} ou \eng{at present}.
\item \textbf{La confusion play/do/go} : \eng{go} s'emploie uniquement avec les activités en \emph{-ing} (\eng{go swimming}), \eng{play} avec les ballons et les jeux (\eng{play football}), \eng{do} avec les disciplines (\eng{do karate}). Écris \eng{go swim} et c'est faux : il faut \eng{go swimming}.
\item \textbf{\eng{assist} pour « assister à »} : \eng{assist} veut dire « aider ». Pour « assister à un match », utilise \eng{attend} ou \eng{watch} : \eng{We attended the match.}
\item \textbf{L'orthographe des dérivés} : \eng{winner} avec deux \emph{n}, \eng{competition} (pas \eng{competion}), \eng{entertainment} (pas \eng{entertainement}), \eng{beginning} avec deux \emph{n}, \eng{swimmer} avec deux \emph{m}. Une faute sur ces mots fréquents est toujours pénalisée.
\end{enumerate}
\end{attention}

\begin{savaistu}[Culture : Le sport au Cameroun et dans le monde anglophone]
Le Cameroun est une nation sportive majeure en Afrique. Le football est roi : les \eng{Indomitable Lions} (Lions Indomptables) ont remporté la CAN 5 fois (1984, 1988, 2000, 2002, 2017) et ont marqué l'histoire à la Coupe du Monde 1990 (quart de finale). Roger Milla et Samuel Eto'o sont des légendes.
La \eng{Mount Cameroon Race of Hope} (Course de l'Espoir) est une course de montagne internationale unique, inscrite au patrimoine de l'athlétisme mondial.
Dans le monde anglophone, le \eng{cricket} (Royaume-Uni, Inde, Australie), le \eng{rugby} (Nouvelle-Zélande, Afrique du Sud), le \eng{baseball} (USA, Japon) et le \eng{American football} (USA) sont des sports culturels majeurs. Les \eng{Olympic Games} (Jeux Olympiques) et les \eng{Commonwealth Games} (Jeux du Commonwealth) sont les grands rendez-vous multi-sports où le Cameroun brille (boxe, haltérophilie, lutte).
\end{savaistu}

\newpage

\coursec{Exercices}

\courssub{Je vérifie mes connaissances}

\begin{exercice}[QCM : les bons verbes des loisirs]
Pour chaque phrase, choisis la bonne réponse (A, B ou C).
\begin{enumerate}
\item We \_\_\_\_\_\_ swimming every Saturday at the municipal pool. \quad A. play \quad B. do \quad C. go
\item She \_\_\_\_\_\_ karate twice a week at the club in Biyem-Assi. \quad A. plays \quad B. does \quad C. goes
\item They \_\_\_\_\_\_ chess during the break. \quad A. play \quad B. do \quad C. go
\item My father \_\_\_\_\_\_ jogging every morning before work. \quad A. plays \quad B. does \quad C. goes
\item The students \_\_\_\_\_\_ board games in the school club on Fridays. \quad A. play \quad B. do \quad C. go
\end{enumerate}
\end{exercice}

\begin{exercice}[Vrai ou faux ? Justifie ta réponse]
Lis les affirmations suivantes sur le vocabulaire des loisirs. Écris \eng{True} ou \eng{False} et justifie en une phrase en anglais.
\begin{enumerate}
\item \eng{“Actually” means “actuellement” in English.}
\item \eng{We say “to do athletics”, not “to play athletics”.}
\item \eng{“To assist a match” is the correct way to say “assister à un match”.}
\item \eng{“To have fun” means “s'amuser”.}
\item \eng{“Eventually” means “éventuellement”.}
\end{enumerate}
\end{exercice}

\begin{exercice}[Vocabulaire : trouve le mot]
Complète chaque définition par le mot anglais du thème des loisirs qui convient. Choisis dans la liste : \eng{hobby — leisure — tournament — spectator — teammate — entertainment}.
\begin{enumerate}
\item An activity you do regularly for pleasure in your free time: \_\_\_\_\_\_.
\item The time when you are not working or studying: \_\_\_\_\_\_.
\item A competition in which several players or teams play against each other: \_\_\_\_\_\_.
\item A person who watches a match or a show: \_\_\_\_\_\_.
\item A person who plays in the same team as you: \_\_\_\_\_\_.
\item Shows, films or games that amuse people: \_\_\_\_\_\_.
\end{enumerate}
\end{exercice}

\begin{exercice}[Familles de mots : la bonne forme]
Complète chaque phrase avec la forme correcte du mot entre parenthèses.
\begin{enumerate}
\item The \_\_\_\_\_\_ between the two schools was very exciting. (compete)
\item After a long week, I need some \_\_\_\_\_\_. (relax)
\item The concert was really \_\_\_\_\_\_; we all loved it. (enjoy)
\item My sister is a very \_\_\_\_\_\_ dancer. (talent)
\item Football is a very \_\_\_\_\_\_ activity in Cameroon. (popularity)
\item He plays the guitar \_\_\_\_\_\_. (beautiful)
\end{enumerate}
\end{exercice}

\courssub{Je m'entraîne}

\begin{exercice}[Traduction de phrases]
Traduis les phrases suivantes en anglais, en faisant attention aux faux amis et aux collocations.
\begin{enumerate}
\item Le samedi, je fais du sport avec mes amis au quartier.
\item Nous avons assisté au match de football au stade de Yaoundé.
\item En fait, elle préfère la lecture aux jeux vidéo.
\item Ils font de la natation à la piscine tous les dimanches.
\item Mon frère joue aux échecs mieux que moi.
\end{enumerate}
\end{exercice}

\begin{exercice}[Corrige les erreurs des francophones]
Chaque phrase contient une erreur typique d'un élève francophone. Réécris la phrase correctement.
\begin{enumerate}
\item \eng{I make sport every weekend.}
\item \eng{We assisted to the match last Sunday.}
\item \eng{I am agree with you: dancing is fun.}
\item \eng{She plays piano very well.}
\item \eng{Actually, I am studying for my exam.}
\item \eng{They do football in the evening.}
\end{enumerate}
\end{exercice}

\begin{exercice}[Texte à trous]
Complète le texte suivant avec les mots de la liste : \eng{go — play — hobby — takes part — relaxing — tournament — fans — do}.

\begin{textebox}[Arlette's free time]
Arlette is a student at Lycée de Nkol-Eton. Her favourite \_\_\_\_\_\_ is handball. She \_\_\_\_\_\_ in the school team and trains three times a week. At the weekend, she and her friends sometimes \_\_\_\_\_\_ swimming at the municipal pool. When she is tired, she prefers \_\_\_\_\_\_ activities like reading novels or listening to makossa music. Last month, her school organised a big \_\_\_\_\_\_ and many \_\_\_\_\_\_ came to watch the final. Arlette also likes to \_\_\_\_\_\_ board games with her little brother in the evening.
\end{textebox}
\end{exercice}

\begin{exercice}[Dialogue d'invitation : remets en ordre]
Les répliques de ce dialogue entre deux camarades sont mélangées. Remets-les dans le bon ordre (numérote de 1 à 6).
\begin{enumerate}
\item[\eng{a.}] \eng{Great! Let's meet at 4 p.m. at the sports complex.}
\item[\eng{b.}] \eng{Hi Sandra! Would you like to play table tennis with me this afternoon?}
\item[\eng{c.}] \eng{Perfect. See you there, and don't forget your racket!}
\item[\eng{d.}] \eng{I'd love to, but I have homework to finish first.}
\item[\eng{e.}] \eng{No problem. What about later, around four o'clock?}
\item[\eng{f.}] \eng{That suits me fine. Where shall we meet?}
\end{enumerate}
\end{exercice}

\courssub{Je me prépare au Bac}

\begin{exercice}[Compréhension de l'écrit et questions de langue]
Lis le texte suivant, puis réponds aux questions en anglais, par des phrases complètes.

\begin{textebox}[Weekend games in Douala]
Every Saturday afternoon, the esplanade near Bonanjo market in Douala comes alive. Dozens of young people gather there to take part in recreational activities. Some play football on the dusty ground, while others prefer board games under the mango trees. Songo'o, a seventeen-year-old student, never misses these meetings. “Playing with my friends helps me forget the stress of school,” he says with a smile. His sister Clarisse, however, is not interested in ball games. She belongs to a dance group that rehearses traditional and modern dances every weekend. Their mother, Mrs Etame, sells roasted corn nearby and watches the young people with pleasure. “When I was young,” she explains, “we also had our games, but today the children are lucky: the town council has built a small sports ground with real goalposts.” For the whole neighbourhood, these Saturday games are more than entertainment: they bring people together and keep the teenagers away from trouble.
\end{textebox}

\textbf{A. Comprehension questions}
\begin{enumerate}
\item Where do the young people of the text meet every Saturday? (2 marks)
\item Mention two recreational activities practised in the text. (2 marks)
\item Why does Songo'o like playing with his friends? Quote from the text. (2 marks)
\item What is Clarisse's hobby? (2 marks)
\item According to Mrs Etame, why are today's children lucky? (2 marks)
\end{enumerate}

\textbf{B. Language work}
\begin{enumerate}
\item Find in the text a word that means “practise again to improve”. (1 mark)
\item Give the noun corresponding to the verb \eng{entertain}. (1 mark)
\item Complete with \eng{play}, \eng{do} or \eng{go}: \eng{Clarisse and her friends \_\_\_\_\_\_ dancing every weekend.} (1 mark)
\item Correct the mistake: \eng{The teenagers assists to the games every Saturday.} (1 mark)
\item Turn into the plural: \eng{The spectator watches the match.} (1 mark)
\end{enumerate}
\end{exercice}

\begin{exercice}[Traduction]
Traduis le passage suivant en anglais. (10 marks)

\begin{textebox}[Texte à traduire]
Pendant les vacances, les jeunes de mon village participent à de nombreuses activités récréatives. Certains jouent au football sur le terrain du chef, d'autres font de la lutte traditionnelle. Le soir, les familles se réunissent pour écouter des contes et chanter ensemble. En fait, ces moments de loisirs sont très importants : ils renforcent l'amitié entre les jeunes et transmettent notre culture aux enfants.
\end{textebox}
\end{exercice}

\begin{exercice}[Expression écrite : composition]
\textbf{Topic:} \eng{Describe your favourite recreational activity. Say what it is, when and where you practise it, who you do it with, and why you enjoy it.}

Rédige une composition en anglais de \textbf{80 à 100 mots}. Tu dois :
\begin{enumerate}
\item présenter ton activité préférée (nom, lieu, fréquence) ;
\item employer au moins trois collocations correctes avec \eng{play}, \eng{do} ou \eng{go} ;
\item utiliser au moins deux expressions idiomatiques vues dans la leçon (par exemple \eng{to have fun}, \eng{to take part in}, \eng{to look forward to}) ;
\item soigner l'orthographe et la ponctuation.
\end{enumerate}
(10 marks)
\end{exercice}

\newpage

\coursec{Corrigés des exercices}

\begin{corrige}[Exercice 1 — QCM : les bons verbes des loisirs]
\begin{enumerate}
\item \textbf{C. go} — \eng{We go swimming every Saturday at the municipal pool.} « Nous allons nager tous les samedis à la piscine municipale. » Règle : \cle{go + verbe en -ing} pour les activités de déplacement ou de plein air (\eng{go swimming, go jogging, go fishing}).
\item \textbf{B. does} — \eng{She does karate twice a week at the club in Biyem-Assi.} « Elle fait du karaté deux fois par semaine au club de Biyem-Assi. » Règle : \cle{do + sport individuel ou de combat} (\eng{do karate, do judo, do gymnastics}). Attention à la 3\up{e} personne du singulier : \eng{she does}, prononcé « deuz ».
\item \textbf{A. play} — \eng{They play chess during the break.} « Ils jouent aux échecs pendant la récréation. » Règle : \cle{play + jeux et sports avec ballon ou adversaire} (\eng{play chess, play football, play cards}).
\item \textbf{C. goes} — \eng{My father goes jogging every morning before work.} « Mon père fait du jogging tous les matins avant le travail. » Même règle que la phrase 1 : \eng{go + -ing} ; à la 3\up{e} personne, \eng{go} devient \eng{goes}.
\item \textbf{A. play} — \eng{The students play board games in the school club on Fridays.} « Les élèves jouent à des jeux de société au club de l'école le vendredi. » \eng{Board games} = jeux de société, donc \eng{play}.
\end{enumerate}
\textbf{Point de vigilance :} on ne dit jamais \eng{play swimming} ni \eng{do chess}. Retiens le trio : \eng{play} (jeux, sports d'équipe, instruments), \eng{do} (sports individuels, arts martiaux), \eng{go} (activités en \eng{-ing}).
\end{corrige}

\begin{corrige}[Exercice 2 — Vrai ou faux ? Justifie ta réponse]
\begin{enumerate}
\item \textbf{False.} \eng{“Actually” means “in fact” or “really”, not “actuellement”; “actuellement” is translated by “currently” or “at present”.} Faux ami classique : \eng{actually} = « en fait », jamais « actuellement ».
\item \textbf{True.} \eng{Athletics is an individual sport, so we say “to do athletics”, like “to do gymnastics”.} Attention : \eng{athletics} se construit au singulier (\eng{Athletics is...}).
\item \textbf{False.} \eng{“To assist” means “to help”; the correct expression is “to attend a match” or “to watch a match”.} \eng{Assist} est un faux ami de « assister à » ; de plus, \eng{attend} se construit \textbf{sans préposition} : \eng{attend a match}, jamais \eng{attend to a match}.
\item \textbf{True.} \eng{“To have fun” means “s'amuser, prendre du plaisir”.} Exemple : \eng{We had fun at the party.} « Nous nous sommes bien amusés à la fête. » Passé irrégulier : \eng{have -- had -- had}.
\item \textbf{False.} \eng{“Eventually” means “finalement, à la fin”; “éventuellement” is translated by “possibly” or “if necessary”.} Autre faux ami à retenir.
\end{enumerate}
\end{corrige}

\begin{corrige}[Exercice 3 — Vocabulaire : trouve le mot]
\begin{enumerate}
\item \textbf{hobby} — un passe-temps. Pluriel : \eng{hobbies} (un \eng{-y} précédé d'une consonne devient \eng{-ies}).
\item \textbf{leisure} — les loisirs, le temps libre. Attention à la prononciation : « lè-jeur ». Expression utile : \eng{leisure activities} « activités de loisirs ».
\item \textbf{tournament} — un tournoi. Prononciation : « teur-ne-ment ».
\item \textbf{spectator} — un spectateur (celui qui regarde ; ne pas confondre avec \eng{player}, le joueur, ni avec \eng{fan}, le supporter).
\item \textbf{teammate} — un coéquipier. Mot composé : \eng{team} (équipe) + \eng{mate} (camarade), comme \eng{classmate} (camarade de classe).
\item \textbf{entertainment} — le divertissement, les distractions. Nom dérivé du verbe \eng{entertain} « divertir ».
\end{enumerate}
\end{corrige}

\begin{corrige}[Exercice 4 — Familles de mots : la bonne forme]
\begin{enumerate}
\item \textbf{competition} — \eng{The competition between the two schools was very exciting.} « La compétition entre les deux écoles était très passionnante. » Nom en \eng{-tion} dérivé du verbe \eng{compete}.
\item \textbf{relaxation} — \eng{After a long week, I need some relaxation.} « Après une longue semaine, j'ai besoin de détente. » Nom après le déterminant \eng{some}.
\item \textbf{enjoyable} — \eng{The concert was really enjoyable; we all loved it.} « Le concert était vraiment agréable ; nous l'avons tous adoré. » Adjectif en \eng{-able} après le verbe \eng{be}.
\item \textbf{talented} — \eng{My sister is a very talented dancer.} « Ma sœur est une danseuse très douée. » Adjectif en \eng{-ed} dérivé du nom \eng{talent}.
\item \textbf{popular} — \eng{Football is a very popular activity in Cameroon.} « Le football est une activité très populaire au Cameroun. » Adjectif ; \eng{popularity} est le nom correspondant.
\item \textbf{beautifully} — \eng{He plays the guitar beautifully.} « Il joue de la guitare magnifiquement. » Adverbe en \eng{-ly} pour modifier le verbe \eng{plays}.
\end{enumerate}
\textbf{Point de vigilance :} demande-toi toujours quelle nature de mot la phrase exige : nom, adjectif ou adverbe. Un adverbe modifie un verbe ; un adjectif modifie un nom ou suit le verbe \eng{be}.
\end{corrige}

\begin{corrige}[Exercice 5 — Traduction de phrases]
\begin{enumerate}
\item \eng{On Saturdays, I do sport with my friends in the neighbourhood.} (ou, à l'américaine, \eng{I play sports with my friends}). Jamais \eng{I make sport} : « faire du sport » = \eng{do sport}.
\item \eng{We attended the football match at the Yaoundé stadium.} (ou \eng{We watched the football match...}). Jamais \eng{assisted to} : \eng{assist} = aider, et \eng{attend} ne prend pas de préposition.
\item \eng{Actually, she prefers reading to video games.} (ou \eng{In fact, she prefers books to video games.}) \eng{Prefer ... to ...} : préposition \eng{to}, jamais \eng{than}. Ici « en fait » se traduit bien par \eng{actually}.
\item \eng{They go swimming at the pool every Sunday.} Collocation obligatoire : \eng{go swimming}.
\item \eng{My brother plays chess better than me.} (ou, plus soutenu, \eng{better than I do}). \eng{Play chess}, sans article devant les jeux et les sports.
\end{enumerate}
\end{corrige}

\begin{corrige}[Exercice 6 — Corrige les erreurs des francophones]
\begin{enumerate}
\item \eng{I make sport every weekend.} $\rightarrow$ \eng{I do sport every weekend.} (ou \eng{I play sports}). « Faire du sport » ne se traduit pas par \eng{make}.
\item \eng{We assisted to the match last Sunday.} $\rightarrow$ \eng{We attended the match last Sunday.} (ou \eng{We watched the match}). Double erreur : \eng{assist} = aider (faux ami), et \eng{attend} se construit sans préposition.
\item \eng{I am agree with you: dancing is fun.} $\rightarrow$ \eng{I agree with you: dancing is fun.} \eng{Agree} est un verbe en anglais, pas un adjectif : pas de \eng{be} devant.
\item \eng{She plays piano very well.} $\rightarrow$ \eng{She plays the piano very well.} Devant un instrument de musique, l'article \eng{the} est obligatoire : \eng{play the piano, play the guitar}.
\item \eng{Actually, I am studying for my exam.} $\rightarrow$ \eng{At present, I am studying for my exam.} (ou \eng{Currently, ...}). Faux ami : \eng{actually} = « en fait » ; « actuellement » = \eng{currently, at present}.
\item \eng{They do football in the evening.} $\rightarrow$ \eng{They play football in the evening.} Sport avec ballon : \eng{play}.
\end{enumerate}
\end{corrige}

\begin{corrige}[Exercice 7 — Texte à trous]
Texte complété :

\eng{Arlette is a student at Lycée de Nkol-Eton. Her favourite \textbf{hobby} is handball. She \textbf{takes part} in the school team and trains three times a week. At the weekend, she and her friends sometimes \textbf{go} swimming at the municipal pool. When she is tired, she prefers \textbf{relaxing} activities like reading novels or listening to makossa music. Last month, her school organised a big \textbf{tournament} and many \textbf{fans} came to watch the final. Arlette also likes to \textbf{play} board games with her little brother in the evening.}

« Arlette est élève au Lycée de Nkol-Eton. Son passe-temps préféré est le handball. Elle fait partie de l'équipe de l'école et s'entraîne trois fois par semaine. Le week-end, elle et ses amis vont parfois nager à la piscine municipale. Quand elle est fatiguée, elle préfère les activités relaxantes comme lire des romans ou écouter de la musique makossa. Le mois dernier, son école a organisé un grand tournoi et de nombreux supporters sont venus regarder la finale. Arlette aime aussi jouer aux jeux de société avec son petit frère le soir. »

\textbf{Justifications :} \eng{hobby} (passe-temps) ; \eng{takes part in} (participe à, fait partie de) ; \eng{go swimming} (collocation \eng{go + -ing}) ; \eng{relaxing} (adjectif en \eng{-ing} : « qui détend », à ne pas confondre avec \eng{relaxed} « détendu ») ; \eng{tournament} (tournoi) ; \eng{fans} (supporters) ; \eng{play board games} (collocation avec \eng{play}).
\end{corrige}

\begin{corrige}[Exercice 8 — Dialogue d'invitation : remets en ordre]
Ordre correct : \textbf{1. b — 2. d — 3. e — 4. f — 5. a — 6. c}

Dialogue reconstitué :
\begin{enumerate}
\item[\eng{1.}] \eng{Hi Sandra! Would you like to play table tennis with me this afternoon?}
\item[\eng{2.}] \eng{I'd love to, but I have homework to finish first.}
\item[\eng{3.}] \eng{No problem. What about later, around four o'clock?}
\item[\eng{4.}] \eng{That suits me fine. Where shall we meet?}
\item[\eng{5.}] \eng{Great! Let's meet at 4 p.m. at the sports complex.}
\item[\eng{6.}] \eng{Perfect. See you there, and don't forget your racket!}
\end{enumerate}
\textbf{Logique :} invitation (\eng{Would you like...?}) $\rightarrow$ acceptation avec réserve (\eng{I'd love to, but...}) $\rightarrow$ proposition d'un autre horaire (\eng{What about...?}) $\rightarrow$ question sur le lieu $\rightarrow$ proposition du lieu $\rightarrow$ conclusion. Retiens les formules : \eng{Would you like to...?} « Voudrais-tu... ? », \eng{I'd love to} « Avec plaisir », \eng{What about...?} « Et si... ? », \eng{That suits me fine} « Cela me convient parfaitement ».
\end{corrige}

\begin{corrige}[Exercice 9 — Compréhension de l'écrit et questions de langue]
\textbf{A. Comprehension questions} (barème indicatif : 2 points par réponse correcte et bien formulée)
\begin{enumerate}
\item \eng{They meet on the esplanade near Bonanjo market in Douala.} (2 marks)
\item \eng{They play football and board games.} (on accepte aussi \eng{dancing}). (2 marks)
\item \eng{Because “playing with my friends helps me forget the stress of school”.} (2 marks — la citation exacte est exigée)
\item \eng{Clarisse's hobby is dancing: she belongs to a dance group that rehearses traditional and modern dances.} (2 marks)
\item \eng{Because the town council has built a small sports ground with real goalposts.} (2 marks)
\end{enumerate}

\textbf{B. Language work} (1 point par item)
\begin{enumerate}
\item \textbf{rehearse} — « répéter, s'exercer pour s'améliorer ». Prononciation : « ri-heurse ».
\item \textbf{entertainment} — nom dérivé du verbe \eng{entertain} « divertir » (suffixe \eng{-ment}).
\item \eng{Clarisse and her friends \textbf{go} dancing every weekend.} (\eng{go + -ing} ; sujet pluriel, donc verbe sans \eng{-s}.)
\item \eng{The teenagers assists to the games every Saturday.} $\rightarrow$ \eng{The teenagers attend the games every Saturday.} (ou \eng{watch the games}). Double erreur : faux ami \eng{assist} et \eng{-s} de la 3\up{e} personne sur un sujet pluriel.
\item \eng{The spectator watches the match.} $\rightarrow$ \eng{The spectators watch the match.} Attention : au pluriel, le nom prend un \eng{-s} mais le verbe le perd.
\end{enumerate}
\textbf{Points de vigilance :} réponds toujours par une phrase complète ; recopie les citations sans les modifier ; vérifie l'accord sujet-verbe (pluriel = verbe sans \eng{-s}).
\end{corrige}

\begin{corrige}[Exercice 10 — Traduction]
\textbf{Proposition de traduction :}

\eng{During the holidays, the young people of my village take part in many recreational activities. Some play football on the chief's field, while others do traditional wrestling. In the evening, families gather to listen to folk tales and sing together. In fact, these leisure moments are very important: they strengthen friendship among the young people and pass on our culture to the children.}

\textbf{Barème indicatif (10 points) :} 2 points par phrase correctement traduite (5 phrases), avec déduction de 0,5 à 1 point par faute de grammaire, de collocation ou de faux ami.

\textbf{Points de vigilance :}
\begin{itemize}
\item « participent à » = \eng{take part in}, pas \eng{participate to} (\eng{participate in} est correct mais moins naturel ici).
\item « font de la lutte traditionnelle » = \eng{do traditional wrestling} (sport individuel : \eng{do}).
\item « se réunissent » = \eng{gather} ou \eng{get together}.
\item « En fait » = \eng{In fact} ou \eng{Actually} ; attention, \eng{actually} ne signifie jamais « actuellement » — ici le sens est bien « en fait ».
\item « transmettent » = \eng{pass on ... to ...} ou \eng{hand down ... to ...}.
\end{itemize}
\end{corrige}

\begin{corrige}[Exercice 11 — Expression écrite : composition]
\textbf{Proposition de rédaction modèle (environ 95 mots) :}

\eng{My favourite recreational activity is basketball. I play it every Wednesday and Saturday at my school's sports ground in Yaoundé, with my teammates from the school club. Before each match, we do warm-up exercises, and sometimes we go jogging to stay fit. I really enjoy this sport because I have fun with my friends and it helps me forget the stress of exams. Last term, our team took part in a big tournament and we won second prize. I always look forward to the next training session. Basketball is more than a hobby for me: it is a passion!}

\textbf{Barème indicatif (10 points) :}
\begin{itemize}
\item Respect de la consigne et du nombre de mots (80--100) : 2 points.
\item Trois collocations correctes (\eng{play basketball, do warm-up exercises, go jogging}) : 3 points.
\item Deux expressions idiomatiques (\eng{have fun, take part in, look forward to}) : 2 points.
\item Correction grammaticale (accords, temps, articles) : 2 points.
\item Orthographe et ponctuation : 1 point.
\end{itemize}

\textbf{Points de vigilance :} évite \eng{make sport}, \eng{I am agree}, \eng{play piano} sans \eng{the} ; utilise le présent simple pour les habitudes (\eng{I play every Saturday}) ; soigne la ponctuation anglaise (pas d'espace avant les deux-points en anglais) ; \eng{look forward to} est suivi d'un nom ou d'un verbe en \eng{-ing} : \eng{I look forward to seeing you}.
\end{corrige}

\newpage

\coursec{Fiche bilan}

\begin{popfigure}
\begin{tikzpicture}[
  mindmap,
  grow cyclic,
  every node/.style={concept, align=center, text=white, font=\small},
  root concept/.append style={concept color=popblue, font=\small\bfseries, minimum size=3.1cm},
  level 1/.append style={level distance=3.6cm, sibling angle=51.4},
  level 1 concept/.append style={minimum size=2.4cm, font=\scriptsize},
  text width=2.1cm
]
\node[root concept, align=center]{Recreational activities\\ \emph{loisirs}}
  child[concept color=poporange]{ node[align=center]{Activities\\ play, do, go\\ sports, games, music} }
  child[concept color=popgreen]{ node[align=center]{Collocations\\ play football\\ do yoga\\ go swimming} }
  child[concept color=poppurple]{ node[align=center]{Idioms\\ have fun\\ kill time\\ join in} }
  child[concept color=poppink]{ node[align=center]{False friends\\ \eng{actually} $\neq$ actuellement} }
  child[concept color=popgold]{ node[align=center]{Word families\\ relax, relaxation, relaxing} }
  child[concept color=popblue!70]{ node[align=center]{Pronunciation\\ leisure\\ ``lé-jeur''} }
  child[concept color=popmuted]{ node[align=center]{In context\\ dialogue, text, writing} };
\end{tikzpicture}
\legende{Carte mentale de la leçon : le vocabulaire des activités récréatives.}
\end{popfigure}

\begin{bilan}
\textbf{1. Lexique clé (anglais -- français)}
\begin{itemize}\setlength{\itemsep}{1pt}
  \item \eng{a hobby} : un passe-temps ; \eng{leisure/free time} : les loisirs, le temps libre
  \item \eng{to relax / to unwind} : se détendre ; \eng{to have fun / to enjoy oneself} : s'amuser
  \item \eng{to take part in / to join in} : participer à ; \eng{to practise a sport} : pratiquer un sport
  \item \eng{a board game} : un jeu de société ; \eng{a match / a game} : un match / une partie
  \item \eng{a team member} : un membre d'équipe ; \eng{a fan / a supporter} : un supporter
  \item \eng{indoor / outdoor activities} : activités en salle / en plein air
  \item \eng{to watch a film / to listen to music} : regarder un film / écouter de la musique
\end{itemize}

\textbf{2. Collocations à connaître par c\oe ur}
\begin{center}
\begin{tabularx}{\linewidth}{|l|X|X|}\hline
\rowcolor{popblueL} \textbf{Verbe} & \textbf{Emploi} & \textbf{Exemples} \\ \hline
\eng{play} & sports avec ballon et jeux & \eng{play football, play cards, play the guitar} \\ \hline
\eng{do} & activités individuelles, sans ballon & \eng{do karate, do gymnastics, do yoga} \\ \hline
\eng{go} & verbes en \emph{-ing} & \eng{go swimming, go dancing, go hiking} \\ \hline
\end{tabularx}
\end{center}

\textbf{3. Expressions idiomatiques utiles}
\begin{itemize}\setlength{\itemsep}{1pt}
  \item \eng{to kill time} : tuer le temps ; \eng{to be keen on} : être passionné de
  \item \eng{to be fond of} : aimer beaucoup ; \eng{to let off steam} : se défouler
  \item \eng{It is my cup of tea} : c'est mon truc ; \eng{to get together} : se retrouver
\end{itemize}

\textbf{4. Faux amis et pièges}
\begin{itemize}\setlength{\itemsep}{1pt}
  \item \eng{actually} = en fait ($\neq$ actuellement = \eng{currently, nowadays})
  \item \eng{a library} = une bibliothèque ($\neq$ librairie = \eng{a bookshop})
  \item \eng{sensible} = raisonnable ($\neq$ sensible = \eng{sensitive})
  \item \eng{to assist} = aider ; assister à = \eng{to attend}
  \item On dit \eng{to play football}, jamais \eng{to play at football}.
\end{itemize}

\textbf{5. Familles de mots et prononciation}
\begin{itemize}\setlength{\itemsep}{1pt}
  \item \eng{relax} (v.) $\rightarrow$ \eng{relaxation} (n.) $\rightarrow$ \eng{relaxing / relaxed} (adj.)
  \item \eng{enjoy} (v.) $\rightarrow$ \eng{enjoyment} (n.) $\rightarrow$ \eng{enjoyable} (adj.)
  \item \eng{compete} (v.) $\rightarrow$ \eng{competition} (n.) $\rightarrow$ \eng{competitive} (adj.)
  \item Prononciation : \eng{leisure} « lé-jeur », \eng{choir} « kouaïeur », \eng{guitar} « gui-tar ».
\end{itemize}

\textbf{6. Méthode express}
\begin{enumerate}\setlength{\itemsep}{1pt}
  \item Repérer le verbe de la collocation : ballon/jeu $\rightarrow$ \eng{play} ; effort individuel $\rightarrow$ \eng{do} ; \emph{-ing} $\rightarrow$ \eng{go}.
  \item Vérifier le faux ami : le mot anglais ressemble-t-il au français ? Si oui, contrôler le sens.
  \item À l'écrit, varier : \eng{enjoy, be keen on, be fond of, take part in} au lieu de répéter \eng{like}.
\end{enumerate}
\end{bilan}

\begin{aretenir}[Checklist avant le Bac]
\begin{itemize}\setlength{\itemsep}{2pt}
  \item[$\square$] Je connais au moins 15 mots du champ lexical des loisirs avec leur traduction.
  \item[$\square$] Je choisis correctement entre \eng{play}, \eng{do} et \eng{go} sans hésiter.
  \item[$\square$] Je sais que \eng{play} + instrument prend \eng{the} : \eng{play the piano}.
  \item[$\square$] Je distingue \eng{actually} (en fait) et \eng{currently} (actuellement).
  \item[$\square$] Je ne traduis pas « assister à » par \eng{to assist} mais par \eng{to attend}.
  \item[$\square$] J'emploie \eng{to take part in} et \eng{to join in} pour dire « participer ».
  \item[$\square$] Je connais 3 expressions idiomatiques : \eng{kill time}, \eng{be keen on}, \eng{let off steam}.
  \item[$\square$] Je forme les dérivés : \eng{relax $\rightarrow$ relaxation $\rightarrow$ relaxing}.
  \item[$\square$] Je prononce correctement \eng{leisure} (« lé-jeur ») et \eng{guitar} (« gui-tar »).
  \item[$\square$] Je peux écrire 8 lignes sur mes loisirs en variant les verbes et les collocations.
\end{itemize}
\end{aretenir}

\findecours

\end{document}
